% Utilisation d’un tableur — Informatique 3ème — Cours signé OnBuch+
% Programme officiel MINESEC. Compiler avec : tectonic N06-utilisation-tableur.tex
\newcommand{\DOCMATIERE}{Informatique}
\newcommand{\DOCNIVEAU}{3ème}
\newcommand{\DOCMODULE}{Informatique – classe de 3ème}
\newcommand{\DOCLECON}{N06}
\newcommand{\DOCTITRE}{Utilisation d’un tableur}
% OnBuch+ — préambule « cours signés » ENRICHI (compilé avec tectonic / XeLaTeX)
% Système visuel « Pop » d'OnBuch+ : fond crème (jamais blanc), bordures encre
% épaisses, ombres « dures » décalées, titres Archivo Black en capitales.
% Version MATHÉMATIQUES : graphes (pgfplots/tikz), boîtes pédagogiques
% (définition, propriété, méthode, attention, le savais-tu, exercice résolu,
% exercice, TP, bilan), page de garde et signature OnBuch+ ancrée en bas.
%
% Macros attendues dans main.tex AVANT \input{preamble} :
%   \DOCMATIERE  \DOCNIVEAU  \DOCMODULE  \DOCLECON (numéro)  \DOCTITRE
\documentclass[11pt]{article}

\usepackage{fontspec}
\usepackage[french]{babel}
\usepackage[a4paper,margin=2cm,top=3.4cm,bottom=2.8cm,headheight=1.4cm,footskip=1.3cm]{geometry}
\usepackage{amsmath,amssymb,amsfonts}
\usepackage{mathtools} % \xmapsto, \coloneqq... souvent utilisés par les modèles
\usepackage{cancel} % simplifications barrées (bilans de forces)
% \abs{z} (module, valeur absolue) et \norm{u} (norme), utilisés par les modèles
\providecommand{\abs}{}\let\abs\relax\DeclarePairedDelimiter{\abs}{\lvert}{\rvert}
\providecommand{\norm}{}\let\norm\relax\DeclarePairedDelimiter{\norm}{\lVert}{\rVert}
\usepackage[table]{xcolor} % [table] : fournit aussi \rowcolors (alternance de lignes)
\usepackage{pagecolor}
\usepackage{tikz}
\usetikzlibrary{arrows.meta,positioning,calc,shapes.geometric,shapes.misc,shapes.arrows,shapes.symbols,decorations.pathmorphing,decorations.pathreplacing,decorations.markings,patterns,shadows,mindmap,trees,fit,backgrounds,matrix,intersections,angles,quotes}
\usepackage{pgfplots}
\usepackage[version=4]{mhchem} % équations nucléaires \ce{^{235}_{92}U}
\usepackage[european,siunitx]{circuitikz} % schémas électriques
\pgfplotsset{compat=1.18}
\usepgfplotslibrary{fillbetween,groupplots} % aires entre courbes (calcul intégral)
\usepackage{enumitem}
\usepackage{fancyhdr}
% [most] charge toutes les bibliothèques tcolorbox (listings, minted, raster,
% documentation...) et gonfle inutilement la mémoire TeX sur un document long
% (~300 boîtes) : on ne charge que ce qui est réellement utilisé.
\usepackage[skins,breakable]{tcolorbox}
\usepackage{graphicx}
\usepackage{colortbl}
\usepackage{booktabs}
\usepackage{tabularx}
\usepackage{diagbox} % case d'en-tête diagonale des tableaux à double entrée
% Colonnes extensibles centrée / gauche / droite (C, L, R), souvent utilisées par les modèles
\newcolumntype{C}{>{\centering\arraybackslash}X}
\newcolumntype{L}{>{\raggedright\arraybackslash}X}
\newcolumntype{R}{>{\raggedleft\arraybackslash}X}
\usepackage{array}
\usepackage{multirow}
\usepackage{multicol}
\usepackage{siunitx}
% parse-numbers=false : accepte aussi \SI{2,5 \times 10^{-3}}{...}
\sisetup{locale=FR,output-decimal-marker={,},group-minimum-digits=4,parse-numbers=false}
\DeclareSIUnit\ohm{\ensuremath{\symup{\Omega}}} % U+2126 absent de la police
\DeclareSIUnit\division{div} % réglages d'oscilloscope (ms/div, V/div)
\DeclareSIUnit\tr{tr} % tours (vitesse de rotation, ex. \SI{1500}{\tr\per\minute})
\DeclareSIUnit\bit{bit}
\DeclareSIUnit\byte{o} % octet
\DeclareSIUnit\inch{po} % pouce
\DeclareSIUnit\ppp{ppp} % points par pouce (résolution d'image)
\usepackage{qrcode}
% Blocs de code source (PHP, C, HTML, JavaScript, SQL) — spécifique à la
% matière Informatique : listings + habillage aux couleurs OnBuch+ (voir le
% style \lstset ci-dessous, à côté des autres boîtes pédagogiques).
\usepackage{listings}
% \degree : commande fréquemment utilisée par les modèles pour un angle ou une
% température en dehors de \SI{}{} (non fournie par les paquets chargés).
\providecommand{\degree}{^{\circ}}
% \qed (fin de démonstration), utilisé par les modèles sans amsthm
\providecommand{\qed}{\hfill$\square$}
% \barre : double barre des valeurs interdites dans un tableau de variations (tkz-tab)
\providecommand{\barre}{\ensuremath{\|}}
% environnement proof (amsthm non chargé) utilisé par les modèles
\ifdefined\proof\else\newenvironment{proof}[1][Démonstration]{\par\noindent\textit{#1.}\ }{\qed\par}\fi
\usepackage{lastpage}
\usepackage{hyperref}
\hypersetup{hidelinks}

% ============================================================
% Palette « Pop » OnBuch+ — mêmes hex que la classe P (plus_ui.dart)
% ============================================================
\definecolor{poporange}{HTML}{F59321}
\definecolor{popdark}{HTML}{E07A0C}
\definecolor{popcream}{HTML}{FFF6E8}
\definecolor{popink}{HTML}{1C1714}
\definecolor{poppurple}{HTML}{7A5AE0}
\definecolor{popgreen}{HTML}{2E9E6A}
\definecolor{poppink}{HTML}{E0457B}
\definecolor{popblue}{HTML}{2F7DE1}
\definecolor{popgold}{HTML}{C98A12}
\definecolor{popmuted}{HTML}{8A7F76}
\definecolor{popline}{HTML}{EADFD2}
\definecolor{popgray}{HTML}{8A7F76}
\colorlet{popgrey}{popgray}
% Alias tolérants (le modèle nomme parfois les couleurs différemment)
\colorlet{popwhite}{white}
\colorlet{popblack}{popink}
\colorlet{popred}{poppink}
\colorlet{popyellow}{popgold}
% Teintes claires pour fonds de tableaux / zones
\colorlet{popblueL}{popblue!10!white}
\colorlet{poporangeL}{poporange!14!white}
\colorlet{popgreenL}{popgreen!12!white}
\colorlet{poppurpleL}{poppurple!12!white}
\colorlet{poppinkL}{poppink!12!white}
\colorlet{popgoldL}{popgold!14!white}

\pagecolor{popcream}

% ============================================================
% Typographie
% ============================================================
\defaultfontfeatures{Ligatures=TeX}
\setmainfont{PlusJakartaSans-Medium.ttf}[Path=../../../fonts/,BoldFont=PlusJakartaSans-Bold.ttf]
% Maths en sans-serif (Fira Math) avec chiffres et lettres droites de Plus
% Jakarta, pour que les formules se fondent dans le texte.
\usepackage[math-style=ISO,bold-style=ISO]{unicode-math}
\setmathfont{FiraMath-Regular.otf}[Path=../../../fonts/]
\setmathfont{PlusJakartaSans-Medium.ttf}[Path=../../../fonts/,range={up/num,up/{Latin,latin},\mathup}]
\usepackage{icomma} % virgule décimale sans espace en mode math
% Caractères mathématiques Unicode absents des polices de texte (Plus Jakarta,
% Archivo Black) : rendus par leur équivalent mathématique, en texte comme en
% formule — sinon ils s'affichent en carré vide (ex. « Arithmétique dans ℤ »).
\usepackage{newunicodechar}
\newunicodechar{ℝ}{\ensuremath{\mathbb{R}}}
\newunicodechar{ℤ}{\ensuremath{\mathbb{Z}}}
\newunicodechar{ℂ}{\ensuremath{\mathbb{C}}}
\newunicodechar{ℕ}{\ensuremath{\mathbb{N}}}
\newunicodechar{ℚ}{\ensuremath{\mathbb{Q}}}
\newunicodechar{𝔻}{\ensuremath{\mathbb{D}}}
\newunicodechar{α}{\ensuremath{\alpha}}
\newunicodechar{β}{\ensuremath{\beta}}
\newunicodechar{γ}{\ensuremath{\gamma}}
\newunicodechar{δ}{\ensuremath{\delta}}
\newunicodechar{ε}{\ensuremath{\varepsilon}}
\newunicodechar{θ}{\ensuremath{\theta}}
\newunicodechar{λ}{\ensuremath{\lambda}}
\newunicodechar{μ}{\ensuremath{\mu}}
\newunicodechar{σ}{\ensuremath{\sigma}}
\newunicodechar{τ}{\ensuremath{\tau}}
\newunicodechar{φ}{\ensuremath{\varphi}}
\newunicodechar{ψ}{\ensuremath{\psi}}
\newunicodechar{ω}{\ensuremath{\omega}}
\newunicodechar{Σ}{\ensuremath{\Sigma}}
% majuscules produites par \MakeUppercase dans les titres (α-aminés -> Α)
\newunicodechar{Α}{\ensuremath{\alpha}}
\newunicodechar{Β}{\ensuremath{\beta}}
\newunicodechar{Γ}{\ensuremath{\Gamma}}
\newunicodechar{Δ}{\ensuremath{\Delta}}
\newunicodechar{Ε}{\ensuremath{\varepsilon}}
\newunicodechar{Θ}{\ensuremath{\Theta}}
\newunicodechar{Λ}{\ensuremath{\Lambda}}
\newunicodechar{Μ}{\ensuremath{\mu}}
\newunicodechar{Τ}{\ensuremath{\tau}}
\newunicodechar{Φ}{\ensuremath{\Phi}}
\newunicodechar{Ψ}{\ensuremath{\Psi}}
\newunicodechar{Ω}{\ensuremath{\Omega}}
\newunicodechar{↦}{\ensuremath{\mapsto}}
\newunicodechar{∈}{\ensuremath{\in}}
\newunicodechar{∉}{\ensuremath{\notin}}
\newunicodechar{∧}{\ensuremath{\wedge}}
\newunicodechar{⇔}{\ensuremath{\Leftrightarrow}}
\newunicodechar{⟺}{\ensuremath{\Longleftrightarrow}}
\newunicodechar{⇒}{\ensuremath{\Rightarrow}}
\newunicodechar{⟹}{\ensuremath{\Longrightarrow}}
\newunicodechar{∘}{\ensuremath{\circ}}
\newunicodechar{∪}{\ensuremath{\cup}}
\newunicodechar{∩}{\ensuremath{\cap}}
\newunicodechar{≡}{\ensuremath{\equiv}}
\newunicodechar{⊂}{\ensuremath{\subset}}
\newunicodechar{⊥}{\ensuremath{\perp}}
\newunicodechar{∃}{\ensuremath{\exists}}
\newunicodechar{∀}{\ensuremath{\forall}}
\newunicodechar{∛}{\ensuremath{\sqrt[3]{\,}}}
\newunicodechar{∜}{\ensuremath{\sqrt[4]{\,}}}
\newunicodechar{⃗}{\ensuremath{{}^{\to}}}
\newunicodechar{ⁿ}{\ifmmode^{n}\else\textsuperscript{\ensuremath{n}}\fi}
\newunicodechar{ˣ}{\ifmmode^{x}\else\textsuperscript{\ensuremath{x}}\fi}
\newunicodechar{ᵇ}{\ifmmode^{b}\else\textsuperscript{\ensuremath{b}}\fi}
\newunicodechar{ʳ}{\ifmmode^{r}\else\textsuperscript{\ensuremath{r}}\fi}
\newunicodechar{ᵘ}{\ifmmode^{u}\else\textsuperscript{\ensuremath{u}}\fi}
\newunicodechar{ʸ}{\ifmmode^{y}\else\textsuperscript{\ensuremath{y}}\fi}
\newunicodechar{ᵃ}{\ifmmode^{a}\else\textsuperscript{\ensuremath{a}}\fi}
\newunicodechar{ᵉ}{\ifmmode^{e}\else\textsuperscript{\ensuremath{e}}\fi}
\newunicodechar{ᵏ}{\ifmmode^{k}\else\textsuperscript{\ensuremath{k}}\fi}
\newunicodechar{ᵖ}{\ifmmode^{p}\else\textsuperscript{\ensuremath{p}}\fi}
\newunicodechar{ᴺ}{\ifmmode^{N}\else\textsuperscript{\ensuremath{N}}\fi}
\newunicodechar{ᶜ}{\ifmmode^{c}\else\textsuperscript{\ensuremath{c}}\fi}
\newunicodechar{ʰ}{\ifmmode^{h}\else\textsuperscript{\ensuremath{h}}\fi}
\newunicodechar{ᵠ}{\ifmmode^{q}\else\textsuperscript{\ensuremath{q}}\fi}
\newunicodechar{ₙ}{\ifmmode_{n}\else\textsubscript{\ensuremath{n}}\fi}
\newunicodechar{ₐ}{\ifmmode_{a}\else\textsubscript{\ensuremath{a}}\fi}
\newunicodechar{ᵢ}{\ifmmode_{i}\else\textsubscript{\ensuremath{i}}\fi}
\newunicodechar{ᵣ}{\ifmmode_{r}\else\textsubscript{\ensuremath{r}}\fi}
\newunicodechar{ₖ}{\ifmmode_{k}\else\textsubscript{\ensuremath{k}}\fi}
\newunicodechar{ᵦ}{\ifmmode_{\beta}\else\textsubscript{\ensuremath{\beta}}\fi}
\newunicodechar{ₓ}{\ifmmode_{x}\else\textsubscript{\ensuremath{x}}\fi}
\newfontfamily\popdisplay{ArchivoBlack-Regular.ttf}[Path=../../../fonts/]
\newfontfamily\poplogo{SpaceGrotesk-Bold.ttf}[Path=../../../fonts/]
\newfontfamily\popsemi{PlusJakartaSans-SemiBold.ttf}[Path=../../../fonts/]
\newfontfamily\popxbold{PlusJakartaSans-ExtraBold.ttf}[Path=../../../fonts/]
\newfontfamily\popxboldls{PlusJakartaSans-ExtraBold.ttf}[Path=../../../fonts/,LetterSpace=3.0]

\setlist[enumerate]{leftmargin=1.7em,itemsep=5pt,topsep=4pt}
\setlist[itemize]{leftmargin=1.7em,itemsep=4pt,topsep=4pt,label={\color{poporange}\textbullet}}
\setlength{\parindent}{0pt}
\setlength{\parskip}{5pt}
\renewcommand{\arraystretch}{1.25}

% pgfplots : style Pop par défaut
\pgfplotsset{
  every axis/.append style={axis line style={popink,line width=1pt},tick style={popink},
    grid style={popline,line width=0.5pt},label style={font=\small},tick label style={font=\footnotesize}},
  popcurve/.style={poporange,line width=1.8pt,no markers,smooth},
}
% Même style disponible dans un \draw TikZ simple (hors pgfplots)
\tikzset{popcurve/.style={poporange,line width=1.8pt}}
\tikzset{popgrid/.style={popline,very thin}}
\colorlet{popcurve}{poporange} % le nom du style est parfois utilisé comme couleur

% ============================================================
% Pill / squiggle
% ============================================================
\newcommand{\poppill}[3]{%
  \tikz[baseline=-0.55ex]\node[draw=popink,line width=1.2pt,rounded corners=2.6pt,%
    fill=#2,inner xsep=6.5pt,inner ysep=3pt]{%
    \popxboldls\fontsize{7}{7}\selectfont\color{#3}\MakeUppercase{#1}};%
}
\newcommand{\popsquiggle}[1]{%
  \tikz[baseline=-0.4ex]{
    \draw[#1,line width=2.6pt,line cap=round]
      plot [smooth] coordinates {(0,0.09) (0.34,0.01) (0.68,0.21) (1.02,0.01) (1.36,0.10)};
  }%
}
% Mot-clé surligné (marqueur orange sous le texte)
\newcommand{\cle}[1]{{\popxbold\color{popdark}#1}}

% ============================================================
% En-tête / pied
% ============================================================
\pagestyle{fancy}
\fancyhf{}
\renewcommand{\headrulewidth}{0pt}
\renewcommand{\footrulewidth}{0pt}
\lhead{\raisebox{-0.15ex}{\poplogo\fontsize{13}{13}\selectfont\color{popink}OnBuch\color{poporange}+}}
\rhead{\poppill{\DOCMATIERE\ \textbullet\ \DOCNIVEAU\ \textbullet\ Leçon \DOCLECON}{white}{popink}}
\lfoot{\popxbold\fontsize{7.6}{7.6}\selectfont\color{popmuted}\MakeUppercase{Cours signé OnBuch+}\ \textbullet\ {\popsemi\DOCTITRE}}
\rfoot{\tikz[baseline=-0.6ex]\node[rounded corners=8.5pt,fill=popink,minimum height=17pt,minimum width=17pt,inner xsep=5pt,inner sep=0]{\popxbold\fontsize{8}{8}\selectfont\color{popcream}\thepage\,/\,\pageref*{LastPage}};}

% ============================================================
% Titres
% ============================================================
\newcommand{\chaptitre}[2]{%
  \begin{tcolorbox}[enhanced,colback=poporange,colframe=popink,boxrule=2.6pt,arc=11pt,
      drop shadow=popink,left=15pt,right=15pt,top=12pt,bottom=11pt,before skip=3pt,after skip=18pt]
    \poppill{#2}{popink}{popcream}\par\vspace{9pt}
    {\raggedright\hyphenpenalty=10000\exhyphenpenalty=10000\popdisplay\fontsize{22}{24}\selectfont\color{popcream}\MakeUppercase{#1}\par}
  \end{tcolorbox}
}
% Section numérotée : pastille orange avec le numéro + titre Archivo
\newcounter{popsec}
\newcommand{\coursec}[1]{%
  \stepcounter{popsec}\setcounter{popsub}{0}%
  \par\vspace{16pt}\noindent
  \begin{minipage}{\linewidth}
  \tikz[baseline=-0.7ex]\node[draw=popink,line width=1.6pt,fill=poporange,rounded corners=4pt,minimum size=22pt,inner sep=0]{\popdisplay\fontsize{12}{12}\selectfont\color{popcream}\thepopsec};%
  \hspace{8pt}{\popdisplay\fontsize{15}{17}\selectfont\color{popink}\MakeUppercase{#1}}\par
  \vspace{4pt}\noindent{\color{poporange}\rule{36pt}{3.6pt}}
  \end{minipage}\par\nopagebreak\vspace{8pt}
}
\newcounter{popsub}
\newcommand{\courssub}[1]{%
  \stepcounter{popsub}%
  \par\vspace{9pt}\noindent{\popxbold\fontsize{11.5}{13}\selectfont\color{popink}{\color{poporange}\thepopsec.\thepopsub}\hspace{6pt}#1}\par\nopagebreak\vspace{4pt}
}

% ============================================================
% Boîtes pédagogiques — même recette : fond blanc, bordure épaisse colorée,
% ombre dure encre, badge plein. Titre optionnel [#1].
% ============================================================
\tcbset{popbase/.style={breakable,enhanced,colback=white,boxrule=2.3pt,arc=9pt,
  shadow={3pt}{-3pt}{0pt}{popink},left=14pt,right=14pt,top=11pt,bottom=11pt,before skip=11pt,after skip=15pt,
  fontupper=\popsemi\fontsize{10.6}{14.5}\selectfont\color{popink},
  fontlower=\popsemi\fontsize{10.6}{14.5}\selectfont\color{popink}}}
% Coche dessinée (le glyphe ✓ n'existe pas dans les polices du document)
\newcommand{\popcheck}[1]{\tikz[baseline=-0.5ex]\draw[#1,line width=1.6pt,line cap=round,line join=round](-0.13,0.0)--(-0.04,-0.09)--(0.13,0.1);}
\providecommand{\coche}{\popcheck{popgreen}}
\providecommand{\resultat}[1]{\par\smallskip\noindent\fcolorbox{popgreen}{popgreenL}{\parbox{\dimexpr\linewidth-2\fboxsep-2\fboxrule\relax}{\textbf{Résultat.} #1}}\par\smallskip}
\newcommand{\popboxhead}[3]{\poppill{#1}{#2}{white}\ifx\relax#3\relax\else\hspace{6pt}{\popxbold\fontsize{10.6}{12}\selectfont\color{popink}#3}\fi\par\vspace{7pt}}

\newtcolorbox{aretenir}[1][]{popbase,colframe=popgreen,before upper={\popboxhead{À retenir}{popgreen}{#1}}}
\newtcolorbox{exemplebox}[1][]{popbase,colframe=poporange,before upper={\popboxhead{Exemple}{poporange}{#1}}}
\newtcolorbox{definition}[1][]{popbase,colframe=popblue,before upper={\popboxhead{Définition}{popblue}{#1}}}
\newtcolorbox{propriete}[1][]{popbase,colframe=poppurple,before upper={\popboxhead{Propriété}{poppurple}{#1}}}
\newtcolorbox{methode}[1][]{popbase,colframe=popgold,before upper={\popboxhead{Méthode}{popgold}{#1}}}
\newtcolorbox{attention}[1][]{popbase,colframe=poppink,before upper={\popboxhead{Attention}{poppink}{#1}}}
\newtcolorbox{experience}[1][]{popbase,colframe=popblue,colback=popblueL,before upper={\popboxhead{Expérience}{popblue}{#1}}}
% « Le savais-tu ? » : carte violette pleine (contexte, culture, Cameroun)
\newtcolorbox{savaistu}[1][]{popbase,colback=poppurple,colframe=popink,
  fontupper=\popsemi\fontsize{10.4}{14.2}\selectfont\color{white},
  before upper={\poppill{Le savais-tu\,?}{white}{poppurple}\ifx\relax#1\relax\else\hspace{6pt}{\popxbold\color{white}#1}\fi\par\vspace{7pt}}}
% Exercice résolu : énoncé en haut, \tcblower puis la solution sur fond crème
\newcounter{popexo}\newcounter{popexr}
\newtcolorbox{exoresolu}[1][]{popbase,colframe=poporange,colbacklower=popcream,
  segmentation style={popink,line width=1.2pt,dashed},
  before upper={\stepcounter{popexr}\popboxhead{Exercice résolu \thepopexr}{poporange}{#1}},
  before lower={\poppill{Solution}{popink}{popcream}\par\vspace{6pt}}}
% Exercice (non résolu en place) — la correction est dans la partie Corrigés
\newtcolorbox{exercice}[1][]{popbase,colframe=popink,
  before upper={\stepcounter{popexo}\popboxhead{Exercice \thepopexo}{popink}{#1}}}
% Corrigé
\newtcolorbox{corrige}[1][]{popbase,colframe=popgreen,colback=popgreenL,
  before upper={\popboxhead{Corrigé}{popgreen}{#1}}}
% Objectifs / prérequis / bilan
\newtcolorbox{objectifs}{popbase,colframe=popink,colback=poporangeL,
  before upper={\popboxhead{Objectifs de la leçon}{popink}{}}}
\newtcolorbox{prerequis}{popbase,colframe=popink,colback=popblueL,
  before upper={\popboxhead{Prérequis}{popblue}{}}}
\newtcolorbox{bilan}{popbase,colframe=popink,colback=white,boxrule=2.8pt,
  before upper={\popboxhead{Fiche bilan}{poporange}{L'essentiel à connaître pour l'examen}}}
% Figure encadrée avec légende
\newtcolorbox{popfigure}[1][]{enhanced,breakable=false,colback=white,colframe=popline,boxrule=1.4pt,arc=8pt,
  left=10pt,right=10pt,top=10pt,bottom=8pt,before skip=10pt,after skip=12pt,halign=center,
  lower separated=false,halign lower=center,
  fontlower=\popsemi\fontsize{9}{11.5}\selectfont\color{popmuted},
  #1}
\newcommand{\legende}[1]{\par\vspace{6pt}{\popsemi\fontsize{9}{11.5}\selectfont\color{popmuted}#1}}

% ============================================================
% Bloc de code source — \begin{lstlisting}[language=Php]...\end{lstlisting}
% (ou language=C, HTML, JavaScript, SQL ; option omise = pseudo-code brut).
% Même recette visuelle que les figures (\popfigure) : cadre encre épais,
% fond clair (popcream), légende possible juste après via \legende{...}
% comme pour une figure (listings ne sait pas arrondir ses coins : on garde
% un cadre rectangulaire, cohérent avec les autres tableaux du document).
% CONTRAT : les modèles ne doivent utiliser QUE \begin{lstlisting}[...] tel
% quel (pas de \begin{tcblisting}, pas de \verb pour un bloc multi-lignes).
% ============================================================
\lstdefinestyle{poplst}{
  backgroundcolor=\color{popcream},
  basicstyle=\popsemi\fontsize{8.6}{11}\selectfont\color{popink},
  keywordstyle=\bfseries\color{popblue},
  commentstyle=\itshape\color{popmuted},
  stringstyle=\color{popgreen},
  numberstyle=\tiny\color{popmuted},
  identifierstyle=\color{popink},
  frame=l,
  rulecolor=\color{popink},
  framerule=1.6pt,
  framesep=7pt,
  framexleftmargin=7pt,
  framexrightmargin=7pt,
  xleftmargin=2pt,
  xrightmargin=2pt,
  breaklines=true,
  breakatwhitespace=false,
  showstringspaces=false,
  showspaces=false,
  showtabs=false,
  tabsize=2,
  columns=flexible,
  keepspaces=true,
  numbers=none,
  aboveskip=10pt,
  belowskip=10pt,
  captionpos=b,
}
\lstset{style=poplst, language=PHP}
% La distribution TeX embarquée ne fournit pas le fichier de langue
% JavaScript de listings (lstlang*.dat) : on la redéfinit nous-mêmes
% pour que \begin{lstlisting}[language=JavaScript] fonctionne.
\lstdefinelanguage{JavaScript}{
  keywords={break,case,catch,class,const,continue,debugger,default,delete,do,
    else,export,extends,finally,for,function,if,import,in,instanceof,let,new,
    return,static,super,switch,this,throw,try,typeof,var,void,while,with,yield,
    async,await,of,get,set},
  keywordstyle=\bfseries\color{popblue},
  ndkeywords={true,false,null,undefined,NaN,Infinity},
  ndkeywordstyle=\bfseries\color{poporange},
  identifierstyle=\color{popink},
  sensitive=true,
  comment=[l]{//},
  morecomment=[s]{/*}{*/},
  string=[b]",
  morestring=[b]',
  morestring=[b]`,
}
% Même limitation pour CSS et les fichiers de configuration (apacheconf,
% nginx, ini...) : ils ne sont pas fournis. CSS et un style générique de
% type "config" (commentaires # ou ;, pas de mots-clés) couvrent les besoins.
\lstdefinelanguage{CSS}{
  keywords={color,background,background-color,margin,padding,border,
    border-radius,font,font-size,font-weight,font-family,width,height,
    display,position,top,left,right,bottom,float,clear,text-align,
    line-height,list-style,cursor,overflow,box-shadow,transition,
    flex,grid,align-items,justify-content,z-index},
  keywordstyle=\bfseries\color{popblue},
  identifierstyle=\color{popink},
  sensitive=false,
  comment=[s]{/*}{*/},
  string=[b]",
  morestring=[b]',
}
\lstdefinelanguage{apacheconf}{
  sensitive=false,
  morecomment=[l]{\#},
}
\lstdefinelanguage{Apache}{
  sensitive=false,
  morecomment=[l]{\#},
}
\lstdefinelanguage{text}{sensitive=false}
\lstdefinelanguage{powershell}{
  keywords={New-Object,New-SmbShare,Get-Acl,Set-Acl,Get-ChildItem,Set-Content,
    Get-Content,Write-Host,ForEach-Object,Where-Object,Select-Object,
    Test-Path,Remove-Item,Copy-Item,Move-Item,Start-Process,Stop-Process,
    If,Else,ElseIf,ForEach,While,Function,Return,Try,Catch},
  keywordstyle=\bfseries\color{popblue},
  identifierstyle=\color{popink},
  sensitive=false,
  comment=[l]{\#},
  string=[b]",
  morestring=[b]',
}
\lstdefinelanguage{ini}{
  sensitive=false,
  morecomment=[l]{;},
  morecomment=[l]{\#},
}

% ============================================================
% Signature OnBuch+
% ============================================================
\newcommand{\poplogoinline}[1]{{\poplogo\fontsize{#1}{#1}\selectfont\color{popink}OnBuch\color{poporange}+}}
\newcommand{\signatureonbuch}{%
  \begin{tcolorbox}[enhanced,colback=popink,colframe=popink,boxrule=2.2pt,arc=10pt,
      drop shadow=poporange,left=14pt,right=14pt,top=10pt,bottom=10pt,before skip=0pt,after skip=0pt]
    \tikz[baseline=-0.7ex]\node[circle,fill=popgreen,draw=popcream,line width=1.3pt,minimum size=22pt,inner sep=0]{\popcheck{white}};%
    \hspace{9pt}\begin{minipage}[c]{0.62\linewidth}
      {\popxbold\fontsize{12}{13}\selectfont\color{popcream}Signé\ \textbullet\ {\poplogo OnBuch\color{poporange}+}}\par\vspace{2pt}
      {\raggedright\popsemi\fontsize{8}{10.5}\selectfont\color{popline}Ludovic A.\\\MakeUppercase{Conforme au programme officiel MINESEC}\par}
    \end{minipage}\hfill
    {\poplogo\fontsize{22}{22}\selectfont\color{popcream}OnBuch\color{poporange}+}
  \end{tcolorbox}%
}

% ============================================================
% Page de garde
% #1 titre · #2 matière · #3 niveau · #4 module · #5 n° leçon · #6 durée ·
% #7 description / objectifs (texte ou itemize)
% ============================================================
\newcommand{\pagegarde}[7]{%
  \thispagestyle{empty}
  \newgeometry{a4paper,margin=1.7cm,top=1.6cm,bottom=1.4cm}%
  \begin{tikzpicture}[remember picture,overlay]
    \fill[poporange!18] ([xshift=-3.2cm,yshift=-3.6cm]current page.north east) circle (4.6cm);
    \fill[poppurple!14] ([xshift=2.4cm,yshift=7.6cm]current page.south west) circle (3.8cm);
  \end{tikzpicture}%
  {\poplogo\fontsize{30}{30}\selectfont\color{popink}OnBuch\color{poporange}+}\hfill
  \raisebox{0.6ex}{\poppill{Cours signé \textbullet\ Premium}{popink}{popcream}}\par
  \vspace{1pt}\popsquiggle{poppurple}\par
  \vspace{8pt}
  \begin{tcolorbox}[enhanced,colback=poporange,colframe=popink,boxrule=2.8pt,arc=13pt,
      drop shadow=popink,left=18pt,right=18pt,top=12pt,bottom=13pt,after skip=10pt]
    \poppill{#2 \textbullet\ #3}{popink}{popcream}\hspace{5pt}\poppill{#4}{white}{popink}\par\vspace{8pt}
    {\popdisplay\fontsize{14}{14}\selectfont\color{popink}LEÇON #5}\par\vspace{4pt}
    {\raggedright\hyphenpenalty=10000\exhyphenpenalty=10000\popdisplay\fontsize{25}{27}\selectfont\color{popcream}\MakeUppercase{#1}\par}
    \vspace{8pt}
    {\popxbold\fontsize{9.5}{9.5}\selectfont\color{popink}\MakeUppercase{Durée conseillée : #6}}
  \end{tcolorbox}
  \begin{tcolorbox}[enhanced,colback=white,colframe=popink,boxrule=2.2pt,arc=10pt,
      drop shadow=popink,left=16pt,right=16pt,top=10pt,bottom=10pt,after skip=8pt]
    \poppill{Dans cette leçon}{poppurple}{white}\par\vspace{5pt}
    \popsemi\fontsize{9.3}{12.6}\selectfont\color{popink}#7
  \end{tcolorbox}
  \noindent\begin{minipage}[t]{0.487\linewidth}
    \begin{tcolorbox}[enhanced,colback=popgreen,colframe=popink,boxrule=2.2pt,arc=10pt,
        drop shadow=popink,left=11pt,right=11pt,top=10pt,bottom=10pt,height=3.1cm,valign=center]
      \tcbox[colback=white,colframe=popink,boxrule=1.3pt,arc=5pt,boxsep=0pt,nobeforeafter,
        left=3pt,right=3pt,top=3pt,bottom=3pt,valign=center]{\qrcode[height=1.9cm]{https://play.google.com/store/apps/details?id=com.luvvix.onbuch&hl=fr}}%
      \hspace{8pt}\begin{minipage}[c]{0.5\linewidth}
        {\popdisplay\fontsize{11}{12.5}\selectfont\color{white}\MakeUppercase{Télécharge OnBuch}}\par\vspace{4pt}
        {\raggedright\popsemi\fontsize{8.4}{11}\selectfont\color{white}Gratuit \textbullet\ Google Play\\Tuteur IA, annales, concours, résultats\par}
      \end{minipage}
    \end{tcolorbox}
  \end{minipage}\hfill
  \begin{minipage}[t]{0.487\linewidth}
    \begin{tcolorbox}[enhanced,colback=poppurple,colframe=popink,boxrule=2.2pt,arc=10pt,
        drop shadow=popink,left=11pt,right=11pt,top=10pt,bottom=10pt,height=3.1cm,valign=center]
      \tikz[baseline=-0.7ex]\node[circle,fill=white,draw=popink,line width=1.3pt,minimum size=1.6cm,inner sep=0]{\popdisplay\fontsize{24}{24}\selectfont\color{poppurple}+};%
      \hspace{8pt}\begin{minipage}[c]{0.6\linewidth}
        {\popdisplay\fontsize{11}{12.5}\selectfont\color{white}\MakeUppercase{Passe à OnBuch+}}\par\vspace{4pt}
        {\raggedright\popsemi\fontsize{8.4}{11}\selectfont\color{white}Tous les cours signés, Tuteur IA illimité, fiches PDF — onglet OnBuch+ de l'app\par}
      \end{minipage}
    \end{tcolorbox}
  \end{minipage}
  \vspace{8pt}
  \popcontenu
  \vfill
  \signatureonbuch
  \restoregeometry
  \newpage
}

% Rangée « Ce cours contient » (page de garde)
\newcommand{\poptile}[3]{% #1 couleur #2 gros texte #3 légende
  \begin{tcolorbox}[enhanced,colback=#1,colframe=popink,boxrule=2pt,arc=9pt,shadow={2.5pt}{-2.5pt}{0pt}{popink},
      left=6pt,right=6pt,top=9pt,bottom=9pt,halign=center,valign=center,height=1.75cm,width=\linewidth,nobeforeafter]
    {\popdisplay\fontsize{12.5}{13}\selectfont\color{white}\MakeUppercase{#2}}\par\vspace{3pt}
    {\centering\popsemi\fontsize{7.8}{9.5}\selectfont\color{white}#3\par}
  \end{tcolorbox}}
\newcommand{\popcontenu}{%
  \par\noindent{\popxboldls\fontsize{7.5}{7.5}\selectfont\color{popmuted}CE COURS SIGNÉ CONTIENT}\par\vspace{7pt}
  \noindent\begin{minipage}[t]{0.235\linewidth}\poptile{popblue}{Cours}{complet, illustré, pas à pas}\end{minipage}\hfill
  \begin{minipage}[t]{0.235\linewidth}\poptile{poporange}{Méthodes}{et exercices résolus}\end{minipage}\hfill
  \begin{minipage}[t]{0.235\linewidth}\poptile{poppink}{Exercices}{type examen + corrigés}\end{minipage}\hfill
  \begin{minipage}[t]{0.235\linewidth}\poptile{popgreen}{Fiche bilan}{l'essentiel à retenir}\end{minipage}\par}

% Sommaire
\newtcolorbox{sommaire}{enhanced,colback=white,colframe=popink,boxrule=2.2pt,arc=10pt,
  drop shadow=popink,left=18pt,right=18pt,top=13pt,bottom=13pt,before skip=6pt,after skip=20pt,
  before upper={\poppill{Sommaire}{poppurple}{white}\par\vspace{9pt}\popsemi\fontsize{10.6}{14.5}\selectfont\color{popink}}}

% Signature de fin de cours — poussée tout en bas de la dernière page
\newcommand{\findecours}{%
  \par\vfill\nopagebreak
  \begin{center}\popsquiggle{poporange}\end{center}\vspace{4pt}
  \signatureonbuch
}
\AtBeginDocument{\def\llbracket{\mathopen{[\mkern-3.5mu[}}\def\rrbracket{\mathclose{]\mkern-3.5mu]}}} % intervalles d'entiers (glyphe ⟦ absent de la police)
% Glyphes absents de Fira Math : redéfinis à partir de glyphes disponibles
\AtBeginDocument{%
  \def\setminus{\mathbin{\backslash}}\let\smallsetminus\setminus
  \def\vdots{\mathord{\vbox{\baselineskip3.2pt\lineskiplimit0pt\kern4pt\hbox{.}\hbox{.}\hbox{.}}}}%
  \def\ddots{\mathinner{\mkern1mu\raise6.5pt\hbox{.}\mkern2mu\raise3.6pt\hbox{.}\mkern2mu\raise0.7pt\hbox{.}\mkern1mu}}%
  \def\triangle{\mathord{\scalebox{0.9}{$\symup{\Delta}$}}}%
  \def\nabla{\mathord{\raisebox{\depth}{\scalebox{1}[-1]{$\symup{\Delta}$}}}}%
  \def\checkmark{\popcheck{popdark}}%
  \def\star{\mathbin{\ast}}\def\bigstar{\mathord{\ast}}%
  \def\diamond{\mathbin{\scalebox{0.6}{$\Diamond$}}}%
  \def\bigcup{\mathop{\mathchoice{\raisebox{-0.3em}{\scalebox{1.7}{$\cup$}}}{\scalebox{1.25}{$\cup$}}{\cup}{\cup}}\displaylimits}%
  \def\bigcap{\mathop{\mathchoice{\raisebox{-0.3em}{\scalebox{1.7}{$\cap$}}}{\scalebox{1.25}{$\cap$}}{\cap}{\cap}}\displaylimits}%
  \def\lesseqqgtr{\mathrel{\lessgtr}}\def\bigoplus{\mathop{\oplus}\displaylimits}%
}
\providecommand{\ssi}{si et seulement si} % abréviation fréquente
\AtBeginDocument{\providecommand{\wideparen}{\overparen}} % arc de cercle

%% FIGFIX-BEGIN v3 — correctifs globaux des figures (docs/onbuch-generation/tools/figfix)
\usepackage{adjustbox}\usepackage{etoolbox}
% toute figure TikZ/pgfplots plus large que la ligne est réduite à la largeur disponible
\BeforeBeginEnvironment{tikzpicture}{\begin{adjustbox}{max width=\linewidth}}
\AfterEndEnvironment{tikzpicture}{\end{adjustbox}}
% légendes pgfplots lisibles par-dessus les courbes
\pgfplotsset{every axis/.append style={legend style={fill=white,fill opacity=0.92,text opacity=1,draw=black!15,inner sep=3pt}}}
% étiquettes d'arêtes (syntaxe edge["..."]) avec fond blanc
\tikzset{every edge quotes/.append style={fill=white,inner sep=1.2pt}}
% bulles de cartes mentales : texte à la ligne dans la bulle, bulle un peu plus grande
\tikzset{every concept/.append style={text width=2.0cm,align=center,minimum size=2.3cm,inner sep=2pt}}
%% FIGFIX-END
\begin{document}

\pagegarde{Utilisation d’un tableur}{Informatique}{3ème}{Informatique – classe de 3ème}{N06}{2 séances (fiche de progression harmonisée)}{Cette leçon initie à l'utilisation avancée d'un tableur (LibreOffice Calc / Excel) pour automatiser les calculs de gestion courante. Elle couvre les formules mathématiques, les fonctions logiques (SI), la mise en forme conditionnelle et la mise en page pour l'impression de rapports.
\vspace{4pt}
\begin{multicols}{2}\small
\begin{itemize}
\item À la fin de cette leçon, je sais saisir des données (texte, nombres, dates) et les mettre en forme (polices, bordures, alignement) pour obtenir un tableau lisible.
\item Je sais construire une formule de calcul simple en utilisant les opérateurs arithmétiques (+, -, *, /, \^{}) et respecter la priorité opératoire.
\item Je sais distinguer et utiliser les références relatives, absolues (\$) et mixtes pour recopier une formule sans erreur.
\item Je sais utiliser les fonctions mathématiques courantes (SOMME, MOYENNE, MAX, MIN, NB, NBVAL) sur des plages de cellules.
\item Je sais exploiter la fonction logique SI pour prendre une décision automatisée (ex: mention, remise, stock bas).
\item Je sais imbriquer plusieurs fonctions SI ou combiner SI avec ET/OU pour gérer des cas multiples (ex: grille de notation, tranches de TVA).
\end{itemize}\end{multicols}}

\begin{sommaire}
\begin{multicols}{2}
\begin{enumerate}
\item 1. Environnement et Mise en forme du classeur
\item 2. Formules, Opérateurs et Références de cellules
\item 3. Fonctions mathématiques et statistiques usuelles
\item 4. Fonctions logiques : SI, ET, OU, NON et imbrication
\item 5. Mise en forme conditionnelle et Visualisation graphique
\item 6. Mise en page, Aperçu et Impression du rapport
\item Activité d'intégration
\item Méthodes et exercices résolus
\item Exercices
\item Corrigés des exercices
\item Fiche bilan
\end{enumerate}
\end{multicols}
\end{sommaire}

\begin{objectifs}
\begin{itemize}
\item À la fin de cette leçon, je sais saisir des données (texte, nombres, dates) et les mettre en forme (polices, bordures, alignement) pour obtenir un tableau lisible.
\item Je sais construire une formule de calcul simple en utilisant les opérateurs arithmétiques (+, -, *, /, \^{}) et respecter la priorité opératoire.
\item Je sais distinguer et utiliser les références relatives, absolues (\$) et mixtes pour recopier une formule sans erreur.
\item Je sais utiliser les fonctions mathématiques courantes (SOMME, MOYENNE, MAX, MIN, NB, NBVAL) sur des plages de cellules.
\item Je sais exploiter la fonction logique SI pour prendre une décision automatisée (ex: mention, remise, stock bas).
\item Je sais imbriquer plusieurs fonctions SI ou combiner SI avec ET/OU pour gérer des cas multiples (ex: grille de notation, tranches de TVA).
\item Je sais appliquer une mise en forme conditionnelle pour mettre en évidence visuellement des valeurs cibles (ex: ventes > objectif en vert, stock < seuil en rouge).
\item Je sais configurer la mise en page (marges, orientation, en-têtes/pieds de page, zone d'impression) pour imprimer un rapport professionnel sur une page A4.
\item Je sais créer et personnaliser un graphique (histogramme, camembert, courbe) adapté aux données pour en faciliter l'analyse.
\end{itemize}
\end{objectifs}

\begin{prerequis}
\begin{itemize}
\item Notion de cellule, ligne, colonne, feuille de calcul et classeur (vu en 4ème).
\item Manipulation de base de la souris et du clavier (sélection, copier-coller, glisser-déposer).
\item Règles de priorité des opérateurs arithmétiques en mathématiques (parenthèses, puissances, mult/div, add/soustr).
\item Notion de variable et de condition (SI... ALORS... SINON) en algorithmique (4ème).
\item Calcul de pourcentages et lecture de graphiques simples (mathématiques 5ème/4ème).
\end{itemize}
\end{prerequis}

\newpage

\coursec{Environnement et Mise en forme du classeur}

Monsieur Kamga, gérant d'une boutique de téléphonie à Douala, note chaque soir ses ventes sur un cahier : modèles, quantités, prix unitaires. En fin de mois, il passe des heures à additionner les colonnes à la calculatrice pour connaître son chiffre d'affaires, identifier le téléphone le plus vendu et calculer la TVA à reverser. Une erreur de saisie l'oblige tout recommencer. Comment l'aider à automatiser ces calculs, à visualiser ses ventes par un graphique et à présenter un rapport professionnel en quelques clics ? C'est tout l'enjeu de cette leçon : transformer un cahier papier en un \cle{classeur} dynamique, fiable et lisible.

\courssub{Découverte de l'interface du tableur}

Un tableur (LibreOffice Calc, Microsoft Excel) se présente comme une grille immense de \cle{cellules}. Chaque cellule est identifiée par une \cle{adresse} (ou référence) composée d'une lettre de colonne et d'un numéro de ligne (ex. : \texttt{B3}). L'interface standard regroupe des zones fonctionnelles qu'il faut savoir nommer pour communiquer efficacement.

\begin{popfigure}
\begin{tikzpicture}[
    box/.style={draw=popdark, thick, rectangle, rounded corners=3pt, fill=popcream, inner sep=4pt},
    labelbox/.style={draw=popblue, thick, rectangle, rounded corners=3pt, fill=popblueL, inner sep=3pt, font=\small\bfseries},
    arrow/.style={->, >=Stealth, thick, poporange},
    font=\small
]
% --- Title Bar ---
\node[box, minimum width=16cm, minimum height=0.8cm] (title) at (0, 7.5) {};
\node[labelbox, anchor=north west] at (title.north west) {Barre de titre : Nom du fichier (ex. Ventes\_Kamga.ods)};
% --- Menu Bar ---
\node[box, minimum width=16cm, minimum height=0.7cm] (menu) at (0, 6.7) {};
\node[labelbox, anchor=north west] at (menu.north west) {Barre de menus : Fichier, Edition, Affichage, Insertion, Format, Outils, Donn\'ees, Fen\^etre, Aide};
% --- Toolbars ---
\node[box, minimum width=16cm, minimum height=0.7cm] (stdbar) at (0, 5.9) {};
\node[labelbox, anchor=north west] at (stdbar.north west) {Barre d'outils Standard : Nouveau, Ouvrir, Enregistrer, Imprimer, Copier, Coller, Annuler/R\'etablir};
\node[box, minimum width=16cm, minimum height=0.7cm] (fmtbar) at (0, 5.1) {};
\node[labelbox, anchor=north west] at (fmtbar.north west) {Barre de formatage : Police, Taille, Gras/Italique/Soulign\'e, Alignement, Fusionner, Format nombre, Bordures, Couleur fond};
% --- Formula Bar ---
\node[box, minimum width=16cm, minimum height=0.9cm] (formula) at (0, 4.1) {};
\node[labelbox, anchor=north west] at (formula.north west) {Barre de formule};
\node[draw=popgreen, thick, rounded corners, fill=popgreenL, inner sep=2pt, font=\ttfamily\small] at (-7.2, 4.1) {Zone de nom : B3};
\node[draw=poppurple, thick, rounded corners, fill=poppurpleL, inner sep=2pt, font=\ttfamily\small] at (-4.5, 4.1) {$\sum$ Assistant};
\node[draw=poppink, thick, rounded corners, fill=poppinkL, inner sep=2pt, minimum width=10cm] at (1.5, 4.1) {Zone de saisie : =B3*C3};
% --- Column Headers ---
\foreach \col/\x in {A/-7.2, B/-5.8, C/-4.4, D/-3.0, E/-1.6, F/-0.2} {
    \node[draw=popdark, thick, fill=popblueL, minimum width=1.4cm, minimum height=0.6cm] at (\x, 3.2) {\bfseries \col};
}
\node[draw=popdark, thick, fill=popblueL, minimum width=1.4cm, minimum height=0.6cm] at (1.2, 3.2) {\bfseries ...};
% --- Row Headers & Grid ---
\foreach \row/\y in {1/2.5, 2/1.8, 3/1.1, 4/0.4, 5/-0.3} {
    \node[draw=popdark, thick, fill=popblueL, minimum width=1.2cm, minimum height=0.6cm] at (-8.0, \y) {\bfseries \row};
}
% Grid cells
\foreach \col/\x in {A/-7.2, B/-5.8, C/-4.4, D/-3.0, E/-1.6, F/-0.2} {
    \foreach \row/\y in {1/2.5, 2/1.8, 3/1.1, 4/0.4, 5/-0.3} {
        \node[draw=popline, thin, minimum width=1.4cm, minimum height=0.6cm] (cell-\col-\row) at (\x, \y) {};
    }
}
% Highlight active cell B3
\node[draw=poporange, very thick, minimum width=1.4cm, minimum height=0.6cm] at (-5.8, 1.1) {};
\node[font=\small\bfseries, poporange] at (-5.8, 1.1) {Cellule active};
% Sheet Tabs
\node[draw=popdark, thick, fill=popmuted!18, minimum width=16cm, minimum height=0.6cm] (sheets) at (0, -1.2) {};
\node[labelbox, anchor=south west] at (sheets.south west) {Onglets de feuilles : Feuille1, Feuille2, Feuille3 | + (Nouvelle feuille) | Barre de d\'efilement horizontal | Zoom | Modes d'affichage};
% Annotations with arrows
\draw[arrow] (-8.5, 7.5) -- ++(-1.5, 0) node[left, align=right, font=\small] {Nom du fichier\\et application};
\draw[arrow] (-8.5, 4.1) -- ++(-1.5, 0) node[left, align=right, font=\small] {Adresse de la\\cellule active};
\draw[arrow] (8.5, 4.1) -- ++(1.5, 0) node[right, align=left, font=\small] {Contenu r\'eel\\formule ou valeur};
\draw[arrow] (-8.5, 3.2) -- ++(-1.5, 0) node[left, align=right, font=\small] {En-t\^etes\\(Lettres)};
\draw[arrow] (-8.5, 1.1) -- ++(-1.5, 0) node[left, align=right, font=\small] {En-t\^etes\\(Chiffres)};
\draw[arrow] (0, -1.2) -- ++(0, -1) node[below, align=center, font=\small] {Navigation entre\\les feuilles};
\end{tikzpicture}
\legende{Interface de LibreOffice Calc : identification des zones principales (barres de menus, d'outils, de formule, grille de cellules, onglets).}
\end{popfigure}

\begin{definition}[Cellule active]
C'est l'intersection d'une ligne (repérée par un \textbf{chiffre}) et d'une colonne (repérée par une \textbf{lettre}) où se trouve le curseur de saisie. Son adresse (ex. : \texttt{B3}) s'affiche dans la \cle{zone de nom} (à gauche de la barre de formule). Seule la cellule active reçoit la saisie clavier.
\end{definition}

\courssub{Saisie et types de données : étiquettes, valeurs, formules}

Dans une cellule, on ne « tape » pas n'importe comment : le tableur interprète le contenu selon sa nature. Il existe \textbf{trois types fondamentaux} de données.

\begin{propriete}[Les trois types de données dans une cellule]
\begin{itemize}
    \item \textbf{Étiquette (Texte) :} Suite de caractères (lettres, chiffres, symboles) non calculable. Alignée à \textbf{gauche} par défaut. Ex. : \texttt{Galaxy A54}, \texttt{Total}, \texttt{'01/01/2024} (apostrophe \texttt{'} forcée pour traiter une date comme du texte).
    \item \textbf{Valeur (Nombre, Date, Heure) :} Donnée numérique exploitable en calcul. Alignée à \textbf{droite} par défaut.
        \begin{itemize}
            \item Entier / Décimal : \texttt{50000}, \texttt{12,5}.
            \item Date : stockée comme un \textbf{numéro de série} (nombre de jours depuis le 30/12/1899). Ex. : \texttt{45292} pour le 01/01/2024.
            \item Heure : stockée comme une \textbf{fraction de jour} (0,5 = 12h00).
        \end{itemize}
    \item \textbf{Formule :} Expression commençant \textbf{obligatoirement par le signe égal \texttt{=}}. Elle combine opérateurs (\texttt{+ - * / \^{}}), références de cellules, constantes et fonctions. Ex. : \texttt{=B3*C3}, \texttt{=SOMME(D3:D10)}. Le tableur affiche le \textbf{résultat} mais stocke la \textbf{formule} (visible dans la barre de formule).
\end{itemize}
\end{propriete}

\begin{attention}[Affichage \emph{vs} Valeur réelle : le piège du pourcentage]
Confondre le contenu affiché et la valeur stockée est une source majeure d'erreurs.
\begin{itemize}
    \item Si vous saisissez \texttt{15\%} dans une cellule, le tableur \textbf{stocke 0,15} (la valeur réelle) et affiche \texttt{15\%} (grâce au format).
    \item Si vous saisissez \texttt{15} puis appliquez le format \%, l'affichage devient \texttt{1500\%} (car $15 = 1500\%$).
    \item \textbf{Règle d'or :} Les calculs utilisent \textbf{toujours la valeur réelle} (0,15), pas l'affichage (15\%).
\end{itemize}
\end{attention}

\begin{exemplebox}[Formatage Monétaire FCFA : Exemple chiffré]
Dans la cellule \texttt{C5}, on saisit le prix \texttt{50000}. On applique le format \textit{Monétaire} avec symbole \texttt{FCFA} (via \texttt{Ctrl+1} $\to$ Nombre $\to$ Monétaire $\to$ Symbole : \texttt{FCFA}).
\begin{itemize}
    \item \textbf{Affichage :} \texttt{50 000 FCFA} (avec espace insécable comme séparateur de milliers).
    \item \textbf{Valeur pour calcul :} \texttt{50000} (nombre pur, sans symbole, sans espace).
\end{itemize}
C'est cette distinction qui permet d'écrire \texttt{=C5*D5} pour obtenir le montant total sans erreur de conversion.
\end{exemplebox}

\courssub{Techniques de sélection : cibler pour agir}

Toute action de mise en forme, suppression ou copie commence par une \cle{sélection}. Maîtriser les raccourcis clavier fait gagner un temps précieux.

\begin{methode}[Sélectionner efficacement]
\begin{enumerate}
    \item \textbf{Cellule unique :} Clic gauche sur la cellule.
    \item \textbf{Plage contiguë (bloc rectangulaire) :} Clic sur la première cellule (ex. \texttt{A1}), maintenir \texttt{Shift} (Maj), clic sur la dernière (ex. \texttt{F10}). \emph{Astuce :} Clic-glisser souris.
    \item \textbf{Plages non contiguës :} Sélectionner la première plage, maintenir \texttt{Ctrl}, sélectionner la deuxième plage, etc. Utile pour mettre en forme les en-têtes \textbf{et} la ligne des totaux simultanément.
    \item \textbf{Ligne(s) entière(s) :} Clic sur le numéro de ligne (en-tête gauche). \texttt{Shift} + flèches pour étendre. \texttt{Ctrl} + clic pour lignes disjointes.
    \item \textbf{Colonne(s) entière(s) :} Clic sur la lettre de colonne (en-tête haut). Même logique que lignes.
    \item \textbf{Toute la feuille :} Clic sur le \textbf{carreau d'angle} (coin supérieur gauche, entre \texttt{A} et \texttt{1}) ou \texttt{Ctrl + A} (deux fois si curseur dans un tableau).
\end{enumerate}
\end{methode}

\courssub{Mise en forme des cellules : police, alignement, fusion}

Un tableau lisible guide l'œil : en-têtes distincts, données alignées logiquement, titres centrés.

\begin{propriete}[Principaux attributs de mise en forme (Onglet \textit{Format} $\to$ \textit{Cellules} ou \texttt{Ctrl+1})]
\begin{itemize}
    \item \textbf{Police :} Nom (Liberation Sans, DejaVu Sans), Taille (10-12 pt pour données, 14-16 pt titres), Style (Gras pour en-têtes/totaux), Couleur.
    \item \textbf{Alignement (Onglet \textit{Alignement}) :}
        \begin{itemize}
            \item Horizontal : Gauche (texte), Droite (nombres), Centré (en-têtes, cases à cocher), Justifié.
            \item Vertical : Haut, Centre (recommandé pour lignes hautes), Bas.
            \item \textbf{Renvoi à la ligne automatique :} Essentiel pour les libellés longs sans élargir la colonne indûment.
        \end{itemize}
    \item \textbf{Fusionnement :} Regroupe plusieurs cellules adjacentes en une seule grande cellule. \textbf{Attention :} Seule la valeur de la cellule en haut à gauche est conservée ; les autres sont effacées.
\end{itemize}
\end{propriete}

\begin{methode}[Centrer un titre sur plusieurs colonnes (Fusionner \& Centrer)]
\begin{enumerate}
    \item Saisir le titre dans la cellule la plus à gauche de la plage visée (ex. \texttt{A1}).
    \item Sélectionner la plage couvrant toutes les colonnes du tableau (ex. \texttt{A1:F1}).
    \item Ouvrir \textit{Format} $\to$ \textit{Cellules} (ou \texttt{Clic droit} $\to$ \textit{Formater les cellules} ou \texttt{Ctrl+1}).
    \item Onglet \textbf{Alignement} $\to$ Cocher \textbf{Fusionner les cellules}.
    \item Horizontal : \textbf{Centré} ; Vertical : \textbf{Centré}.
    \item Valider. Appliquer \textbf{Gras} et taille \textbf{14-16 pt} via la barre de formatage.
\end{enumerate}
\end{methode}

\begin{savaistu}[Le raccourci universel : \texttt{Ctrl + 1}]
Dans \textbf{tous} les tableurs majeurs (LibreOffice Calc, Excel, Google Sheets), \texttt{Ctrl + 1} ouvre instantanément la boîte de dialogue \textit{Formater les cellules} avec tous ses onglets : \textit{Nombre, Police, Effets de police, Alignement, Bordures, Arrière-plan, Protection}. C'est le raccourci « couteau suisse » du tableur.
\end{savaistu}

\courssub{Bordures, couleurs de fond et formatage des nombres}

La structure visuelle sépare l'information : l'en-tête, le corps, les totaux.

\begin{propriete}[Structurer visuellement : Bordures et Arrière-plan]
\begin{itemize}
    \item \textbf{Bordures (Onglet \textit{Bordures}) :} Choisir \textit{Style de trait} (fin, double, pointillé), \textit{Couleur}, puis cliquer sur les icônes de position (Extérieur, Intérieur, Diagonales) ou dessiner dans l'aperçu.
    \item \textbf{Règle pro :} Bordure \textbf{double trait fin} en bas des totaux ; trait \textbf{fin continu} pour la grille intérieure ; \textbf{pas de bordure} sur les cellules vides hors tableau.
    \item \textbf{Arrière-plan (Onglet \textit{Arrière-plan}) :} Couleur de fond pour les en-têtes (ex. \texttt{popblueL} ou bleu foncé avec police blanche) et la ligne des totaux (ex. \texttt{popgoldL} ou gris clair). Éviter les couleurs saturées qui fatiguent l'œil à l'impression.
\end{itemize}
\end{propriete}

\begin{propriete}[Formatage des nombres (Onglet \textit{Nombre}) : Adapter l'affichage au sens]
\begin{center}
\begin{tabularx}{\linewidth}{|l|X|X|}\hline
\rowcolor{popblueL}
\textbf{Catégorie} & \textbf{Usage typique (Boutique Kamga)} & \textbf{Exemple d'affichage} \\\hline
\textbf{Nombre} & Quantités, numéros de série & \texttt{12}, \texttt{1 234,56} (décimales paramétrables) \\\hline
\textbf{Monétaire} & Prix unitaires, Totaux, Chiffre d'affaires & \texttt{50 000 FCFA}, \texttt{1 250 000 FCFA} \\\hline
\textbf{Pourcentage} & Taux TVA (19,25\%), Marge, Remise & \texttt{19,25\%} (stocke 0,1925) \\\hline
\textbf{Date} & Date de vente, Date réception stock & \texttt{01/01/2024} (court) ou \texttt{Lundi 1 Janvier 2024} (long) \\\hline
\textbf{Heure} & Heure de vente (pour analyse pics) & \texttt{14:30} ou \texttt{14:30:05} \\\hline
\textbf{Personnalisé} & Formats spécifiques (ex. code article) & \texttt{TEL-001} (format \texttt{TEL-000}) \\\hline
\end{tabularx}
\end{center}
\end{propriete}

\begin{popfigure}
\begin{tikzpicture}[
    cell/.style={draw=popline, thin, minimum width=2.2cm, minimum height=0.7cm, anchor=north west, font=\small},
    header/.style={cell, fill=popblue, font=\bfseries\small, text=white},
    total/.style={cell, fill=popgold, font=\bfseries\small},
    title/.style={cell, fill=poppurple, font=\bfseries\small\Large, text=white, minimum width=13.2cm},
    empty/.style={cell, fill=white},
    node distance=0
]
% --- AVANT MISE EN FORME ---
\node[empty] (a1) at (0,0) {Journal de ventes};
\node[empty] (a2) at (2.2,0) {};
\node[empty] (a3) at (4.4,0) {};
\node[empty] (a4) at (6.6,0) {};
\node[empty] (a5) at (8.8,0) {};
\node[empty] (a6) at (11.0,0) {};
\node[empty] (b1) at (0,-0.7) {Date};
\node[empty] (b2) at (2.2,-0.7) {Modèle};
\node[empty] (b3) at (4.4,-0.7) {Qté};
\node[empty] (b4) at (6.6,-0.7) {Prix U};
\node[empty] (b5) at (8.8,-0.7) {Total};
\node[empty] (b6) at (11.0,-0.7) {TVA};
\node[empty] (c1) at (0,-1.4) {01/01/24};
\node[empty, align=right] (c2) at (2.2,-1.4) {Galaxy A14};
\node[empty, align=right] (c3) at (4.4,-1.4) {2};
\node[empty, align=right] (c4) at (6.6,-1.4) {120000};
\node[empty, align=right] (c5) at (8.8,-1.4) {240000};
\node[empty, align=right] (c6) at (11.0,-1.4) {46200};
\node[empty] (d1) at (0,-2.1) {02/01/24};
\node[empty, align=right] (d2) at (2.2,-2.1) {Tecno Spark 10};
\node[empty, align=right] (d3) at (4.4,-2.1) {1};
\node[empty, align=right] (d4) at (6.6,-2.1) {95000};
\node[empty, align=right] (d5) at (8.8,-2.1) {95000};
\node[empty, align=right] (d6) at (11.0,-2.1) {18287,5};
\node[empty] (e1) at (0,-2.8) {Total};
\node[empty] (e2) at (2.2,-2.8) {};
\node[empty, align=right] (e3) at (4.4,-2.8) {3};
\node[empty] (e4) at (6.6,-2.8) {};
\node[empty, align=right] (e5) at (8.8,-2.8) {335000};
\node[empty, align=right] (e6) at (11.0,-2.8) {64487,5};
\node[font=\bfseries, popdark] at (6.6, 0.5) {\textbf{AVANT : Brut, illisible, pas de distinction}};
% --- APRÈS MISE EN FORME (décalé vers le bas) ---
\begin{scope}[yshift=-5.5cm]
\node[title] (a1p) at (0,0) {JOURNAL DE VENTES - BOUTIQUE KAMGA - JANVIER 2024};
\node[header] (b1p) at (0,-0.7) {Date};
\node[header] (b2p) at (2.2,-0.7) {Modèle};
\node[header] (b3p) at (4.4,-0.7) {Qté};
\node[header] (b4p) at (6.6,-0.7) {Prix Unit. (FCFA)};
\node[header] (b5p) at (8.8,-0.7) {Total (FCFA)};
\node[header] (b6p) at (11.0,-0.7) {TVA (FCFA)};
\node[cell] (c1p) at (0,-1.4) {01/01/2024};
\node[cell, align=left] (c2p) at (2.2,-1.4) {Galaxy A14};
\node[cell, align=center] (c3p) at (4.4,-1.4) {2};
\node[cell, align=right] (c4p) at (6.6,-1.4) {120 000};
\node[cell, align=right] (c5p) at (8.8,-1.4) {240 000};
\node[cell, align=right] (c6p) at (11.0,-1.4) {46 200};
\node[cell] (d1p) at (0,-2.1) {02/01/2024};
\node[cell, align=left] (d2p) at (2.2,-2.1) {Tecno Spark 10};
\node[cell, align=center] (d3p) at (4.4,-2.1) {1};
\node[cell, align=right] (d4p) at (6.6,-2.1) {95 000};
\node[cell, align=right] (d5p) at (8.8,-2.1) {95 000};
\node[cell, align=right] (d6p) at (11.0,-2.1) {18 287,5};
\node[total] (e1p) at (0,-2.8) {TOTAL};
\node[total] (e2p) at (2.2,-2.8) {};
\node[total, align=center] (e3p) at (4.4,-2.8) {3};
\node[total] (e4p) at (6.6,-2.8) {};
\node[total, align=right] (e5p) at (8.8,-2.8) {335 000};
\node[total, align=right] (e6p) at (11.0,-2.8) {64 487,5};
% Draw double bottom border for total row
\draw[popdark, double, thick] (0, -3.5) -- (13.2, -3.5);
\end{scope}
\node[font=\bfseries, popgreen] at (6.6, -5.0) {\textbf{APRÈS : Titre fusionné, en-têtes colorés, alignements logiques, format FCFA, bordure double total}};
\end{tikzpicture}
\legende{Exemple « Journal de ventes » : comparaison Avant / Après mise en forme professionnelle.}
\end{popfigure}

\courssub{Ajustement, insertion/suppression et gestion des feuilles}

Un tableau vivant évolue : on ajoute des lignes, on élargit des colonnes, on organise par onglets.

\begin{methode}[Ajuster la largeur des colonnes et la hauteur des lignes]
\begin{enumerate}
    \item \textbf{Manuel :} Placer la souris sur le séparateur d'en-têtes (ex. entre \texttt{C} et \texttt{D}) $\to$ le curseur devient $\leftrightarrow$ $\to$ Glisser.
    \item \textbf{Auto-ajustement (Double-clic) :} Double-clic sur le même séparateur $\to$ La colonne s'adapte au contenu le plus large de la sélection (ou de la colonne si une seule cellule sélectionnée). \textbf{Le plus utilisé au quotidien.}
    \item \textbf{Hauteur de ligne :} Même principe sur les séparateurs horizontaux (numéros de lignes). Utile si \textit{Renvoi à la ligne auto} activé.
    \item \textbf{Format précis :} Clic droit sur en-tête $\to$ \textit{Largeur de colonne...} / \textit{Hauteur de ligne...} $\to$ Valeur en cm ou points.
\end{enumerate}
\end{methode}

\begin{propriete}[Gestion structurelle : Lignes, Colonnes, Feuilles]
\begin{itemize}
    \item \textbf{Insérer :} Clic droit sur en-tête (ligne/colonne) $\to$ \textit{Insérer} (au-dessus / à gauche). Raccourci : \texttt{Ctrl + +} (pavé numérique) ou \textit{Feuille} $\to$ \textit{Insérer des lignes/colonnes}.
    \item \textbf{Supprimer :} Clic droit $\to$ \textit{Supprimer}. \textbf{Attention :} Supprime le contenu \textbf{et} décale les cellules (les formules adaptent leurs références grâce aux références relatives, voir Section 2).
    \item \textbf{Effacer le contenu seulement :} Sélection $\to$ Touche \texttt{Suppr} (Del) $\to$ Choisir \textit{Tout}, \textit{Formats}, \textit{Commentaires}... Les cellules restent vides mais en place.
    \item \textbf{Feuilles (Onglets bas) :}
        \begin{itemize}
            \item \textbf{Renommer :} Double-clic sur l'onglet $\to$ Nom explicite (ex. \texttt{Janvier}, \texttt{Stock}, \texttt{Recap}).
            \item \textbf{Colorier l'onglet :} Clic droit sur onglet $\to$ \textit{Colorier l'onglet} (code couleur par mois/service).
            \item \textbf{Déplacer/Copier :} Glisser-déposer l'onglet (maintenir \texttt{Ctrl} pour copier).
            \item \textbf{Insérer/Supprimer feuille :} Clic droit sur onglet $\to$ \textit{Insérer} / \textit{Supprimer}.
        \end{itemize}
\end{itemize}
\end{propriete}

\begin{exoresolu}[Mise en forme du tableau de bord de M. Kamga]
\textbf{Énoncé :} M. Kamga a saisi ses ventes brutes de janvier dans la feuille \texttt{Feuille1} (plage \texttt{A1:F10}). Il veut : (1) Renommer l'onglet en \texttt{Janvier} couleur bleu. (2) Insérer une ligne en haut pour un titre fusionné sur 6 colonnes « VENTES JANVIER 2024 - BOUTIQUE KAMGA », police 16, gras, centré, fond bleu marine, police blanche. (3) Mettre en forme la ligne d'en-têtes (ligne 2 maintenant) : gras, fond bleu clair, bordures fines, alignement centré. (4) Formater colonnes \texttt{D, E, F} en Monétaire FCFA (0 décimale). (5) Auto-ajuster toutes les colonnes.

\textbf{Solution détaillée :}
\begin{enumerate}
    \item \textbf{Onglet :} Double-clic sur \texttt{Feuille1} $\to$ taper \texttt{Janvier} $\to$ \texttt{Entrée}. Clic droit onglet $\to$ \textit{Colorier l'onglet} $\to$ choisir bleu.
    \item \textbf{Titre :} Clic droit sur n° ligne 1 $\to$ \textit{Insérer des lignes au-dessus}. Sélection \texttt{A1:F1}. \texttt{Ctrl+1} $\to$ \textit{Alignement} : Cocher \textit{Fusionner}, Horizontal \textit{Centré}, Vertical \textit{Centré}. \textit{Police} : \textit{Liberation Sans}, \textit{Gras}, \textit{Taille 16}, \textit{Couleur blanche}. \textit{Arrière-plan} : Bleu marine (popdark). \textit{Bordures} : Aucune (ou contour fin blanc). Saisir le texte.
    \item \textbf{En-têtes (Ligne 2) :} Sélection \texttt{A2:F2}. \texttt{Ctrl+1} $\to$ \textit{Police} : Gras. \textit{Arrière-plan} : Bleu clair (popblueL). \textit{Bordures} : Style fin, Couleur popdark, bouton \textit{Extérieur} + \textit{Intérieur}. \textit{Alignement} : Horizontal Centré, Vertical Centré, \textit{Renvoi à la ligne auto} coché.
    \item \textbf{Format FCFA :} Sélection \texttt{D3:F10} (données). \texttt{Ctrl+1} $\to$ \textit{Nombre} : Catégorie \textit{Monétaire}. Format : symbole \texttt{FCFA}, séparateur milliers espace, Décimales : 0. Valider. (Équivalent \textit{Personnalisé} : \texttt{\#\#0 "FCFA"}).
    \item \textbf{Auto-ajustement :} Clic sur carreau d'angle (sélection tout) $\to$ Double-clic sur n'importe quel séparateur vertical d'en-têtes (A/B, B/C...). Toutes les colonnes s'adaptent.
\end{enumerate}
\textbf{Résultat :} Un tableau professionnel, prêt pour l'entrée des formules (Section 2) et l'impression (Section 6).
\end{exoresolu}

\begin{exercice}[Entraînement : « Inventaire Accessoires »]
Vous gérez le stock d'accessoires (coques, chargeurs, écouteurs) de la boutique.
\begin{enumerate}
    \item Ouvrir un nouveau classeur. Renommer \texttt{Feuille1} en \texttt{Stock\_Accessoires} (onglet vert).
    \item En \texttt{A1:E1}, saisir le titre « INVENTAIRE ACCESSOIRES - FEVRIER 2024 », fusionner, centrer, police 14, gras, fond orange, police blanche.
    \item En ligne 2, créer les en-têtes : \texttt{A2:Code}, \texttt{B2:Désignation}, \texttt{C2:Qté Stock}, \texttt{D2:Prix Achat (FCFA)}, \texttt{E2:Valeur Stock (FCFA)}. Mise en forme : gras, fond jaune clair, bordures, centré, renvoi à la ligne auto.
    \item Saisir 5 articles fictifs (codes \texttt{ACC001} à \texttt{ACC005}), quantités (entiers), prix achat (entiers FCFA).
    \item Formater colonnes \texttt{D} et \texttt{E} en Monétaire FCFA (0 décimale).
    \item Auto-ajuster toutes les colonnes. Enregistrer sous \texttt{Inventaire\_Kamga.ods}.
\end{enumerate}
\end{exercice}

\begin{corrige}[Exercice n°1]
\begin{enumerate}
    \item Onglet : Double-clic \texttt{Feuille1} $\to$ \texttt{Stock\_Accessoires} $\to$ Clic droit onglet $\to$ Colorier $\to$ Vert.
    \item Ligne 1 : Insérer ligne si besoin (déjà ligne 1). Sélection \texttt{A1:E1} $\to$ \texttt{Ctrl+1} $\to$ Alignement (Fusionner, Centré/Centré) $\to$ Police (14, Gras, Blanche) $\to$ Arrière-plan (Orange). Saisie titre.
    \item Ligne 2 : Saisie en-têtes. Sélection \texttt{A2:E2} $\to$ \texttt{Ctrl+1} $\to$ Police (Gras) $\to$ Arrière-plan (Jaune clair) $\to$ Bordures (Extérieur+Intérieur, fin) $\to$ Alignement (Centré, Centré, Renvoi auto).
    \item Saisie données lignes 3 à 7. Respecter types : Texte (Code, Désignation), Nombre (Qté), Monétaire (Prix, Valeur).
    \item Colonnes D:E : Sélection \texttt{D3:E7} (ou D:E entier) $\to$ \texttt{Ctrl+1} $\to$ Nombre $\to$ Monétaire $\to$ FCFA $\to$ Décimales 0.
    \item Carreau d'angle $\to$ Double-clic séparateur vertical. Fichier $\to$ Enregistrer sous $\to$ \texttt{Inventaire\_Kamga.ods}.
\end{enumerate}
\end{corrige}

\coursec{Formules, Opérateurs et Références de cellules}

Dans la section précédente, tu as appris à te repérer dans la fenêtre d'un tableur (barre de formule, cellule active, feuilles) et à mettre en forme ton classeur : couleurs, bordures, formats de nombre. Mais un tableur qui ne fait que \emph{présenter} des nombres n'est guère plus utile qu'une feuille de papier. Sa véritable puissance, c'est le \textbf{calcul automatique} : tu écris une fois une formule, et le logiciel calcule, recalcule et recopie tout seul, même sur mille lignes. C'est cette magie que nous allons maintenant maîtriser.

Imaginons la boutique de provisions de Mme Etoundi, à Mvog-Ada. Elle vend 40 articles différents et veut calculer pour chacun le montant total (prix unitaire $\times$ quantité), puis ajouter la TVA. Faut-il faire 80 calculs à la main ? Non : une seule formule, bien écrite, recopiée en quelques secondes, fait tout le travail. Encore faut-il comprendre \emph{comment} le tableur interprète ce que tu écris.

\courssub{Écrire une formule : le signe égal et les opérateurs}

\begin{definition}[Formule]
Une \cle{formule} est une expression saisie dans une cellule, qui demande au tableur d'effectuer un \textbf{calcul} et d'afficher le \textbf{résultat}. Toute formule commence \textbf{obligatoirement} par le signe égal \texttt{=}. C'est ce signe qui dit au logiciel : « attention, ce qui suit n'est pas du texte, c'est un calcul à effectuer ».
\end{definition}

Une formule peut contenir des \textbf{nombres} (constantes), des \textbf{références de cellules} (comme \texttt{B3}), des \textbf{opérateurs} et des \textbf{fonctions} (que nous verrons à la section 3). Les opérateurs arithmétiques disponibles sont :

\begin{center}
\begin{tabularx}{\linewidth}{|c|l|X|l|}
\hline
\rowcolor{popblueL} \textbf{Symbole} & \textbf{Opération} & \textbf{Exemple de formule} & \textbf{Résultat} \\
\hline
\texttt{+} & Addition & \texttt{=2500+1500} & 4000 \\
\hline
\texttt{-} & Soustraction & \texttt{=5000-1200} & 3800 \\
\hline
\texttt{*} & Multiplication & \texttt{=10*5000} & 50000 \\
\hline
\texttt{/} & Division & \texttt{=10000/4} & 2500 \\
\hline
\texttt{\^{}} & Puissance (exposant) & \texttt{=2\^{}10} & 1024 \\
\hline
\texttt{\%} & Pourcentage (divise par 100) & \texttt{=8000*15\%} & 1200 \\
\hline
\end{tabularx}
\end{center}

\begin{attention}[Le piège du signe égal oublié]
L'erreur la plus fréquente du débutant : taper \texttt{B3*C3} \textbf{sans} le signe \texttt{=}. Le tableur considère alors la saisie comme du \emph{texte} et affiche bêtement « B3*C3 » dans la cellule, sans rien calculer. Si ta cellule affiche exactement ce que tu as tapé, vérifie d'abord que la formule commence bien par \texttt{=}.
\end{attention}

\courssub{La priorité des opérations : dans quel ordre le tableur calcule-t-il ?}

Que vaut \texttt{=2+3*4} ? Si l'on calcule de gauche à droite, on obtient $5 \times 4 = 20$. Mais le tableur, comme en mathématiques, respecte des \textbf{priorités} : la multiplication passe avant l'addition, donc le résultat est $2 + 12 = 14$.

\begin{propriete}[Ordre de priorité des opérateurs]
Le tableur évalue une formule dans l'ordre suivant :
\begin{enumerate}
\item les \textbf{parenthèses} \texttt{( )} (les plus internes d'abord) ;
\item les \textbf{puissances} \texttt{\^{}} et le \textbf{pourcentage} \texttt{\%} ;
\item les \textbf{multiplications} \texttt{*} et \textbf{divisions} \texttt{/} ;
\item les \textbf{additions} \texttt{+} et \textbf{soustractions} \texttt{-}.
\end{enumerate}
À priorité égale (par exemple \texttt{*} et \texttt{/}), le calcul se fait \textbf{de gauche à droite}.
\end{propriete}

\begin{exemplebox}[Comparaison de deux formules voisines]
\begin{itemize}
\item \texttt{=A1+B1*C1} avec A1 $= 100$, B1 $= 10$, C1 $= 5$ donne $100 + (10 \times 5) = 150$.
\item \texttt{=(A1+B1)*C1} avec les mêmes valeurs donne $(100 + 10) \times 5 = 550$.
\end{itemize}
Les parenthèses changent tout ! Pour calculer une moyenne de deux notes, il faut écrire \texttt{=(B2+C2)/2} et non \texttt{=B2+C2/2} (qui diviserait seulement C2 par 2).
\end{exemplebox}

\begin{attention}[Parenthèses : ton meilleur outil de sécurité]
Quand tu hésites sur l'ordre de calcul, \textbf{mets des parenthèses}. Elles ne coûtent rien et rendent ta formule à la fois correcte et lisible. Retiens aussi que le tableur n'accepte que les parenthèses \texttt{( )} : pas de crochets ni d'accolades comme en mathématiques. Pour des parenthèses imbriquées, on écrit par exemple \texttt{=((A1+B1)*2-C1)/3}.
\end{attention}

\courssub{Les références de cellules : le vrai secret du tableur}

Revenons à la boutique de Mme Etoundi. En B3 se trouve la quantité vendue (10), en C3 le prix unitaire (5000 FCFA). Pour le montant, on peut écrire \texttt{=10*5000}... mais si demain le prix passe à 5500 FCFA, il faudra réécrire la formule ! La bonne pratique consiste à écrire \texttt{=B3*C3} : la formule va \emph{chercher} les valeurs dans les cellules. Si le contenu de C3 change, le résultat se met à jour \textbf{automatiquement}.

\begin{aretenir}[Règle d'or]
Dans une formule, utilise toujours des \cle{références de cellules} (\texttt{B3}, \texttt{C3}...) plutôt que des valeurs fixes tapées en dur. Ainsi, quand une donnée change, tous les calculs qui en dépendent se recalculent tout seuls. C'est le principe fondamental du tableur : \textbf{on sépare les données des calculs}.
\end{aretenir}

Mais une question se pose : si je recopie la formule \texttt{=B3*C3} vers le bas (pour l'article suivant, en ligne 4), que devient-elle ? Tout dépend du \emph{type} de référence utilisé.

\textbf{La référence relative (B3).} Par défaut, toute référence est \cle{relative} : elle désigne une cellule \emph{par sa position par rapport à la formule}. \texttt{=B3*C3} écrite en D3 signifie en réalité : « multiplie la cellule située 2 colonnes à gauche par celle située 1 colonne à gauche ». Recopiée en D4, elle devient donc \texttt{=B4*C4} ; recopiée en E5 (une colonne à droite, deux lignes plus bas), elle devient \texttt{=C5*D5}. C'est exactement ce qu'on veut pour une facture : chaque ligne calcule avec ses propres données.

\textbf{La référence absolue (\$B\$3).} Parfois, une valeur doit servir à \emph{tous} les calculs : le taux de TVA, par exemple, écrit une seule fois en F2. Si l'on recopie une formule contenant \texttt{F2} vers le bas, elle devient F3, F4... et pointe vers des cellules vides ! Pour « clouer » la référence sur place, on ajoute des dollars : \texttt{\$F\$2}. Le symbole \texttt{\$} devant la lettre verrouille la colonne, le \texttt{\$} devant le chiffre verrouille la ligne. Recopiée n'importe où, \texttt{\$F\$2} désigne toujours la cellule F2.

\textbf{Les références mixtes (\$B3 ou B\$3).} On peut ne verrouiller qu'une seule des deux coordonnées : \texttt{\$B3} fixe la colonne B mais laisse la ligne glisser ; \texttt{B\$3} fixe la ligne 3 mais laisse la colonne glisser. C'est l'outil idéal pour construire une table de multiplication ou une grille tarifaire, où l'on recopie une formule à la fois vers le bas \emph{et} vers la droite.

\begin{propriete}[Théorème de la recopie]
Une formule contenant des références relatives se transpose lors de la recopie en \textbf{conservant le décalage relatif} (la distance en lignes et en colonnes) entre la cellule de la formule et chaque cellule référencée. Toute partie de référence précédée du symbole \texttt{\$} (colonne et/ou ligne) reste quant à elle \textbf{inchangée} pendant la recopie.
\end{propriete}

\begin{methode}[Verrouiller une référence avec la touche F4]
Plutôt que de taper les dollars à la main :
\begin{enumerate}
\item double-clique sur la cellule contenant la formule (ou clique dans la barre de formule) ;
\item place le curseur sur la référence à modifier (par exemple \texttt{F2}) ;
\item appuie sur la touche \texttt{F4} (Windows) ou \texttt{Fn+F4} / \texttt{Cmd+T} (Mac selon version) ;
\item chaque pression fait défiler les quatre formes : \texttt{F2} $\rightarrow$ \texttt{\$F\$2} $\rightarrow$ \texttt{F\$2} $\rightarrow$ \texttt{\$F2} $\rightarrow$ \texttt{F2} ;
\item valide avec \texttt{Entrée} quand la forme voulue est affichée.
\end{enumerate}
\end{methode}

\begin{exemplebox}[Facture avec TVA : l'exemple de référence]
Mme Etoundi écrit le taux de TVA en F2 : 19,25 \%. En G3, elle calcule le prix TTC de l'article de la ligne 3 (prix HT en E3) :
\begin{lstlisting}
G3 : =E3*(1+$F$2)
\end{lstlisting}
Puis elle recopie G3 vers le bas jusqu'en G10. Grâce au théorème de la recopie :
\begin{itemize}
\item \texttt{E3} (relative) devient \texttt{E4}, \texttt{E5}, ..., \texttt{E10} : chaque ligne utilise son propre prix HT ;
\item \texttt{\$F\$2} (absolue) reste \texttt{\$F\$2} partout : toutes les lignes utilisent le même taux.
\end{itemize}
Si le taux de TVA change un jour, il suffit de modifier la seule cellule F2 : les huit montants TTC se recalculent instantanément.
\end{exemplebox}

\begin{popfigure}
\begin{tikzpicture}[scale=0.92, every node/.style={font=\small}]
% Grille source
\draw[popline, step=0.9] (0,0) grid (4.5,2.7);
\foreach \x/\l in {0.45/A, 1.35/B, 2.25/C, 3.15/D, 4.05/E}
  \node[above, popmuted] at (\x,2.7) {\l};
\foreach \y/\l in {2.25/2, 1.35/3, 0.45/4}
  \node[left, popmuted] at (0,\y) {\l};
\node[fill=popblueL, minimum width=0.85cm, minimum height=0.85cm] at (3.15,2.25) {\footnotesize =B2*C2};
\node[fill=popblueL, minimum width=0.85cm, minimum height=0.85cm] at (3.15,1.35) {\footnotesize =B3*C3};
\node[fill=popblueL, minimum width=0.85cm, minimum height=0.85cm] at (3.15,0.45) {\footnotesize =B4*C4};
\draw[-{Stealth}, very thick, popblue] (3.6,2.0) -- (3.6,0.7);
\node[popblue, right, align=left] at (4.6,1.35) {Recopie vers le bas\\ (références \textbf{relatives})\\ B2 $\to$ B3 $\to$ B4};
\node[popdark, above] at (2.25,3.1) {\textbf{Référence relative : tout glisse}};

% Grille absolue
\begin{scope}[xshift=10.2cm]
\draw[popline, step=0.9] (0,0) grid (4.5,2.7);
\foreach \x/\l in {0.45/E, 1.35/F, 2.25/G, 3.15/H, 4.05/I}
  \node[above, popmuted] at (\x,2.7) {\l};
\foreach \y/\l in {2.25/2, 1.35/3, 0.45/4}
  \node[left, popmuted] at (0,\y) {\l};
\node[fill=poporangeL, minimum width=0.85cm, minimum height=0.85cm] at (1.35,2.25) {\footnotesize 19,25\%};
\node[fill=popgreenL, minimum width=0.85cm, minimum height=0.85cm] at (2.25,2.25) {\footnotesize =E2*\$F\$2};
\node[fill=popgreenL, minimum width=0.85cm, minimum height=0.85cm] at (2.25,1.35) {\footnotesize =E3*\$F\$2};
\node[fill=popgreenL, minimum width=0.85cm, minimum height=0.85cm] at (2.25,0.45) {\footnotesize =E4*\$F\$2};
\draw[-{Stealth}, very thick, popgreen] (2.7,2.0) -- (2.7,0.7);
\draw[-{Stealth}, dashed, thick, poporange] (1.35,1.9) -- (1.35,0.8);
\node[popdark, above] at (2.25,3.1) {\textbf{Référence absolue : \$F\$2 reste fixe}};
\end{scope}
\end{tikzpicture}
\legende{Lors de la recopie vers le bas, les références relatives glissent avec la formule (à gauche), tandis que la référence absolue \$F\$2 (à droite, en orange) pointe toujours vers la même cellule.}
\end{popfigure}

\courssub{Recopier une formule : la poignée de recopie}

Comment recopier concrètement une formule sur 10, 100 ou 1000 lignes ? Inutile de tout retaper : le tableur offre la \cle{poignée de recopie}.

\begin{methode}[Utiliser la poignée de recopie]
\begin{enumerate}
\item clique sur la cellule contenant la formule (elle devient la cellule active, entourée d'un cadre) ;
\item repère le \textbf{petit carré noir} situé au coin inférieur droit du cadre : c'est la poignée de recopie ;
\item place le pointeur de la souris dessus : il se transforme en fine croix noire \texttt{+} ;
\item \textbf{glisse} vers le bas (ou vers la droite) jusqu'à la dernière cellule voulue, puis relâche ;
\item astuce de pro : si une colonne voisine est déjà remplie, un \textbf{double-clic} sur la poignée recopie automatiquement la formule jusqu'en bas du tableau.
\end{enumerate}
\end{methode}

\begin{experience}[TP guidé : la facture de la boutique]
\textbf{Objectif :} construire une mini-facture avec recopie de formules et taux de TVA verrouillé.\\
\textbf{Matériel :} un ordinateur avec Microsoft Excel ou LibreOffice Calc.
\begin{enumerate}
\item Ouvre un classeur neuf. Saisis : en A1 « Article », B1 « Qté », C1 « Prix unit. », D1 « Montant HT », F1 « TVA » et en F2 la valeur 19,25\%.
\item Remplis trois lignes d'articles : (Cahier, 10, 500), (Stylo, 25, 100), (Règle, 5, 300).
\item En D3, tape la formule \texttt{=B3*C3} et valide. Vérifie le résultat : 5000.
\item Recopie D3 vers D4 et D5 avec la poignée de recopie. Clique sur D4 : la barre de formule affiche \texttt{=B4*C4}. Le théorème de la recopie est confirmé !
\item En E3, tape \texttt{=D3*(1+\$F\$2)} puis recopie vers le bas. Vérifie que \$F\$2 n'a pas bougé.
\item Modifie maintenant le prix du cahier (C3 passe à 600). Observe : D3 et E3 se recalculent instantanément.
\end{enumerate}
\textbf{Interprétation :} les références relatives suivent la recopie, la référence absolue reste fixe, et toute modification d'une donnée se propage automatiquement à tous les calculs qui en dépendent.
\end{experience}

\courssub{Comprendre et corriger les erreurs de calcul}

Quand une formule est incorrecte, le tableur ne se fâche pas : il affiche un \textbf{code d'erreur} qui commence par le symbole \texttt{\#}. Apprends à les lire comme un médecin lit des symptômes.

\begin{center}
\begin{tabularx}{\linewidth}{|l|X|X|}
\hline
\rowcolor{popblueL} \textbf{Code} & \textbf{Cause probable} & \textbf{Solution} \\
\hline
\texttt{\#DIV/0!} & Division par zéro ou par une cellule vide (ex. \texttt{=A1/B1} avec B1 vide). & Vérifier le diviseur ; le compléter ou empêcher la division si la cellule est vide. \\
\hline
\texttt{\#VALEUR!} & La formule utilise du texte là où un nombre est attendu (ex. \texttt{=A1*2} avec A1 contenant « dix »). & Corriger le contenu de la cellule : saisir un nombre, pas du texte. \\
\hline
\texttt{\#REF!} & La formule pointe vers une cellule qui a été supprimée (ligne ou colonne effacée). & Réécrire la formule avec une référence valide ; annuler la suppression si possible (Ctrl+Z). \\
\hline
\texttt{\#NOM?} & Nom de fonction mal orthographié ou nom de plage inconnu (ex. \texttt{=SOMME(A1;B1)} au lieu de \texttt{=SOMME(A1;B1)} si la fonction est \texttt{SOMME} mais tapé \texttt{SOM}). & Corriger l'orthographe du nom de la fonction. \\
\hline
\texttt{\#\#\#\#} & La colonne est trop étroite pour afficher le nombre (ce n'est pas une erreur de calcul !). & Élargir la colonne (double-clic sur le bord droit de l'en-tête de colonne). \\
\hline
\end{tabularx}
\end{center}

\begin{exoresolu}[Construire une grille de prix avec références mixtes]
\textbf{Énoncé.} Un cybercafé de Bafoussam facture la connexion selon la durée. En A3:A5 figurent les durées (1 h, 2 h, 3 h) et en B2:D2 les tarifs horaires selon la tranche de la journée (300, 250, 200 FCFA). On veut remplir la grille B3:D5 où chaque cellule contient durée $\times$ tarif, en n'écrivant \textbf{qu'une seule formule} en B3, recopiée ensuite sur toute la grille. Quelle formule faut-il écrire ?
\tcblower
\textbf{Solution détaillée.} En B3, le calcul est \texttt{=A3*B2}. Réfléchissons à ce qui doit bouger lors de la recopie :
\begin{itemize}
\item la durée est toujours lue en \textbf{colonne A} : la colonne ne doit jamais changer, mais la ligne doit glisser quand on recopie vers le bas $\Rightarrow$ on écrit \texttt{\$A3} ;
\item le tarif est toujours lu en \textbf{ligne 2} : la ligne ne doit jamais changer, mais la colonne doit glisser quand on recopie vers la droite $\Rightarrow$ on écrit \texttt{B\$2}.
\end{itemize}
La formule unique à saisir en B3 est donc :
\begin{lstlisting}
B3 : =$A3*B$2
\end{lstlisting}
Recopiée vers le bas puis vers la droite (ou l'inverse), elle devient par exemple \texttt{=\$A5*D\$2} en D5, ce qui donne $3 \times 200 = 600$ FCFA. Toute la grille (9 cellules) est remplie à partir d'une seule formule : c'est toute la puissance des références mixtes.
\end{exoresolu}

\begin{savaistu}[Le symbole @ dans les formules modernes]
Dans les versions récentes d'Excel, lorsque tes données sont formatées comme un \emph{tableau structuré} (avec des colonnes nommées), tu peux rencontrer des formules comme \texttt{=[@Quantité]*[@Prix]}. Le symbole \texttt{@} signifie « la valeur de cette colonne, sur la ligne courante » : c'est une \textbf{référence implicite}. Elle rend les formules plus lisibles — on lit des mots au lieu de B3 et C3 — mais le principe reste exactement celui que tu viens d'apprendre : référencer des cellules plutôt que taper des valeurs.
\end{savaistu}

\begin{exercice}[À toi de jouer]
Le bulletin trimestriel d'un élève de 3ème contient en B2, C2 et D2 ses trois notes d'informatique (12, 15 et 9 sur 20), et en E2 le total des coefficients (4).
\begin{enumerate}
\item Écris en F2 la formule qui calcule la moyenne de l'élève, en utilisant uniquement des références de cellules et en respectant les priorités.
\item Cette formule est recopiée vers le bas pour 30 élèves. Les références doivent-elles être relatives ou absolues ? Justifie.
\item En G2, on veut écrire l'appréciation « Admis » si la moyenne est au moins 10. Quelle difficulté prévois-tu, sachant que nous n'avons pas encore vu les fonctions logiques ? (Réponse à la section 4 !)
\end{enumerate}
\end{exercice}

\begin{corrige}[Exercice : bulletin de notes]
\begin{enumerate}
\item \texttt{=(B2+C2+D2)/E2} — les parenthèses sont indispensables pour additionner avant de diviser.
\item Relatives : chaque élève a ses propres notes sur sa propre ligne, donc B2, C2, D2, E2 doivent devenir B3, C3... lors de la recopie. Aucun dollar n'est nécessaire.
\item Il faudrait un calcul qui dépend d'une \emph{condition} (moyenne $\geq 10$) : une simple formule arithmétique ne suffit pas ; il faudra la fonction logique \texttt{SI}.
\end{enumerate}
\end{corrige}

\begin{aretenir}[L'essentiel de la section]
\begin{itemize}
\item Toute formule commence par \texttt{=} et respecte les priorités : parenthèses, puissances, puis $\times$ et $\div$, puis $+$ et $-$.
\item On référence les cellules plutôt que de taper des valeurs : les calculs se mettent alors à jour automatiquement.
\item Référence \textbf{relative} (\texttt{B3}) : glisse à la recopie. Référence \textbf{absolue} (\texttt{\$B\$3}) : reste fixe. Référence \textbf{mixte} (\texttt{\$B3} ou \texttt{B\$3}) : verrouille une seule coordonnée. La touche \texttt{F4} (ou raccourci Mac dédié) fait basculer entre les formes.
\item La \textbf{poignée de recopie} (petit carré noir) permet de dupliquer une formule par glisser ou double-clic.
\item Les codes \texttt{\#DIV/0!}, \texttt{\#VALEUR!}, \texttt{\#REF!}, \texttt{\#NOM?} et \texttt{\#\#\#\#} indiquent respectivement une division par zéro, du texte dans un calcul, une référence supprimée, un nom de fonction inconnu et une colonne trop étroite.
\end{itemize}
\end{aretenir}

\coursec{Fonctions mathématiques et statistiques usuelles}

Après avoir maîtrisé l'écriture de formules simples et la gestion des références relatives, absolues et mixtes, nous franchissons une étape décisive : l'utilisation des \textbf{fonctions prédéfinies}. Une fonction est une formule « toute faite », programmée par les concepteurs du tableur, qui effectue un calcul complexe à partir de valeurs que vous lui fournissez. Elle vous évite d'écrire des formules interminables comme \texttt{=B3+B4+B5+...+B50} et réduit drastiquement le risque d'erreur. Dans cette section, nous explorons les fonctions indispensables pour analyser des données numériques : sommes, moyennes, effectifs, extrema, comptages conditionnels et arrondis.

\courssub{Concept de fonction et Assistant Fonction (fx)}

Une fonction se reconnaît à sa syntaxe normalisée : un \textbf{nom} (en majuscules par convention), suivi de \textbf{parenthèses} contenant un ou plusieurs \textbf{arguments} séparés par un \textbf{point-virgule (;)}.

\begin{definition}[Argument d'une fonction]
Une \cle{valeur} fournie à une fonction pour qu'elle opère. Un argument peut être :
\begin{itemize}
    \item une \cle{constante} numérique ou textuelle (ex. \texttt{10}, \texttt{"Vendu"}),
    \item une \cle{référence de cellule} (ex. \texttt{B3}),
    \item une \cle{plage de cellules} (ex. \texttt{B3:B10}),
    \item une \cle{expression} logique (ex. \texttt{">=100"}),
    \item une \cle{autre fonction} (imbrication).
\end{itemize}
\end{definition}

\begin{propriete}[Syntaxe générale d'une fonction]
\[
\texttt{=NOM\_FONCTION(argument1 ; argument2 ; ...)}
\]
\begin{itemize}
    \item Le signe \texttt{=} indique au tableur qu'il s'agit d'un calcul.
    \item Le \textbf{point-virgule (;)} est le séparateur d'arguments standard dans la locale française (virgule décimale). \emph{Attention : dans les locales anglo-saxonnes, c'est la virgule (,) qui sépare les arguments.}
    \item Les arguments \emph{obligatoires} sont sans crochets ; les arguments \emph{optionnels} sont entourés de crochets \texttt{[...]} dans la documentation.
\end{itemize}
\end{propriete}

Pour découvrir les fonctions, leurs arguments et obtenir de l'aide contextuelle, on utilise l'\textbf{Assistant Fonction} (bouton \texttt{fx} à gauche de la barre de formule ou raccourci \texttt{Maj + F3}).

\begin{popfigure}
\begin{tikzpicture}[
    node distance=0.8cm and 1cm,
    box/.style={draw=popblue, thick, rounded corners, fill=popblueL, text width=5cm, align=center, minimum height=1.2cm},
    boxarg/.style={draw=poporange, thick, rounded corners, fill=poporangeL, text width=4.5cm, align=left, minimum height=1cm},
    arrow/.style={-Stealth, thick, popdark}
]
% Main dialog
\node[box] (title) {Assistant Fonction (fx)};
\node[box, below=of title, text width=10cm, minimum height=4cm] (main) {
    \textbf{Rechercher une fonction :} \texttt{SOMME.SI} \textit{(zone de saisie)}\\[0.3cm]
    \textbf{Ou sélectionner une catégorie :} \textit{Liste déroulante : Mathématiques et trigonométrie, Statistiques, Logique, Texte, Date/Heure, Recherche et matrices...}\\[0.3cm]
    \textbf{Sélectionner une fonction :} \textit{Liste filtrée : SOMME, MOYENNE, NB, NBVAL, NB.SI, SOMME.SI, MAX, MIN, ARRONDI...}\\[0.3cm]
    \textit{Description : « Additionne les cellules spécifiées par un critère donné. »}
};
% Arguments panel
\node[boxarg, right=of main, xshift=1cm, minimum height=4cm, align=center] (args){
    \textbf{Arguments de SOMME.SI :}\\[0.2cm]
    \texttt{Plage\_test} \textit{(obligatoire)}\\
    \textit{Plage de cellules à évaluer selon le critère.}\\[0.2cm]
    \texttt{Critère} \textit{(obligatoire)}\\
    \textit{Condition sous forme de nombre, expression ou texte.}\\[0.2cm]
    \texttt{[Somme\_plage]} \textit{(optionnel)}\\
    \textit{Cellules réelles à additionner. Si omis, = Plage\_test.}
};
\draw[arrow] (main.east) -- (args.west) node[midway, above, font=\small] {Sélection};
\node[draw=popgreen, thick, rounded corners, fill=popgreenL, below=of main, text width=10cm, align=center, minimum height=1cm] (result) {
    \textbf{Résultat de la fonction :} \texttt{=SOMME.SI(C3:C10; ">50000")} \quad \textit{Valeur calculée : 320\,000}\\
    \textit{Boutons : OK \quad Annuler \quad Aide sur cette fonction}
};
\draw[arrow] (main.south) -- (result.north);
\end{tikzpicture}
\legende{Boîte de dialogue « Assistant Fonction » pour SOMME.SI : recherche, description, arguments détaillés et prévisualisation du résultat.}
\end{popfigure}

\courssub{Fonctions de base : SOMME, MOYENNE, MAX, MIN, NB, NBVAL}

Ces fonctions sont les piliers de toute analyse statistique élémentaire. Elles ignorent généralement le texte et les cellules vides (sauf NBVAL).

\begin{methode}[Insérer SOMME rapidement (Somme auto)]
\begin{enumerate}
    \item Sélectionnez la cellule \textbf{juste en dessous} d'une colonne de nombres (ou à droite d'une ligne).
    \item Appuyez sur le raccourci \texttt{Alt + =} (ou cliquez sur le bouton \texttt{$\Sigma$ Somme auto} dans l'onglet \textit{Accueil} ou \textit{Formules}).
    \item Le tableur propose une plage (entourée de pointillés). Vérifiez qu'elle est correcte.
    \item Validez par \texttt{Entrée}.
\end{enumerate}
\end{methode}

\begin{exemplebox}[Tableau de notes d'une classe de 3ème]
Considérons les notes sur 20 de 8 élèves en colonne \texttt{B} (lignes 3 à 10).
\begin{center}
\begin{tabularx}{\linewidth}{|l|X|X|}\hline
\textbf{Cellule} & \textbf{Formule} & \textbf{Résultat / Signification} \\\hline
\texttt{B11} & \texttt{=SOMME(B3:B10)} & Somme totale des notes (ex. 112). \\\hline
\texttt{B12} & \texttt{=MOYENNE(B3:B10)} & Moyenne de la classe (ex. 14). \\\hline
\texttt{B13} & \texttt{=MAX(B3:B10)} & Meilleure note (ex. 19). \\\hline
\texttt{B14} & \texttt{=MIN(B3:B10)} & Note la plus faible (ex. 8). \\\hline
\texttt{B15} & \texttt{=NB(B3:B10)} & \textbf{Effectif} : nombre de notes \emph{numériques} (ex. 8). \\\hline
\texttt{B16} & \texttt{=NBVAL(B3:B10)} & Nombre de cellules \emph{non vides} (texte + nombres). \\\hline
\end{tabularx}
\end{center}
Si \texttt{B5} contient \texttt{"Absent"} (texte) : \texttt{NB} renverra 7 (ignore le texte), \texttt{NBVAL} renverra 8 (compte la cellule non vide), \texttt{MOYENNE} calculera sur les 7 notes numériques.
\end{exemplebox}

\begin{attention}[Pièges classiques sur les fonctions de base]
\begin{itemize}
    \item \textbf{Référence circulaire :} Écrire \texttt{=SOMME(B3:B11)} en \texttt{B11} inclut la cellule de total dans le calcul $\rightarrow$ Erreur ou boucle infinie. \textbf{Toujours arrêter la plage avant la cellule de résultat.}
    \item \textbf{Texte dans une plage numérique :} \texttt{MOYENNE} et \texttt{SOMME} ignorent le texte (traité comme 0 pour la somme, exclu du diviseur pour la moyenne). Mais \texttt{NB} ne compte \emph{que} les nombres. Si vous voulez l'effectif réel (élèves présents + absents), utilisez \texttt{NBVAL} sur la colonne des noms ou \texttt{NB.SI} avec critère \texttt{"<>""} (non vide).
    \item \textbf{Séparateur d'arguments :} En France/Cameroun (virgule décimale), c'est le \textbf{point-virgule (;)}. La virgule (,) sert à séparer les milliers dans les nombres (ex. \texttt{1 000 000}) ou les décimales en notation US.
\end{itemize}
\end{attention}

\courssub{Fonctions de comptage conditionnel : NB.SI}

\cle{NB.SI} compte le nombre de cellules dans une plage qui satisfont un \textbf{critère unique}. C'est l'outil premier pour extraire des indicateurs (taux de réussite, nombre de ventes > seuil, effectif par catégorie).

\begin{propriete}[Syntaxe de NB.SI]
\[
\texttt{=NB.SI(plage\_test ; critère)}
\]
\begin{itemize}
    \item \texttt{plage\_test} : cellule(s) à examiner (ex. \texttt{C3:C100}).
    \item \texttt{critère} : condition à respecter. Syntaxes acceptées :
    \begin{itemize}
        \item Nombre : \texttt{10} (égalité stricte).
        \item Texte : \texttt{"Vendu"} (guillemets obligatoires). Insensible à la casse.
        \item Expression logique : \texttt{">=10"}, \texttt{"<0"}, \texttt{"<>""} (non vide), \texttt{"*abc*"} (contient « abc » — \texttt{*} = joker).
        \item Référence de cellule : \texttt{E1} (la cellule E1 contient la valeur seuil, ex. 10).
    \end{itemize}
\end{itemize}
\end{propriete}

\begin{exemplebox}[Analyse de ventes d'un cybercafé (Yaoundé)]
Colonne \texttt{C} : Montants journaliers (FCFA) sur 30 jours (\texttt{C3:C32}).
\begin{itemize}
    \item Jours avec recette $\ge$ 50\,000 FCFA : \texttt{=NB.SI(C3:C32; ">=50000")}.
    \item Jours avec recette exactement 0 (fermeture) : \texttt{=NB.SI(C3:C32; 0)}.
    \item Jours « Pleine affluence » (texte en col. D) : \texttt{=NB.SI(D3:D32; "Pleine affluence")}.
    \item Jours non vides (activité enregistrée) : \texttt{=NB.SI(C3:C32; "<>""")}.
\end{itemize}
\end{exemplebox}

\begin{popfigure}
\begin{tikzpicture}[
    node distance=0.6cm and 0.5cm,
    cell/.style={draw=popline, minimum width=1.2cm, minimum height=0.7cm, anchor=north west, fill=white},
    header/.style={cell, fill=popblue, text=white, font=\bfseries},
    formula/.style={cell, fill=popcream, font=\ttfamily\small},
    result/.style={cell, fill=popgreenL, font=\bfseries},
    label/.style={font=\small, text=popdark, anchor=east}
]
% Grid headers
\node[header] (A1) at (0,0) {A};
\node[header] (B1) at (1.2,0) {B};
\node[header] (C1) at (2.4,0) {C};
\node[header] (D1) at (3.6,0) {D};
\node[header] (E1) at (4.8,0) {E};
% Row 2: Titles
\node[cell, fill=popblueL] (A2) at (0,-0.7) {Élève};
\node[cell, fill=popblueL] (B2) at (1.2,-0.7) {Note /20};
\node[cell, fill=popblueL] (C2) at (2.4,-0.7) {Statut};
\node[cell, fill=popblueL] (D2) at (3.6,-0.7) {Critère};
\node[cell, fill=popblueL] (E2) at (4.8,-0.7) {Résultat};
% Data rows
\foreach \i/\name/\note/\stat in {3/Amina/14/Admis, 4/Boris/8/Recalé, 5/Chantal/16/Admis, 6/David/10/Admis, 7/Elie/7/Recalé, 8/Fatou/12/Admis, 9/Gaston/9/Recalé, 10/Hélène/15/Admis} {
    \pgfmathsetmacro{\y}{-0.7 - (\i-2)*0.7}
    \node[cell] at (0,\y) {\name};
    \node[cell] at (1.2,\y) {\note};
    \node[cell] at (2.4,\y) {\stat};
}
% Criteria & Formulas
\node[label] at (-0.2, -0.7 - 8*0.7 - 0.3) {Seuil :};
\node[cell] at (1.2, -0.7 - 8*0.7 - 0.35) {10};
\node[formula] at (3.6, -0.7 - 8*0.7 - 0.35) {=">=" \& B11};
\node[formula] at (3.6, -0.7 - 9*0.7 - 0.35) {"Admis"};
\node[formula] at (3.6, -0.7 - 10*0.7 - 0.35) {"Recalé"};
% Results
\node[result] at (4.8, -0.7 - 8*0.7 - 0.35) {6};
\node[result] at (4.8, -0.7 - 9*0.7 - 0.35) {5};
\node[result] at (4.8, -0.7 - 10*0.7 - 0.35) {3};
% Formulas in E (small font below result)
\node[formula, font=\ttfamily\tiny, anchor=north] at (4.8, -0.7 - 8*0.7 - 0.35 + 0.15) {=NB.SI(B3:B10;E11)};
\node[formula, font=\ttfamily\tiny, anchor=north] at (4.8, -0.7 - 9*0.7 - 0.35 + 0.15) {=NB.SI(C3:C10;E12)};
\node[formula, font=\ttfamily\tiny, anchor=north] at (4.8, -0.7 - 10*0.7 - 0.35 + 0.15) {=NB.SI(C3:C10;E13)};
% Braces replaced by double arrows (calc library)
\draw[popblue, thick, <->] ($(1.2,-0.7) + (-0.3,0)$) -- ($(1.2,-6.3) + (-0.3,0)$) node[midway, left=3mm, popblue, font=\small] {Plage B3:B10};
\draw[poporange, thick, <->] ($(2.4,-0.7) + (0.3,0)$) -- ($(2.4,-6.3) + (0.3,0)$) node[midway, right=3mm, poporange, font=\small] {Plage C3:C10};
\end{tikzpicture}
\legende{Exemple complet : calcul de la moyenne, max, min, effectif (NB), et comptages conditionnels (NB.SI) sur notes et statuts.}
\end{popfigure}

\courssub{Fonction de somme conditionnelle : SOMME.SI}

Alors que \texttt{NB.SI} \textbf{compte}, \texttt{SOMME.SI} \textbf{additionne} des valeurs correspondant à un critère. La distinction subtile réside dans l'argument optionnel \texttt{somme\_plage}.

\begin{propriete}[Syntaxe de SOMME.SI]
\[
\texttt{=SOMME.SI(plage\_test ; critère ; [somme\_plage])}
\]
\begin{itemize}
    \item \texttt{plage\_test} : cellules évaluées par le critère (même syntaxe que NB.SI).
    \item \texttt{critère} : condition (nombre, texte, expression, référence).
    \item \texttt{somme\_plage} (\textbf{optionnel}) : cellules \textbf{réelles à additionner}. Si omis, \texttt{plage\_test} est additionnée.
\end{itemize}
\end{propriete}

\begin{aretenir}[Distinction cruciale NB.SI vs SOMME.SI]
\begin{center}
\begin{tabularx}{\linewidth}{|l|X|X|}\hline
\textbf{Fonction} & \textbf{Question posée} & \textbf{Plage additionnée} \\\hline
\texttt{NB.SI(A:A; "Vendu")} & Combien de lignes ont « Vendu » en col. A ? & \textit{Aucune (compte seulement)} \\\hline
\texttt{SOMME.SI(A:A; "Vendu")} & Somme des valeurs \textbf{en col. A} où c'est « Vendu ». & \texttt{A:A} (identique à test) \\\hline
\texttt{SOMME.SI(A:A; "Vendu"; B:B)} & Somme des valeurs \textbf{en col. B} où col. A = « Vendu ». & \texttt{B:B} (différente de test) \\\hline
\end{tabularx}
\end{center}
\end{aretenir}

\begin{exemplebox}[Chiffre d'affaires par catégorie (Boutique « Chez Maman »)]
\begin{center}
\begin{tabularx}{\linewidth}{|l|X|X|}\hline
\textbf{Colonne} & \textbf{Contenu} & \textbf{Exemple} \\\hline
A & Produit & « Pain », « Lait », « Savon » \\\hline
B & Catégorie & « Alimentaire », « Hygiène » \\\hline
C & Quantité vendue & 50, 20, 15 \\\hline
D & Prix unitaire & 150, 500, 300 \\\hline
E & Montant (C*D) & 7500, 10000, 4500 \\\hline
\end{tabularx}
\end{center}
\begin{itemize}
    \item CA total « Alimentaire » : \texttt{=SOMME.SI(B2:B100; "Alimentaire"; E2:E100)}.
    \item Quantité totale « Pain » : \texttt{=SOMME.SI(A2:A100; "Pain"; C2:C100)} (somme\_plage = Quantité).
    \item CA total si montant > 5000 : \texttt{=SOMME.SI(E2:E100; ">5000")} (somme\_plage omise = Montant).
\end{itemize}
\end{exemplebox}

\begin{exoresolu}[Analyse des notes de 3ème B]
{Dans une feuille, les notes (0 à 20) sont en \texttt{B3:B22}. La moyenne de la classe est en \texttt{B24}. On veut :\newline
1. L'effectif total (NB).\newline
2. Le nombre de notes $\ge$ moyenne de la classe.\newline
3. La somme des notes $\ge$ 10 (total points « acquis »).\newline
4. La moyenne des notes $\ge$ 10 (moyenne des admis).}
\tcblower
\textbf{Solution détaillée :}
\begin{enumerate}
    \item \texttt{=NB(B3:B22)} $\rightarrow$ 20 (si toutes remplies numériquement).
    \item \texttt{=NB.SI(B3:B22; ">=" \& B24)}. \textbf{Astuce :} Concaténer l'opérateur et la référence par \texttt{\&}. Si B24=11,5, le critère devient \texttt{">=11,5"}.
    \item \texttt{=SOMME.SI(B3:B22; ">=10")}. Ici \texttt{somme\_plage} omis : on additionne les notes elles-mêmes.
    \item \textbf{Piège :} \texttt{MOYENNE.SI} n'existe pas en version de base (apparue plus tard). On calcule : \texttt{=SOMME.SI(B3:B22; ">=10") / NB.SI(B3:B22; ">=10")}.
\end{enumerate}

\end{exoresolu}

\courssub{Calculs de pourcentages et évolutions}

Le tableur ne possède pas de fonction « POURCENTAGE ». Un pourcentage est une \textbf{division} mise en forme.

\begin{propriete}[Règles du pourcentage au tableur]
\begin{enumerate}
    \item \textbf{Part / Total} : \texttt{=Valeur\_Part / Valeur\_Total} $\rightarrow$ Appliquer le format \textbf{\%} (bouton \texttt{\%} ou \texttt{Ctrl+Maj+\%}). Le tableur affiche \texttt{0,25} comme \texttt{25\%}.
    \item \textbf{Évolution (Taux de variation)} : \texttt{=(Nouveau - Ancien) / Ancien} $\rightarrow$ Format \%. Positif = hausse, Négatif = baisse.
    \item \textbf{Application d'un taux} : \texttt{=Montant * Taux} (ex. TVA 19,25\% : \texttt{=Montant * 19,25\%} ou \texttt{=Montant * 0,1925}).
\end{enumerate}
\end{propriete}

\begin{exemplebox}[Bulletin de paie simplifié (Salaire de base : 250\,000 FCFA)]
\begin{center}
\begin{tabularx}{\linewidth}{|l|X|X|}\hline
\textbf{Ligne} & \textbf{Libellé} & \textbf{Formule / Valeur} \\\hline
3 & Salaire de base & 250000 \\\hline
4 & Prime ancienneté (5\%) & \texttt{=B3*5\%} $\rightarrow$ 12\,500 \\\hline
5 & Brut imposable & \texttt{=B3+B4} $\rightarrow$ 262\,500 \\\hline
6 & CNPS (4,2\%) & \texttt{=B5*4,2\%} $\rightarrow$ 11\,025 \\\hline
7 & Net à payer & \texttt{=B5-B6} $\rightarrow$ 251\,475 \\\hline
8 & \textbf{Évolution vs an dernier (230\,000)} & \texttt{=(B3-230000)/230000} $\rightarrow$ \textbf{8,70\%} \\\hline
\end{tabularx}
\end{center}
\textbf{Conseil pro :} Mettez le taux (5\%, 4,2\%, 19,25\%) dans une cellule dédiée (ex. \texttt{\$F\$1}) et référencez-la en absolu (\texttt{=\$F\$1}) pour modifier le taux globalement.
\end{exemplebox}

\courssub{Fonctions d'arrondi : ARRONDI, ARRONDI.SUP, ARRONDI.INF, ENT}

En comptabilité et gestion, l'affichage formaté (ex. 2 décimales) ne change pas la \textbf{valeur réelle} stockée (ex. 10,0049 affiché 10,00). Pour forcer la valeur calculée, on utilise les fonctions d'arrondi.

\begin{propriete}[Syntaxe des arrondis]
\[
\texttt{=ARRONDI(nombre ; nb\_chiffres)} \quad
\texttt{=ARRONDI.SUP(nombre ; nb\_chiffres)} \quad
\texttt{=ARRONDI.INF(nombre ; nb\_chiffres)} \quad
\texttt{=ENT(nombre)}
\]
\begin{itemize}
    \item \texttt{nombre} : valeur ou référence (ex. \texttt{B3}, \texttt{10/3}).
    \item \texttt{nb\_chiffres} : nombre de décimales à conserver.
    \begin{itemize}
        \item $>0$ : décimales (ex. 2 $\rightarrow$ centièmes).
        \item $=0$ : entier le plus proche (ex. \texttt{ARRONDI(3,14;0)=3}).
        \item $<0$ : arrondi aux dizaines, centaines... (ex. -1 $\rightarrow$ dizaine : \texttt{ARRONDI(123;-1)=120}).
    \end{itemize}
    \item \texttt{ARRONDI.SUP} : s'éloigne de zéro (vers $+\infty$ pour positifs, $-\infty$ pour négatifs).
    \item \texttt{ARRONDI.INF} : se rapproche de zéro (tronque).
    \item \texttt{ENT} : partie entière (vers $-\infty$). \texttt{ENT(3,9)=3}, \texttt{ENT(-3,9)=-4}.
\end{itemize}
\end{propriete}

\begin{attention}[Affichage vs Valeur réelle — Erreur majeure]
\begin{lstlisting}[language=, caption={Exemple de divergence affichage/valeur}]
A1 = 10/3          -> Valeur réelle : 3,333333...  Affichage (2 déc.) : 3,33
B1 = ARRONDI(A1;2) -> Valeur réelle : 3,33          Affichage (2 déc.) : 3,33
C1 = A1 * 3        -> Valeur réelle : 10,000000...  Affichage : 10,00
D1 = B1 * 3        -> Valeur réelle : 9,99          Affichage : 9,99  <-- ECART !
\end{lstlisting}
\textbf{Règle d'or :} En comptabilité, \textbf{arrondissez systématiquement les résultats intermédiaires} (TVA, totaux lignes) avec \texttt{ARRONDI(...,2)} avant de les sommer, sinon la somme des lignes $\neq$ total du bas.
\end{attention}

\courssub{Critères avancés dans NB.SI et SOMME.SI}

La puissance de ces fonctions réside dans la souplesse du paramètre \texttt{critère}.

\begin{propriete}[Syntaxes de critères autorisées]
\begin{center}
\begin{tabularx}{\linewidth}{|l|X|l|}\hline
\textbf{Type} & \textbf{Syntaxe dans la formule} & \textbf{Exemple} \\\hline
Valeur exacte (nombre) & \texttt{10} ou \texttt{B1} & \texttt{=NB.SI(C:C; 10)} \\\hline
Valeur exacte (texte) & \texttt{"Texte"} & \texttt{=NB.SI(D:D; "Payé")} \\\hline
Comparaison & \texttt{">100"}, \texttt{"<=B1"}, \texttt{"<>0"} & \texttt{=SOMME.SI(C:C; ">100")} \\\hline
Joker (texte) & \texttt{"*mot*"}, \texttt{"??"}, \texttt{"~*"} & \texttt{=NB.SI(A:A; "*Cam*")} \\\hline
Cellule vide / non vide & \texttt{""}, \texttt{"<>""} & \texttt{=NB.SI(B:B; "<>""")} \\\hline
Critère dynamique (concaténation) & \texttt{">=" \& B1} & \texttt{=NB.SI(Notes; ">=" \& Seuil)} \\\hline
\end{tabularx}
\end{center}
\begin{itemize}
    \item \texttt{*} remplace n'importe quelle suite de caractères (0 ou plus).
    \item \texttt{?} remplace un seul caractère.
    \item \texttt{~} annule le joker suivant (recherche littérale \texttt{*} ou \texttt{?}).
    \item La concaténation \texttt{"operator" \& Cellule} est \textbf{indispensable} pour comparer à une valeur variable (seuil en B1).
\end{itemize}
\end{propriete}

\begin{experience}[TP Guidés : Tableau de bord « Notes 3ème »]
\textbf{Objectif :} Construire un tableau de synthèse statistique automatisé.
\textbf{Données :} Onglet \texttt{Notes} : Col A = Noms (A2:A31), Col B = Devoir 1 (/20), Col C = Devoir 2 (/20), Col D = Examen (/20).
\textbf{Manipulation :}
\begin{enumerate}
    \item Créez un onglet \texttt{Synthese}. En \texttt{A1} : \texttt{"Statistique"}, \texttt{B1} : \texttt{"Devoir 1"}, \texttt{C1} : \texttt{"Devoir 2"}, \texttt{D1} : \texttt{"Examen"}.
    \item \texttt{A2} : \texttt{"Effectif"} $\rightarrow$ \texttt{B2} : \texttt{=NB(Notes!B:B)} (étendez à C2, D2).
    \item \texttt{A3} : \texttt{"Moyenne"} $\rightarrow$ \texttt{B3} : \texttt{=MOYENNE(Notes!B:B)} (format 2 déc.). Étendez.
    \item \texttt{A4} : \texttt{"Maximum"} $\rightarrow$ \texttt{B4} : \texttt{=MAX(Notes!B:B)}. Étendez.
    \item \texttt{A5} : \texttt{"Minimum"} $\rightarrow$ \texttt{B5} : \texttt{=MIN(Notes!B:B)}. Étendez.
    \item \texttt{A6} : \texttt{"Admis ($\ge$10)"} $\rightarrow$ \texttt{B6} : \texttt{=NB.SI(Notes!B:B; ">=10")}. Étendez.
    \item \texttt{A7} : \texttt{"Total points admis"} $\rightarrow$ \texttt{B7} : \texttt{=SOMME.SI(Notes!B:B; ">=10")}. Étendez.
    \item \texttt{A8} : \texttt{"Moyenne admis"} $\rightarrow$ \texttt{B8} : \texttt{=SI(B6=0; 0; B7/B6)} (protège division par 0).
    \item \textbf{Test :} Modifiez une note dans \texttt{Notes}. Vérifiez mise à jour instantanée sur \texttt{Synthese}.
\end{enumerate}
\textbf{Observation :} Les références \texttt{Notes!B:B} (colonne entière) évitent de modifier les plages si on ajoute des élèves.
\textbf{Interprétation :} Ce tableau de bord permet au professeur principal d'avoir une vision immédiate sans recalculer manuellement.
\end{experience}

\begin{savaistu}[Le tableur, ancêtre du « Low Code »]
Le premier tableur grand public, \textbf{VisiCalc} (1979, Apple II), a révolutionné la gestion d'entreprise en remplaçant les feuilles de calcul papier (14 colonnes, gomme et crayon). Dan Bricklin, son inventeur, l'a conçu en regardant un professeur refaire un tableau noir à cause d'un chiffre changé. Aujourd'hui, les fonctions \texttt{NB.SI}, \texttt{SOMME.SI} et leurs versions multiples (\texttt{NB.SI.ENS}, \texttt{SOMME.SI.ENS} — niveau supérieur) sont les « moteurs » de milliers de tableaux de bord RH, commerciaux et scolaires au Cameroun, du cybercafé de quartier à la direction du budget du MINFI. Maîtriser ces fonctions, c'est acquérir une compétence immédiatement monétisable sur le marché de l'emploi local.
\end{savaistu}

\begin{exercice}[Entraînement : Gestion de stock « Epicerie du Coin »]
Feuille \texttt{Stock} : Col A = Produit, Col B = Catégorie, Col C = Prix unitaire, Col D = Quantité, Col E = Valeur stock (C*D).
\begin{enumerate}
    \item Écrire la formule en \texttt{E2} (à recopier vers le bas) pour calculer la valeur stock.
    \item En \texttt{H1} : \texttt{"Seuil alerte"}, \texttt{H2} : \texttt{5}. Formule en \texttt{I2} : Nombre de produits dont Quantité $\le$ Seuil (référence absolue sur H2).
    \item Formule en \texttt{I3} : Valeur totale du stock pour la catégorie « Boissons » (col B).
    \item Formule en \texttt{I4} : Prix moyen unitaire des produits « Alimentaire » (astuce : SOMME.SI / NB.SI).
    \item On veut arrondir la valeur stock (col E) à la centaine supérieure (ex. 12345 $\rightarrow$ 12400). Formule en \texttt{F2}.
\end{enumerate}
\end{exercice}

\begin{corrige}[Exercice : Gestion de stock « Epicerie du Coin »]
\begin{enumerate}
    \item \texttt{E2} : \texttt{=C2*D2} (recopier vers le bas).
    \item \texttt{I2} : \texttt{=NB.SI(D:D; "<=" \& \$H\$2)}. \textbf{Note :} \texttt{\$H\$2} absolu pour verrouiller le seuil lors de la recopie.
    \item \texttt{I3} : \texttt{=SOMME.SI(B:B; "Boissons"; E:E)}. \texttt{somme\_plage} = E:E (valeur stock).
    \item \texttt{I4} : \texttt{=SOMME.SI(B:B; "Alimentaire"; C:C) / NB.SI(B:B; "Alimentaire")}. Ou \texttt{=MOYENNE.SI(B:B; "Alimentaire"; C:C)} si version récente.
    \item \texttt{F2} : \texttt{=ARRONDI.SUP(E2; -2)}. \texttt{-2} = arrondi à la centaine (100 = $10^2$). \texttt{ARRONDI.SUP} garantit le sur-stockage (prudence).
\end{enumerate}
\end{corrige}

\coursec{Fonctions logiques : SI, ET, OU, NON et imbrication}

Dans la section précédente, nous avons vu comment les fonctions mathématiques et statistiques (\texttt{SOMME}, \texttt{MOYENNE}, \texttt{NB}, \texttt{MIN}, \texttt{MAX}) permettent de \emph{résumer} des données. Mais un tableur devient un véritable outil d'aide à la décision lorsqu'il sait \emph{réagir} différemment selon les valeurs qu'il contient : « Si la moyenne est $\ge 10$, écrire "Admis", sinon écrire "Ajourné" » ; « Si le stock est critique, commander, sinon ne rien faire ». C'est le rôle des \textbf{fonctions logiques}. Elles transforment un classeur passif en un tableau de bord actif, capable d'automatiser la gestion d'une boutique, l'édition de bulletins ou le suivi de stock d'un cybercafé à Yaoundé.

\courssub{La fonction SI : la décision binaire}

La fonction \cle{SI} est la pierre angulaire de la logique dans un tableur. Elle évalue une condition (un \textbf{test logique}) et renvoie une valeur si cette condition est \textbf{VRAI}, une autre si elle est \textbf{FAUX}.

\begin{definition}[Test logique]
Un \textbf{test logique} est une expression qui ne peut prendre que deux valeurs : \cle{VRAI} (équivalent à 1) ou \cle{FAUX} (équivalent à 0). Il compare deux grandeurs (nombres, textes, dates, références de cellules) à l'aide d'un \textbf{opérateur de comparaison}.
\end{definition}

\begin{center}
\begin{tabularx}{\linewidth}{|l|X|}\hline
\rowcolor{popblueL}
\textbf{Opérateur} & \textbf{Signification} \\ \hline
= & Égal à \\ \hline
<> & Différent de \\ \hline
> & Supérieur strictement à \\ \hline
< & Inférieur strictement à \\ \hline
>= & Supérieur ou égal à \\ \hline
<= & Inférieur ou égal à \\ \hline
\end{tabularx}
\end{center}

\begin{propriete}[Syntaxe de la fonction SI]
\texttt{=SI(test\_logique; valeur\_si\_vrai; valeur\_si\_faux)}
\begin{itemize}
    \item \texttt{test\_logique} : toute expression renvoyant VRAI ou FAUX (ex: \texttt{B3>=10}, \texttt{A1="Cameroon"}, \texttt{SOMME(C1:C5)>100}).
    \item \texttt{valeur\_si\_vrai} : résultat affiché si le test est VRAI. Peut être un nombre, un texte (obligatoirement entre \texttt{" "}), une référence de cellule, une formule, ou une chaîne vide \texttt{""} (cellule visuellement vide).
    \item \texttt{valeur\_si\_faux} : résultat affiché si le test est FAUX. Même types possibles. Si omis (ex: \texttt{=SI(A1>10; "Ok"; )}), la fonction renvoie \texttt{FAUX} (ou 0 selon les versions).
\end{itemize}
\end{propriete}

\begin{exemplebox}[Décision simple : Admis ou Ajourné]
Dans une cellule (ex: \texttt{D2}), on veut afficher la décision pour l'élève dont la moyenne est en \texttt{C2}.
\begin{lstlisting}
=SI(C2>=10; "Admis"; "Ajourné")
\end{lstlisting}
\begin{itemize}
    \item Si \texttt{C2} vaut 12,5 $\to$ test VRAI $\to$ affiche \texttt{Admis}.
    \item Si \texttt{C2} vaut 8 $\to$ test FAUX $\to$ affiche \texttt{Ajourné}.
    \item Si \texttt{C2} est vide $\to$ test FAUX (vide < 10) $\to$ affiche \texttt{Ajourné}.
\end{itemize}
\end{exemplebox}

\begin{attention}[Pièges classiques avec SI]
\begin{enumerate}
    \item \textbf{Oublier les guillemets pour le texte} : \texttt{=SI(A1>10; Admis; Recale)} $\to$ Erreur \texttt{\#NOM?} (le tableur croit que \texttt{Admis} est un nom de plage ou de fonction). \textbf{Toujours} écrire \texttt{"Admis"}.
    \item \textbf{Séparateur d'arguments} : Au Cameroun (région francophone), le séparateur standard est le \textbf{point-virgule (;)}. Le point-virgule ne doit \textbf{JAMAIS} précéder la parenthèse fermante finale. \texttt{=SI(A1>10; "Ok"; "Non";)} est une erreur de syntaxe.
    \item \textbf{Comparer du texte} : La comparaison est \emph{généralement} insensible à la casse (\texttt{"yaounde"} = \texttt{"YAOUNDÉ"} renvoie VRAI dans Excel/LibreOffice Calc), mais sensible aux espaces (\texttt{" Yaoundé"} $\neq$ \texttt{"Yaoundé"}). Utilisez \texttt{NETTOYER()} ou \texttt{ESPACES()} si nécessaire.
\end{enumerate}
\end{attention}

\courssub{Combiner plusieurs conditions : ET, OU, NON}

Souvent, la décision dépend de \emph{plusieurs} critères simultanés. « L'élève est admis \textbf{ET} s'il a moins de 5 absences » ; « Le client a une remise \textbf{SI} il est fidèle \textbf{OU} s'il achète en gros volume ». Les fonctions \cle{ET}, \cle{OU} et \cle{NON} renvoient VRAI ou FAUX et s'utilisent \emph{à l'intérieur} du premier argument de \texttt{SI}.

\begin{propriete}[Fonctions ET, OU, NON]
\begin{itemize}
    \item \texttt{ET(cond1; cond2; ...)} : Renvoie \textbf{VRAI} si \textbf{TOUTES} les conditions sont vraies. Une seule fausse $\to$ FAUX.
    \item \texttt{OU(cond1; cond2; ...)} : Renvoie \textbf{VRAI} si \textbf{AU MOINS UNE} condition est vraie. Toutes fausses $\to$ FAUX.
    \item \texttt{NON(cond)} : Renvoie l'\textbf{inverse} de la condition. VRAI $\to$ FAUX, FAUX $\to$ VRAI.
\end{itemize}
\end{propriete}

\begin{exemplebox}[Admission conditionnée : Moyenne ET Assiduité]
Un élève est \texttt{Admis} \textbf{si} (\texttt{Moyenne >= 10} \textbf{ET} \texttt{Absences < 5}).
\begin{lstlisting}
=SI(ET(C2>=10; D2<5); "Admis"; "Ajourné")
\end{lstlisting}
\begin{itemize}
    \item Moyenne 12, Absences 3 $\to$ ET(VRAI; VRAI) = VRAI $\to$ \texttt{Admis}.
    \item Moyenne 12, Absences 6 $\to$ ET(VRAI; FAUX) = FAUX $\to$ \texttt{Ajourné}.
    \item Moyenne 8, Absences 2 $\to$ ET(FAUX; VRAI) = FAUX $\to$ \texttt{Ajourné}.
\end{itemize}
\end{exemplebox}

\begin{propriete}[Équivalences logiques (Lois de De Morgan)]
Utile pour simplifier ou réécrire des conditions complexes sans changer leur sens :
\[
\text{NON(ET(A; B))} \equiv \text{OU(NON(A); NON(B))}
\]
\[
\text{NON(OU(A; B))} \equiv \text{ET(NON(A); NON(B))}
\]
\emph{Exemple :} « Pas (Majeur ET Étudiant) » $\equiv$ « Mineur OU Pas Étudiant ».
\end{propriete}

\courssub{SI imbriqués : choisir parmi plusieurs cas}

Lorsque les cas possibles sont plus de deux (ex: mentions au BEPC : TB, B, AB, Passable, Insuffisant), on \textbf{imbrique} des fonctions \texttt{SI} : le 3ème argument (\texttt{valeur\_si\_faux}) devient lui-même une nouvelle fonction \texttt{SI}.

\begin{methode}[Construire un SI imbriqué (choix séquentiels)]
\begin{enumerate}
    \item \textbf{Ordonner les conditions} de la plus restrictive (la plus difficile à obtenir) à la moins restrictive. Ex: $\ge 16$, puis $\ge 14$, puis $\ge 12$, puis $\ge 10$, sinon.
    \item \textbf{Écrire le premier SI} avec la condition la plus haute. Pour \texttt{valeur\_si\_vrai}, mettre le résultat correspondant (ex: \texttt{"Très Bien"}).
    \item \textbf{À la place de \texttt{valeur\_si\_faux}}, ouvrir un \textbf{nouveau SI} pour la condition suivante.
    \item \textbf{Répéter} jusqu'au dernier cas. Le tout dernier \texttt{valeur\_si\_faux} (le "sinon final") ne contient plus de \texttt{SI}, juste la valeur par défaut (ex: \texttt{"Insuffisant"}).
    \item \textbf{Fermer toutes les parenthèses} : autant de \texttt{)} que de \texttt{SI} ouverts.
\end{enumerate}
\end{methode}

\begin{exemplebox}[Attribution des mentions au BEPC (SI imbriqués 4 niveaux)]
Note en \texttt{B2}. Formule en \texttt{C2} :
\begin{lstlisting}
=SI(B2>=16; "Très Bien";
    SI(B2>=14; "Bien";
        SI(B2>=12; "Assez Bien";
            SI(B2>=10; "Passable"; "Insuffisant")
        )
    )
)
\end{lstlisting}
\emph{Exécution pour B2 = 13 :}
1. \texttt{13>=16} ? FAUX $\to$ va au 2ème SI.
2. \texttt{13>=14} ? FAUX $\to$ va au 3ème SI.
3. \texttt{13>=12} ? VRAI $\to$ renvoie \texttt{"Assez Bien"}. Stop.
\end{exemplebox}

\begin{popfigure}
\begin{tikzpicture}[
    decision/.style={diamond, draw=popdark, thick, fill=poporangeL, text width=2.8cm, align=center, inner sep=2pt},
    result/.style={rectangle, rounded corners, draw=popgreen, thick, fill=popgreenL, text width=2.2cm, align=center, inner sep=3pt},
    start/.style={rectangle, draw=popblue, thick, fill=popblueL, text width=2.5cm, align=center, inner sep=3pt},
    arrow/.style={-Stealth, thick, popdark},
    level distance=1.8cm,
    sibling distance=3.5cm
]
\node[start] (start) {Note en B2};
\node[decision, below of=start] (d1) {B2 >= 16 ?};
\node[result, below left=1.5cm and 3cm of d1] (r1) {Très Bien};
\node[decision, below right=1.5cm and 3cm of d1] (d2) {B2 >= 14 ?};
\node[result, below left=1.5cm and 2.5cm of d2] (r2) {Bien};
\node[decision, below right=1.5cm and 2.5cm of d2] (d3) {B2 >= 12 ?};
\node[result, below left=1.5cm and 2cm of d3] (r3) {Assez Bien};
\node[decision, below right=1.5cm and 2cm of d3] (d4) {B2 >= 10 ?};
\node[result, below left=1.5cm and 1.5cm of d4] (r4) {Passable};
\node[result, below right=1.5cm and 1.5cm of d4] (r5) {Insuffisant};

\draw[arrow] (start) -- (d1);
\draw[arrow] (d1) -- node[left, pos=0.6, font=\small] {VRAI} (r1);
\draw[arrow] (d1) -- node[right, pos=0.6, font=\small] {FAUX} (d2);
\draw[arrow] (d2) -- node[left, pos=0.6, font=\small] {VRAI} (r2);
\draw[arrow] (d2) -- node[right, pos=0.6, font=\small] {FAUX} (d3);
\draw[arrow] (d3) -- node[left, pos=0.6, font=\small] {VRAI} (r3);
\draw[arrow] (d3) -- node[right, pos=0.6, font=\small] {FAUX} (d4);
\draw[arrow] (d4) -- node[left, pos=0.6, font=\small] {VRAI} (r4);
\draw[arrow] (d4) -- node[right, pos=0.6, font=\small] {FAUX} (r5);
\end{tikzpicture}
\legende{Arbre de décision : imbrication de 4 SI pour l'attribution des mentions au BEPC. L'évaluation suit le chemin des flèches "FAUX" jusqu'à trouver un VRAI.}
\end{popfigure}

\begin{attention}[Limites de l'imbrication]
\begin{itemize}
    \item \textbf{Lisibilité} : Au-delà de 3-4 niveaux, la formule devient illisible et difficile à déboguer (erreurs de parenthèses fréquentes).
    \item \textbf{Limite technique} : Excel/LibreOffice limitent à \textbf{64 niveaux} d'imbrication (suffisant pour le BEPC, mais pas pour des échelles fines).
    \item \textbf{Ordre crucial} : Si on inverse l'ordre (tester $\ge 10$ en premier), une note de 17 serait catégorisée "Passable" car le premier test serait VRAI et les suivants ignorés. \textbf{Toujours tester du plus grand au plus petit} (ou du plus petit au plus grand avec des \texttt{ET} pour borner : \texttt{ET(Note>=14; Note<16)}).
\end{itemize}
\end{attention}

\courssub{Alternative moderne : SI.MULTIPLE (équivalent IFS)}

Depuis Excel 2019 et LibreOffice 5.2, la fonction \cle{SI.MULTIPLE} (anglais \texttt{IFS}) permet de gérer une liste de conditions \textbf{sans imbrication}. Elle évalue les conditions dans l'ordre et s'arrête à la première VRAIE.

\begin{propriete}[Syntaxe de SI.MULTIPLE]
\texttt{=SI.MULTIPLE(condition1; résultat1; [condition2; résultat2; ...]; [défaut])}
\begin{itemize}
    \item Les paires \texttt{condition; résultat} s'enchaînent.
    \item Le dernier argument optionnel \texttt{défaut} (sans condition) est la valeur si \textbf{aucune} condition n'est remplie (équivalent du dernier \texttt{valeur\_si\_faux}).
    \item \textbf{Avantage} : Une seule parenthèse ouvrante/fermante, lecture linéaire haut-bas.
\end{itemize}
\end{propriete}

\begin{center}
\begin{tabularx}{\linewidth}{|l|X|X|}\hline
\rowcolor{popblueL}
\textbf{Aspect} & \textbf{SI imbriqué classique} & \textbf{SI.MULTIPLE (Moderne)} \\ \hline
Syntaxe & \texttt{=SI(c1; r1; SI(c2; r2; SI(c3; r3; r4)))} & \texttt{=SI.MULTIPLE(c1; r1; c2; r2; c3; r3; r4)} \\ \hline
Parenthèses & Imbriquées (4 ouvrantes, 4 fermantes) & Plates (1 ouvrante, 1 fermante) \\ \hline
Lecture & De l'intérieur vers l'extérieur & Linéaire : condition 1, puis 2, puis 3... \\ \hline
Maintenance & Difficile (ajouter un cas = tout refaire) & Facile (insérer une paire au milieu) \\ \hline
Disponibilité & Toutes versions (Excel 97+, Calc) & Excel 2019+, LibreOffice 5.2+, Google Sheets \\ \hline
\end{tabularx}
\end{center}

\begin{exemplebox}[Calcul de remise par tranche de quantité]
\textbf{Règle de gestion d'une boutique à Douala :}
\begin{itemize}
    \item Quantité $\ge$ 100 $\to$ Remise 10\%
    \item Quantité $\ge$ 50 $\to$ Remise 5\%
    \item Sinon $\to$ Remise 0\%
\end{itemize}
Quantité en \texttt{B2}.

\textbf{Version SI imbriqué (compatible toutes versions) :}
\begin{lstlisting}
=SI(B2>=100; 10\%; SI(B2>=50; 5\%; 0\%))
\end{lstlisting}

\textbf{Version SI.MULTIPLE (recommandée si dispo) :}
\begin{lstlisting}
=SI.MULTIPLE(B2>=100; 10\%; B2>=50; 5\%; 0\%)
\end{lstlisting}

\textbf{Test :} Si \texttt{B2 = 75}.
\begin{itemize}
    \item SI imbriqué : 75>=100 (Faux) $\to$ 75>=50 (Vrai) $\to$ \textbf{5\%}.
    \item SI.MULTIPLE : teste 100 (Faux), teste 50 (Vrai) $\to$ \textbf{5\%}.
\end{itemize}
\end{exemplebox}

\courssub{Gérer les erreurs dans la logique : SIERREUR}

Une formule logique peut renvoyer une erreur technique (\texttt{\#DIV/0!}, \texttt{\#N/A}, \texttt{\#VALEUR!}, \texttt{\#REF!}) qui « casse » l'affichage et les calculs en aval. \cle{SIERREUR} permet d'intercepter ces erreurs et d'afficher une valeur propre.

\begin{propriete}[Syntaxe de SIERREUR]
\texttt{=SIERREUR(formule; valeur\_si\_erreur)}
\begin{itemize}
    \item Si \texttt{formule} s'évalue normalement $\to$ renvoie son résultat.
    \item Si \texttt{formule} renvoie \textbf{n'importe quelle erreur} $\to$ renvoie \texttt{valeur\_si\_erreur}.
\end{itemize}
\emph{Variante :} \texttt{SI.NON.DISP(formule; valeur)} ne capture \textbf{que} \texttt{\#N/A} (utile pour \texttt{RECHERCHEV}).
\end{propriete}

\begin{exemplebox}[Stock : éviter \#DIV/0! lors du calcul du prix unitaire]
Colonne \texttt{B} = Stock, Colonne \texttt{C} = Valeur totale. Prix unitaire = \texttt{C2/B2}.
Si \texttt{B2 = 0} (stock épuisé) $\to$ \texttt{\#DIV/0!}.
\begin{lstlisting}
=SIERREUR(C2/B2; "Stock épuisé")
\end{lstlisting}
Résultat : Affiche le prix calculé, ou le texte \texttt{"Stock épuisé"} si B2=0 ou vide.
\end{exemplebox}

\courssub{Applications concrètes et Exercices Résolus}

\begin{exoresolu}[Gestion de stock : Alerte commande automatique]
\textbf{Énoncé :} Un gérant de cybercafé suit ses cartouches d'encre. Colonnes : \texttt{A}=Produit, \texttt{B}=Stock Actuel, \texttt{C}=Seuil Alerte, \texttt{D}=Commande Auto (Oui/Non). Écrire la formule en \texttt{D2} pour afficher \texttt{"COMMANDER"} si Stock < Seuil, \texttt{"OK"} sinon. Protéger contre lignes vides.
\textbf{Solution détaillée :}
\begin{enumerate}
    \item \textbf{Analyse} : Test logique = \texttt{B2 < C2}. Résultat VRAI = "COMMANDER", FAUX = "OK". Gérer le cas où B2 ou C2 est vide (ne rien afficher).
    \item \textbf{Formule robuste} (avec ET pour vérifier que les cellules ne sont pas vides) :
\begin{lstlisting}
=SI(ET(B2<>""; C2<>""); SI(B2<C2; "COMMANDER"; "OK"); "")
\end{lstlisting}
    \item \textbf{Explication} :
    \begin{itemize}
        \item \texttt{ET(B2<>""; C2<>"")} : VRAI seulement si les deux cellules ont du contenu.
        \item Si VRAI : on évalue le \texttt{SI} principal (Stock < Seuil).
        \item Si FAUX (ligne vide) : on renvoie \texttt{""} (cellule vide visuelle), pas "OK" ni "COMMANDER".
    \end{itemize}
\end{enumerate}
\end{exoresolu}

\begin{experience}[TP Guidé : Bulletin de notes avec mentions automatiques]
\textbf{Objectif} : Créer un mini-bulletin pour 5 élèves calculant la moyenne, la décision (Admis/Ajourné) et la Mention (TB/B/AB/Passable/Insuffisant) avec \texttt{SI.MULTIPLE} (ou SI imbriqué).
\textbf{Matériel} : LibreOffice Calc ou Excel.
\textbf{Manipulation} :
\begin{enumerate}
    \item Créer les en-têtes en ligne 1 : \texttt{A1:Nom}, \texttt{B1:Note1}, \texttt{C1:Note2}, \texttt{D1:Note3}, \texttt{E1:Moyenne}, \texttt{F1:Décision}, \texttt{G1:Mention}.
    \item Saisir 5 noms et 3 notes chacun (valeurs entre 0 et 20).
    \item En \texttt{E2} : \texttt{=MOYENNE(B2:D2)}. Étirer vers le bas.
    \item En \texttt{F2} (Décision) : \texttt{=SI(E2>=10; "Admis"; "Ajourné")}. Étirer.
    \item En \texttt{G2} (Mention) :
    \begin{itemize}
        \item \textbf{Si SI.MULTIPLE dispo :} \texttt{=SI.MULTIPLE(E2>=16; "TB"; E2>=14; "B"; E2>=12; "AB"; E2>=10; "Passable"; "Insuffisant")}
        \item \textbf{Sinon (SI imbriqué) :} \texttt{=SI(E2>=16; "TB"; SI(E2>=14; "B"; SI(E2>=12; "AB"; SI(E2>=10; "Passable"; "Insuffisant"))))}
    \end{itemize}
    Étirer.
    \item \textbf{Observation} : Modifier une note (ex: passer de 9 à 11) $\to$ Moyenne, Décision et Mention se mettent à jour \textbf{instantanément}.
    \item \textbf{Interprétation} : Le tableur agit comme un \emph{système expert} : il applique le règlement d'examen (décret MINESEC) sans intervention humaine sur chaque ligne. C'est l'automatisation de la gestion administrative scolaire.
\end{enumerate}
\end{experience}

\begin{savaistu}[Le saviez-vous ? L'évaluation paresseuse (Short-circuit)]
La fonction \texttt{SI.MULTIPLE} (et \texttt{SI} classique dans les tableurs modernes) utilise l'\textbf{évaluation paresseuse} : elle teste les conditions \textbf{une par une dans l'ordre} et \textbf{s'arrête immédiatement} dès qu'une condition est VRAIE. Les conditions suivantes ne sont \textbf{jamais évaluées}.
\emph{Conséquence pratique :} On peut mettre en premier une condition qui provoquerait une erreur si elle était évaluée seule. Ex: \texttt{=SI.MULTIPLE(B2=0; "Zéro"; 10/B2>5; "Grand"; "Petit")}. Si B2=0, la 1ère condition est VRAIE $\to$ renvoie "Zéro". La division \texttt{10/B2} (qui ferait \#DIV/0!) n'est \textbf{jamais exécutée}. C'est une astuce de pro pour sécuriser les formules sans \texttt{SIERREUR}.
\end{savaistu}

\begin{exercice}[Entraînement : Calcul de prime de fin d'année]
\textbf{Contexte :} Une PME de Bafoussam verse une prime selon l'ancienneté (colonne \texttt{B}, en années) et la performance (colonne \texttt{C}, note sur 20).
Règles :
\begin{enumerate}
    \item Si Ancienneté $< 2$ ans $\to$ Prime = 0 FCFA (période d'essai).
    \item Sinon, si Performance $\ge 16$ $\to$ Prime = 100 000 FCFA.
    \item Sinon, si Performance $\ge 12$ $\to$ Prime = 50 000 FCFA.
    \item Sinon $\to$ Prime = 10 000 FCFA (prime de fidélité).
\end{enumerate}
\textbf{Questions :}
\begin{enumerate}
    \item Écrire la formule en \texttt{D2} avec des \texttt{SI} imbriqués (version compatible toutes versions).
    \item Réécrire avec \texttt{SI.MULTIPLE}.
    \item Tester avec : (Ancienneté=1; Perf=18), (Ancienneté=5; Perf=14), (Ancienneté=3; Perf=10).
\end{enumerate}
\end{exercice}

\begin{corrige}[Exercice : Calcul de prime]
\begin{enumerate}
    \item \textbf{SI imbriqué} (Attention à l'ordre : ancienneté d'abord, puis performance décroissante) :
\begin{lstlisting}
=SI(B2<2; 0; SI(C2>=16; 100000; SI(C2>=12; 50000; 10000)))
\end{lstlisting}
    \item \textbf{SI.MULTIPLE} :
\begin{lstlisting}
=SI.MULTIPLE(B2<2; 0; C2>=16; 100000; C2>=12; 50000; 10000)
\end{lstlisting}
    \item \textbf{Tests} :
    \begin{itemize}
        \item (1; 18) : B2<2 VRAI $\to$ \textbf{0 FCFA}.
        \item (5; 14) : B2<2 FAUX, C2>=16 FAUX, C2>=12 VRAI $\to$ \textbf{50 000 FCFA}.
        \item (3; 10) : B2<2 FAUX, C2>=16 FAUX, C2>=12 FAUX $\to$ Défaut \textbf{10 000 FCFA}.
    \end{itemize}
\end{enumerate}
\end{corrige}

\courssub{Résumé de la section 4}
\begin{aretenir}[Ce qu'il faut retenir des fonctions logiques]
\begin{itemize}
    \item \textbf{SI(test; Vrai; Faux)} : Décision binaire de base. Test = comparaison avec \texttt{=, <>, >, <, >=, <=}.
    \item \textbf{ET / OU / NON} : Combinent des tests. \texttt{ET} = TOUTES vraies. \texttt{OU} = AU MOINS UNE vraie. \texttt{NON} = inverse.
    \item \textbf{SI imbriqué} : Pour 3+ cas. Ordre : condition la plus restrictive $\to$ la moins restrictive. Fermer toutes les parenthèses.
    \item \textbf{SI.MULTIPLE(cond1; res1; cond2; res2; ...; défaut)} : Moderne, lisible, plate, recommandée (dispo Excel 2019+/LibreOffice 5.2+).
    \item \textbf{SIERREUR(formule; secours)} : Masque \texttt{\#DIV/0!}, \texttt{\#N/A}, etc. Indispensable pour des rapports propres.
    \item \textbf{Application} : Bulletins (mentions), Stocks (commande auto), Facturation (remises), Paie (primes).
\end{itemize}
\end{aretenir}

\coursec{Mise en forme conditionnelle et Visualisation graphique}

Après avoir maîtrisé les fonctions logiques (\texttt{SI}, \texttt{ET}, \texttt{OU}, \texttt{NON}) qui permettent de prendre des décisions \emph{à l'intérieur} d'une cellule, nous abordons maintenant une étape cruciale : \textbf{la communication visuelle des résultats}. Un tableau brut de centaines de lignes est illisible pour un décideur. La \textbf{Mise en Forme Conditionnelle (MFC)} automatise le surlignage des anomalies, des meilleurs vendeurs ou des seuils critiques. Les \textbf{graphiques} transforment les nombres en images parlantes. C'est le passage du \emph{calcul} à l'\emph{information}.

\courssub{Principe et Règles de base de la Mise en Forme Conditionnelle}

\begin{definition}[Mise en Forme Conditionnelle (MFC)]
La MFC est une fonctionnalité qui modifie \textbf{automatiquement} l'apparence d'une ou plusieurs cellules (police, couleur de fond, bordures, icônes, barres) \textbf{en fonction de leur valeur ou du résultat d'une formule}. La mise à jour est dynamique : si la valeur change, le format change instantanément.
\end{definition}

Contrairement à une mise en forme manuelle (sélection $\to$ couleur rouge), la MFC pose une \textbf{règle logique} : \og Si condition vraie $\to$ Appliquer ce format \fg.

\begin{propriete}[Règles types prédéfinies (Menu \texttt{Accueil} $\to$ \texttt{Mise en forme conditionnelle} $\to$ \texttt{Règles de mise en surbrillance des cellules})]
Ces règles couvrent 90\% des besoins courants sans écrire de formule :
\begin{itemize}
    \item \textbf{Supérieur à / Inférieur à / Entre / Égal à} : Seuil numérique (ex: Notes $< 10$, Ventes $> 100\,000$).
    \item \textbf{Texte contenant} : Filtrage textuel (ex: Statut = \og Impayé \fg $\to$ Fond rouge).
    \item \textbf{Date du jour / Cette semaine / Le mois dernier} : Suivi temporel dynamique.
    \item \textbf{Valeurs en double / Uniques} : Détection d'erreurs de saisie (ex: Numéros d'élèves dupliqués).
\end{itemize}
\end{propriete}

\begin{exemplebox}[Application directe : Repérer les notes insuffisantes]
\textbf{Contexte} : Feuille de notes de 3ème, colonne \texttt{D} (Moyennes), lignes 3 à 50.
\textbf{Objectif} : Colorer en rouge clair (fond) / rouge foncé (police) toute moyenne $< 10$.
\textbf{Démarche} :
\begin{enumerate}
    \item Sélectionner la plage \texttt{D3:D50} (la sélection \emph{avant} la règle est cruciale, voir \texttt{attention}).
    \item \texttt{Accueil} $\to$ \texttt{Mise en forme conditionnelle} $\to$ \texttt{Règles de mise en surbrillance des cellules} $\to$ \texttt{Inférieur à...}
    \item Saisir \texttt{10} dans la zone de gauche.
    \item Choisir \texttt{Fond rouge clair, Texte rouge foncé} (ou \texttt{Format personnalisé...} pour affiner).
    \item \texttt{OK}.
\end{enumerate}
\emph{Résultat} : Tout élève sous la moyenne apparaît immédiatement. Si une note est corrigée et passe à 11, le rouge disparaît tout seul.
\end{exemplebox}

\begin{attention}[Erreur classique : Ordre Sélection / Création de règle]
\textbf{Piège} : Créer la règle sur une seule cellule (ex: \texttt{D3}), puis étendre la plage via \texttt{Gérer les règles} $\to$ \texttt{Modifie la plage}.
\textbf{Conséquence} : Les références relatives dans les formules personnalisées ne s'adaptent pas comme on l'espère (ex: \texttt{=D3<10} devient \texttt{=D4<10} sur la ligne 4, mais si la formule est complexe, ça casse).
\textbf{Bonne pratique} : \textbf{Toujours sélectionner la plage cible complète \emph{avant} d'ouvrir la boîte de dialogue "Nouvelle règle"}. La \og Cellule active \fg (souvent la première de la sélection) sert de référence pour écrire la formule.
\end{attention}

\courssub{Règle personnalisée avec formule : Colorer une ligne entière}

C'est la technique professionnelle par excellence. Elle permet de formater la ligne \textbf{A3:F3} (Nom, Produit, Qté, Prix, Total, Région) si la colonne \textbf{E (Total)} dépasse un seuil.

\begin{methode}[Colorer une ligne entière selon une valeur d'une colonne]
\textbf{Objectif} : Mettre en surbrillance toute la ligne si le Chiffre d'Affaires (colonne E) $> 100\,000$ FCFA.
\textbf{Données} : Tableau structuré en \texttt{A3:F50} (En-têtes en ligne 2, données à partir de 3).
\begin{enumerate}
    \item \textbf{Sélectionner toute la plage de données} : \texttt{A3:F50} (pas les en-têtes).
    \item \texttt{Accueil} $\to$ \texttt{Mise en forme conditionnelle} $\to$ \texttt{Nouvelle règle...}
    \item Choisir \texttt{Utiliser une formule pour déterminer pour quelles cellules le format sera appliqué}.
    \item Saisir la formule : \texttt{=\$E3>100000}
    \item Cliquer \texttt{Format...} $\to$ Onglet \texttt{Remplissage} $\to$ Choisir une couleur (ex: Vert clair) $\to$ \texttt{OK} $\to$ \texttt{OK}.
\end{enumerate}
\end{methode}

\begin{propriete}[Le secret du \texttt{\$} (Verrouillage de colonne) dans la MFC]
Dans la formule \texttt{=\$E3>100000} :
\begin{itemize}
    \item \texttt{\$E} (colonne verrouillée) : La condition teste \textbf{toujours} la colonne E (Chiffre d'affaires), quelle que soit la cellule formatée (colonne A, B, C...). C'est le \textbf{point d'ancrage}.
    \item \texttt{3} (ligne \emph{relative}, sans \$) : La ligne s'adapte. Sur la ligne 3, Excel teste \texttt{E3}. Sur la ligne 4, il teste \texttt{E4}. Sur la ligne 50, \texttt{E50}.
\end{itemize}
\emph{Règle mnémotechnique} : \og Je verrouille la colonne où se trouve l'information de décision, je laisse la ligne bouger pour balayer le tableau. \fg
\end{propriete}

\begin{popfigure}
\begin{tikzpicture}[
    node distance=0.8cm and 1.2cm,
    box/.style={rectangle, draw=popdark, thick, rounded corners=3pt, fill=popcream, text width=4.5cm, align=center, font=\small},
    btn/.style={rectangle, draw=popblue, thick, rounded corners=2pt, fill=popblueL, font=\small\bfseries, text width=3.5cm, align=center},
    arr/.style={-Stealth, thick, popdark},
    title/.style={font=\bfseries\large, text=popblue},
    code/.style={font=\ttfamily\small, text=poppurple, fill=poppurple!10, rounded corners=2pt, inner sep=2pt}
]
% Title
\node[title] (t) at (0,4.5) {Boîte de dialogue "Nouvelle règle de mise en forme" (Excel / LibreOffice Calc)};
% Left column: Steps
\node[box, align=center] (s1) at (-5, 3){\textbf{Étape 1 : Type de règle}\\\texttt{Utiliser une formule pour déterminer...}};
\node[box, align=center] (s2) at (-5, 1.5){\textbf{Étape 2 : La Formule}\\\texttt{=\$E3>100000}\\\textit{(E = Colonne test, 3 = 1ère ligne sélectionnée)}};
\node[box, align=center] (s3) at (-5, 0){\textbf{Étape 3 : Format}\\\texttt{Format...} $\to$ \texttt{Remplissage} $\to$ \textit{Vert clair}};
% Right column: Explanation
\node[box, fill=popgreenL, text width=5cm] (expl) at (3.5, 2) {
    \textbf{Pourquoi \$E3 et pas E3 ?}\\
    \texttt{\$E} : Colonne \textbf{fixe} (ancrage).\\
    \texttt{3} : Ligne \textbf{relative} (défilement).\\
    \medskip
    \textit{Si on met \$E\$3 :} Seule la ligne 3 serait testée pour tout le tableau. \textbf{Erreur !}\\
    \textit{Si on met E3 (sans \$) :} Sur la colonne A, on teste A3 ; sur B, B3... \textbf{Erreur !}
};
% Arrow flow
\draw[arr] (s1.south) -- (s2.north);
\draw[arr] (s2.south) -- (s3.north);
\draw[arr, dashed] (s2.east) -- ++(1,0) |- (expl.west) node[midway, above, font=\tiny, text=poporange] {Logique du \$};
% Preview area simulation
\node[box, fill=popblue!10, draw=popblue, text width=6cm, align=left, font=\small] (preview) at (0, -2) {
    \textbf{Aperçu du résultat sur le tableau (A3:F5)} \\
    \begin{tabular}{|l|c|c|c|c|l|}
    \hline \rowcolor{popblueL} \textbf{Nom} & \textbf{Produit} & \textbf{Qté} & \textbf{Prix} & \textbf{Total} & \textbf{Région} \\ \hline
    \cellcolor{popgreenL} Dupont & Riz & 50 & 2500 & \cellcolor{popgreenL} 125000 & \cellcolor{popgreenL} Douala \\ \hline
    Martin & Maïs & 20 & 3000 & 60000 & Yaoundé \\ \hline
    \cellcolor{popgreenL} Nkolo & Huile & 100 & 1500 & \cellcolor{popgreenL} 150000 & \cellcolor{popgreenL} Bafoussam \\ \hline
    \end{tabular}
};
\draw[arr, dashed] (s3.south) -- ++(0,-0.5) -| (preview.north);
\end{tikzpicture}
\legende{Fenêtre "Nouvelle règle" : Utilisation d'une formule avec verrouillage de colonne (\texttt{\$E3}) pour colorer la ligne entière.}
\end{popfigure}

\begin{exoresolu}[MFC avancée : Moyenne de la classe]
\textbf{Énoncé} : Dans un tableau \texttt{A3:C30} (Nom, Matière, Note), colorer en jaune la ligne de tout élève dont la note est \textbf{supérieure à la moyenne de la classe} (colonne C).
\textbf{Solution détaillée} :
\begin{enumerate}
    \item Sélection \texttt{A3:C30}.
    \item Nouvelle règle $\to$ Formule : \texttt{=\$C3>MOYENNE(\$C\$3:\$C\$30)}.
    \item \textbf{Analyse des \$} :
    \begin{itemize}
        \item \texttt{\$C3} : On teste la note de la ligne courante (colonne C verrouillée, ligne relative).
        \item \texttt{\$C\$3:\$C\$30} : La plage de calcul de la moyenne est \textbf{absolue} (verrouillée lignes et colonnes). Elle ne doit pas bouger quand on descend dans le tableau.
    \end{itemize}
    \item Format : Remplissage Jaune.
\end{enumerate}
\textbf{Vérification} : Si la moyenne de la classe est 11,5, l'élève avec 12 devient jaune. Si on change une note et que la moyenne monte à 12,5, le jaune disparaît dynamiquement.
\tcblower
\textbf{Formule finale} : \texttt{=\$C3>MOYENNE(\$C\$3:\$C\$30)}
\end{exoresolu}

\courssub{Visualisation quantitative intégrée : Barres, Échelles, Icônes}

Au-delà du binaire \og Condition remplie / Non remplie \fg, ces trois outils donnent une \textbf{information graduée} directement dans la cellule, sans graphique séparé.

\begin{propriete}[Barres de données (Data Bars) — \texttt{MFC} $\to$ \texttt{Barres de données}]
Affiche un \textbf{histogramme miniature} à l'intérieur de chaque cellule. La longueur de la barre est proportionnelle à la valeur par rapport au Min/Max de la sélection.
\begin{itemize}
    \item \textbf{Usage} : Comparer rapidement des chiffres d'affaires, des stocks, des notes sur effectif.
    \item \textbf{Réglage fin} : \texttt{Autres règles...} $\to$ Type \texttt{Nombre} pour Min/Max (ex: Min=0, Max=20 pour des notes) évite qu'une note 19 fasse une barre "pleine" si le max réel est 19.
    \item \textbf{Affichage} : Option \texttt{Barre pleine} ou \texttt{Dégradé}. Masquer la valeur (afficher barre seule) pour tableau de bord propre.
\end{itemize}
\end{propriete}

\begin{propriete}[Échelles de couleurs (Color Scales) — \texttt{MFC} $\to$ \texttt{Échelles de couleurs}]
Applique un \textbf{dégradé de 2 ou 3 couleurs} sur le fond des cellules selon la valeur (Min $\to$ Milieu $\to$ Max).
\begin{itemize}
    \item \textbf{2 couleurs} (ex: Blanc $\to$ Rouge) : Intensité = distance au minimum. Bien pour "Chaleur" (température, taux de remplissage).
    \item \textbf{3 couleurs} (ex: \textbf{Rouge} $\to$ \textbf{Jaune} $\to$ \textbf{Vert}) : Le \textbf{Milieu} (souvent Médiane ou Moyenne) est Jaune. En dessous $\to$ Rouge (mauvais), Au-dessus $\to$ Vert (bon). \textbf{Standard "Feu tricolore" pour KPI}.
    \item \textbf{Réglage} : \texttt{Autres règles...} permet de fixer Min/Milieu/Max sur \texttt{Nombre}, \texttt{Pourcentage}, \texttt{Percentile}, \texttt{Formule}.
\end{itemize}
\end{propriete}

\begin{propriete}[Jeux d'icônes (Icon Sets) — \texttt{MFC} $\to$ \texttt{Jeux d'icônes}]
Affiche un symbole (flèches, feux tricolores, drapeaux, barres de signal) selon des seuils.
\begin{itemize}
    \item \textbf{Configuration par défaut} : 3 ou 4 intervalles basés sur les percentiles (ex: $\ge 67\%$ $\to$ Flèche verte haut ; $33-67\%$ $\to$ Flèche latérale ; $< 33\%$ $\to$ Flèche rouge bas).
    \item \textbf{Personnalisation pro} : \texttt{Gérer les règles} $\to$ \texttt{Modifier la règle} $\to$ Changer \texttt{Type} en \texttt{Nombre} et fixer les seuils métier (ex: $\ge 100\,000$ $\to$ Vert ; $50\,000-99\,999$ $\to$ Jaune ; $< 50\,000$ $\to$ Rouge).
    \item \textbf{Option "Afficher l'icône seulement"} : Cache le nombre, ne laisse que l'indicateur visuel. Idéal pour tableaux de direction.
\end{itemize}
\end{propriete}

\begin{exemplebox}[Tableau de bord "Ventes Mensuelles" — Combinaison des 3]
\textbf{Données} : Colonne \texttt{B} (Mois), \texttt{C} (Ventes FCFA).
\begin{center}
\begin{tabularx}{\linewidth}{|l|X|X|X|}
\hline
\textbf{Outil} & \textbf{Configuration} & \textbf{Information visuelle immédiate} & \textbf{Cas d'usage} \\ \hline
\textbf{Barres de données} & Min=0, Max=Max(Ventes), Dégradé Vert & \textbf{Amplitude} : "Mars est 2x Janvier" & Classement rapide des mois \\ \hline
\textbf{Échelle 3 couleurs} & Rouge (Min) - Jaune (Moyenne) - Vert (Max) & \textbf{Position relative} : "Juin est dans le vert (au-dessus moy)" & Détection écarts vs norme \\ \hline
\textbf{Jeu d'icônes} & Feux tricolores, Seuils : $\ge 150k$ Vert, $100k-149k$ Jaune, $<100k$ Rouge & \textbf{Statut KPI} : "Objectif atteint / Vigilance / Alerte" & Reporting direction / Seuil de rentabilité \\ \hline
\end{tabularx}
\end{center}
\end{exemplebox}

\courssub{Création et Personnalisation de Graphiques}

Le graphique est la \textbf{traduction visuelle} du tableau pour la présentation orale ou le rapport écrit. Le choix du type n'est pas esthétique : il dépend de la \textbf{nature des variables}.

\begin{propriete}[Choix du type de graphique selon les données]
\begin{center}
\begin{tabularx}{\linewidth}{|l|X|X|l|}
\hline
\textbf{Type de variable X (Abscisses)} & \textbf{Type de variable Y (Ordonnées)} & \textbf{Graphique recommandé} & \textbf{Question répondue} \\ \hline
Qualitative / Catégorielle (Mois, Produits, Régions) & Quantitative (Montants, Quantités) & \textbf{Histogramme (Colonnes groupées)} ou \textbf{Barres horizontales} & \og Qui est le plus grand ? \fg (Comparaison) \\ \hline
Temporelle ordonnée (Jan, Fév, Mar... ou Dates) & Quantitative continue & \textbf{Courbe (Lignes)} & \og Comment évolue ? \fg (Tendance, saisonnalité) \\ \hline
Qualitative (Parts d'un tout = 100\%) & Quantitative (Parts, \%) & \textbf{Camembert (Secteurs)} ou \textbf{Anneau} & \og Quelle part représente chacun ? \fg (Structure) \\ \hline
Quantitative (ex: Age) & Quantitative (ex: Salaire) & \textbf{Nuage de points (XY)} & \og Y a-t-il corrélation ? \fg (Relation) \\ \hline
\end{tabularx}
\end{center}
\end{propriete}

\begin{methode}[Créer un graphique professionnel en 5 étapes]
\begin{enumerate}
    \item \textbf{Sélection intelligente} : Sélectionner \textbf{étiquettes (en-têtes) + valeurs}. Ex: \texttt{A2:B13} (Mois + Ventes). \textit{Astuce} : \texttt{Ctrl + *} (tout le tableau contigu).
    \item \textbf{Insertion ciblée} : Onglet \texttt{Insertion} $\to$ Groupe \texttt{Graphiques} $\to$ Choisir le type (ex: \texttt{Histogramme} $\to$ \texttt{Colonnes groupées 2D}). \textit{Éviter} les versions 3D (lisibilité faussée).
    \item \textbf{Titre explicite} : Cliquer sur "Titre du graphique" $\to$ Saisir : \og \textbf{Évolution des ventes 2024 (FCFA)} \fg (Contexte + Unité + Période).
    \item \textbf{Axes nommés} : \texttt{Éléments du graphique} ($+$) $\to$ \texttt{Titres des axes} $\to$ Horizontal : \og Mois \fg ; Vertical : \og Chiffre d'affaires (FCFA) \fg.
    \item \textbf{Étiquettes de données + Style} : \texttt{Éléments} $\to$ \texttt{Étiquettes de données} $\to$ \texttt{Extérieur} (valeurs sur les barres). \texttt{Style de graphique} (pinceau) $\to$ Palette professionnelle (ex: Style 6 ou 8, couleurs unies, pas arc-en-ciel).
\end{enumerate}
\end{methode}

\begin{exemplebox}[Calcul des angles pour un Camembert — Vérification manuelle]
\textbf{Données} : Ventes 1er Trimestre : Janv = 120, Fév = 150, Mar = 90. Total = 360.
\textbf{Calcul des parts (secteurs)} :
\begin{align*}
\text{Total} &= 120 + 150 + 90 = 360 \\
\text{Part Janv} &= \frac{120}{360} \times 100 = 33{,}3\% \quad \to \quad \text{Angle} = \frac{120}{360} \times 360^\circ = 120^\circ \\
\text{Part Fév} &= \frac{150}{360} \times 100 = 41{,}7\% \quad \to \quad \text{Angle} = 150^\circ \\
\text{Part Mar} &= \frac{90}{360} \times 100 = 25{,}0\% \quad \to \quad \text{Angle} = 90^\circ \\
\text{Vérification} &= 120^\circ + 150^\circ + 90^\circ = 360^\circ \quad \checkmark
\end{align*}
\textbf{Interprétation} : Février représente près de 42\% du CA du trimestre. Le camembert est lisible car seulement 3 parts. \textbf{Au-delà de 6-7 parts, le camembert devient illisible} $\to$ Préférer un histogramme.
\end{exemplebox}

\courssub{Galerie comparative : Trois regards sur les mêmes données}

Le même jeu de données (Ventes mensuelles 2024) raconte trois histoires différentes selon le graphique choisi.

\begin{popfigure}
\begin{tikzpicture}[
    scale=0.75,
    every node/.style={font=\small},
    bar/.style={fill=popblue!70, draw=popblue},
    axis/.style={thick, popdark},
    tick/.style={thin, popdark},
    label/.style={font=\tiny, popdark},
    title/.style={font=\bfseries\small, text=popblue, yshift=5pt}
]
% Data
\def\sales{{120,150,90,180,160,140,110,130,170,190,155,135}}
\def\months{{"J","F","M","A","M","J","J","A","S","O","N","D"}}
\pgfmathsetmacro{\maxval}{200}
\pgfmathsetmacro{\barwidth}{0.35}
\pgfmathsetmacro{\gap}{0.15}
% --- LEFT: HISTOGRAMME ---
\begin{scope}[xshift=0cm]
    \node[title] at (6.5, 8.5) {Histogramme : Comparaison mensuelle};
    % Axes
    \draw[axis] (0,0) -- (13.5,0) node[right, label] {Mois};
    \draw[axis] (0,0) -- (0,9) node[above, label] {CA (kFCFA)};
    % Ticks Y
    \foreach \y/\yl in {0/0, 2.25/50, 4.5/100, 6.75/150, 9/200}
        \draw[tick] (-0.1,\y) -- (0.1,\y) node[left, label] {\yl};
    % Bars
    \foreach \i in {0,...,11} {
        \pgfmathsetmacro{\x}{1 + \i * (\barwidth + \gap) + \barwidth/2}
        \pgfmathsetmacro{\h}{\sales[\i] / \maxval * 9}
        \pgfmathsetmacro{\val}{\sales[\i]}
        \fill[bar] ({\x - \barwidth/2}, 0) rectangle ({\x + \barwidth/2}, \h);
        \node[label, rotate=90, anchor=west, yshift=2pt] at (\x, \h) {\val};
        \node[label] at (\x, -0.3) {\pgfmathparse{\months[\i]}\pgfmathresult};
    }
\end{scope}
% --- CENTER: COURBE ---
\begin{scope}[xshift=15cm]
    \node[title] at (6.5, 8.5) {Courbe : Tendance \& Saisonnalité};
    % Axes
    \draw[axis] (0,0) -- (13.5,0) node[right, label] {Mois};
    \draw[axis] (0,0) -- (0,9) node[above, label] {CA (kFCFA)};
    % Ticks Y
    \foreach \y/\yl in {0/0, 2.25/50, 4.5/100, 6.75/150, 9/200}
        \draw[tick] (-0.1,\y) -- (0.1,\y) node[left, label] {\yl};
    % Markers
    \foreach \i in {0,...,11} {
        \pgfmathsetmacro{\x}{1 + \i * (\barwidth + \gap) + \barwidth/2}
        \pgfmathsetmacro{\h}{\sales[\i] / \maxval * 9}
        \fill[poporange] (\x, \h) circle (2pt);
        \node[label] at (\x, -0.3) {\pgfmathparse{\months[\i]}\pgfmathresult};
    }
\end{scope}
% --- RIGHT: CAMENBERT ---
\begin{scope}[xshift=30cm]
    \node[title] at (0, 8.5) {Camembert : Parts trimestrielles (Annuel)};
    % Quarterly Data
    \pgfmathsetmacro{\qone}{360}   % Q1
    \pgfmathsetmacro{\qtwo}{480}   % Q2
    \pgfmathsetmacro{\qthree}{410} % Q3
    \pgfmathsetmacro{\qfour}{480}  % Q4
    \pgfmathsetmacro{\total}{1730}
    \pgfmathsetmacro{\aone}{\qone/\total*360}
    \pgfmathsetmacro{\atwo}{\qtwo/\total*360}
    \pgfmathsetmacro{\athree}{\qthree/\total*360}
    \pgfmathsetmacro{\afour}{\qfour/\total*360}
    \pgfmathsetmacro{\r}{3.5}
    % Slices
    \fill[popblue!60, draw=white, thick] (0,0) -- (90:\r) arc (90:90-\aone:\r) -- cycle;
    \fill[poporange!60, draw=white, thick] (0,0) -- (90-\aone:\r) arc (90-\aone:90-\aone-\atwo:\r) -- cycle;
    \fill[popgreen!60, draw=white, thick] (0,0) -- (90-\aone-\atwo:\r) arc (90-\aone-\atwo:90-\aone-\atwo-\athree:\r) -- cycle;
    \fill[poppink!60, draw=white, thick] (0,0) -- (90-\aone-\atwo-\athree:\r) arc (90-\aone-\atwo-\athree:90-\aone-\atwo-\athree-\afour:\r) -- cycle;
    % Labels (approximate mid-angles)
    \pgfmathsetmacro{\mone}{90 - \aone/2}
    \pgfmathsetmacro{\mtwo}{90 - \aone - \atwo/2}
    \pgfmathsetmacro{\mthree}{90 - \aone - \atwo - \athree/2}
    \pgfmathsetmacro{\mfour}{90 - \aone - \atwo - \athree - \afour/2}
    \node[font=\tiny, text=white] at (\mone:1.8) {T1: 21\%};
    \node[font=\tiny, text=white] at (\mtwo:1.8) {T2: 28\%};
    \node[font=\tiny, text=white] at (\mthree:1.8) {T3: 24\%};
    \node[font=\tiny, text=white] at (\mfour:1.8) {T4: 28\%};
    % Legend
    \begin{scope}[yshift=-5cm, xshift=-3cm]
        \foreach \c/\t/\y in {popblue!60/T1 (21\%)/0, poporange!60/T2 (28\%)/-0.6, popgreen!60/T3 (24\%)/-1.2, poppink!60/T4 (28\%)/-1.8} {
            \fill[\c] (0,\y) rectangle (0.5,\y+0.4);
            \node[anchor=west, font=\tiny] at (0.7,\y+0.2) {\t};
        }
    \end{scope}
\end{scope}
\end{tikzpicture}
\legende{Galerie des 3 graphiques types pour les ventes 2024. \textbf{Gauche} : Histogramme (comparer mois à mois). \textbf{Centre} : Courbe (voir la baisse mars, pic octobre). \textbf{Droite} : Camembert (structure annuelle par trimestre).}
\end{popfigure}

\courssub{Mise en page, Aperçu et Impression du rapport}

Un graphique parfait mal imprimé (tronqué, sans en-têtes, trop petit) perd toute sa valeur.

\begin{methode}[Préparer l'impression d'un rapport tableur + graphiques]
\begin{enumerate}
    \item \textbf{Mise en page} $\to$ \textbf{Orientation} : \texttt{Paysage} (souvent mieux pour graphiques larges + tableau).
    \item \textbf{Mise en page} $\to$ \textbf{Taille} : \texttt{A4}.
    \item \textbf{Mise en page} $\to$ \textbf{Marges} : \texttt{Étroites} (maximise l'espace) ou \texttt{Personnalisées} (ex: Haut/Bas 1,5 cm, Gauche/Droite 1 cm).
    \item \textbf{Zone d'impression} : Sélectionner le tableau + graphiques $\to$ \texttt{Mise en page} $\to$ \texttt{Zone d'impression} $\to$ \texttt{Définir}.
    \item \textbf{Titres d'impression (CRUCIAL)} : \texttt{Mise en page} $\to$ \texttt{Titres d'impression} $\to$ \texttt{Lignes à répéter en haut} : Sélectionner la ligne d'en-tête (ex: \texttt{\$2:\$2}). \textit{Cela répète les noms de colonnes sur chaque page papier.}
    \item \textbf{Ajustement} : \texttt{Mise en page} $\to$ \texttt{Mise à l'échelle} $\to$ \texttt{Ajuster la feuille sur une page} (Largeur : 1 page, Hauteur : Automatique ou 1 page). \textbf{Vérifier} la lisibilité (police $\ge$ 8 pt).
    \item \textbf{Aperçu avant impression} : \texttt{Fichier} $\to$ \texttt{Imprimer} (ou \texttt{Ctrl+P}). Vérifier : Nombre de pages, coupures, graphiques coupés en deux.
    \item \textbf{Imprimer} : Choisir \texttt{Imprimer la sélection active} ou \texttt{Imprimer les pages actives}.
\end{enumerate}
\end{methode}

\begin{attention}[Graphiques coupés au saut de page]
Si un graphique chevauche une zone d'impression, il sera coupé en deux à l'impression.
\textbf{Solution} :
\begin{itemize}
    \item Déplacer/Redimensionner le graphique pour qu'il tienne dans la zone d'impression définie (lignes pointillées en mode \texttt{Aperçu des sauts de page} : \texttt{Affichage} $\to$ \texttt{Aperçu des sauts de page}).
    \item Ou : Clic droit graphique $\to$ \texttt{Taille et propriétés} $\to$ \texttt{Propriétés} $\to$ Cocher \texttt{Ne pas déplacer ni redimensionner avec les cellules} (fixe sa position absolue sur la feuille).
\end{itemize}
\end{attention}

\begin{savaistu}[Les Sparklines : Le graphique dans la cellule]
Inventées par Edward Tufte (expert visualisation), intégrées nativement dans **Excel (Insertion $\to$ Sparklines)\textbf{ et dispo dans }LibreOffice** via extension.
\begin{itemize}
    \item \textbf{Ligne} : Tendance mini (ex: évolution stock 12 mois dans une cellule).
    \item \textbf{Colonnes} : Comparaison mini.
    \item \textbf{Gain/Perte} : Binaire (Positif/Négatif).
\end{itemize}
\textbf{Usage pro} : Tableau de bord "Une page" pour un DG : 50 lignes produits, colonnes Mois, et une colonne \texttt{Sparkline} résumant l'année. Zéro place perdue, information dense.
\end{savaistu}

\begin{experience}[TP Guidé : Tableau de bord "Boutique École" - MFC + Graphique]
\textbf{Objectif} : Créer un suivi de stock avec alertes visuelles et graphique de valorisation.
\textbf{Durée} : 1 séance (50 min).
\textbf{Logiciel} : LibreOffice Calc ou Excel.
\textbf{Données de départ} (à saisir en \texttt{A1:E10}) :
\begin{center}
\begin{tabular}{|l|l|r|r|r|}\hline
\textbf{Code} & \textbf{Produit} & \textbf{Stock} & \textbf{Prix Unit.} & \textbf{Valorisation} \\ \hline
P001 & Cahiers 100p & 150 & 500 & \\ \hline
P002 & Stylos Bleu & 80 & 150 & \\ \hline
P003 & Règles 30cm & 5 & 200 & \\ \hline
P004 & Compas & 25 & 800 & \\ \hline
P005 & Gommes & 200 & 100 & \\ \hline
P006 & Crayons HB & 12 & 50 & \\ \hline
P007 & Feutres & 30 & 350 & \\ \hline
P008 & Chemises & 60 & 250 & \\ \hline
\end{tabular}
\end{center}

\textbf{Manipulation} :
\begin{enumerate}
    \item \textbf{Formule Valorisation} : En \texttt{E2} : \texttt{=C2*D2} $\to$ Étirer jusqu'\texttt{E9}.
    \item \textbf{MFC 1 (Rupture critique)} : Sélection \texttt{A2:E9} $\to$ Nouvelle règle $\to$ Formule : \texttt{=\$C2<10} $\to$ Format : Fond \textbf{Rouge}, Police \textbf{Blanc Gras}. (Repère P003, P006).
    \item \textbf{MFC 2 (Surstock)} : Même sélection $\to$ Nouvelle règle $\to$ Formule : \texttt{=\$C2>100} $\to$ Format : Fond \textbf{Vert clair}. (Repère P001, P005).
    \item \textbf{Barres de données} : Sélection \texttt{E2:E9} $\to$ MFC $\to$ Barres de données $\to$ Vert dégradé. (Visualise le capital immobilisé par produit).
    \item \textbf{Graphique} : Sélection \texttt{B1:E9} (Produit + Stock + Prix + Valorisation) $\to$ Insertion $\to$ Histogramme $\to$ \textbf{Colonnes groupées}.
    \item \textbf{Personnalisation Graphique} :
    \begin{itemize}
        \item Titre : \og \textbf{Valorisation du stock par produit (FCFA)} \fg.
        \item Axe vertical : \og Montant (FCFA) \fg.
        \item Série "Stock" et "Prix Unit." : Clic droit $\to$ \texttt{Format de la série de données} $\to$ \texttt{Axe secondaire} (pour Stock) ou masquer si trop d'échelles. \textit{Mieux : Graphique combiné (Colonnes Valorisation + Ligne Stock sur axe secondaire)}.
        \item Étiquettes de données sur Valorisation.
    \end{itemize}
    \item \textbf{Mise en page} : Orientation Paysage, Zone d'impression \texttt{A1:Graphique}, Titres d'impression ligne 1, Ajustement 1 page large.
\end{enumerate}

\textbf{Observations attendues} :
\begin{itemize}
    \item Les lignes P003 (Règles, stock 5) et P006 (Crayons, stock 12) sont en \textbf{Rouge/Blanc} $\to$ Commande urgente.
    \item P001 (Cahiers) et P005 (Gommes) en \textbf{Vert} $\to$ Surstock, promo possible.
    \item Barres dans colonne E : Cahiers et Gommes ont les plus longues barres (capital élevé).
    \item Graphique : La barre "Valorisation" domine pour Cahiers et Gommes. La ligne "Stock" (si axe secondaire) montre que Stylos ont un stock correct mais faible valorisation unitaire.
\end{itemize}

\textbf{Interprétation} : Le gérant décide de commander 200 Règles et 100 Crayons (Rouge), et de faire une promo "2 Cahiers achetés = 1 Gratuit" pour écouler le surstock (Vert).
\end{experience}

\begin{exercice}[Entraînement BEPC - MFC et Graphiques]
\textbf{Sujet} : Le fichier \texttt{Notes\_3eme.ods} contient les moyennes annuelles de 30 élèves (Colonnes : \texttt{A}=Nom, \texttt{B}=Moyenne).
\begin{enumerate}
    \item Écrire la formule exacte de MFC (règle personnalisée) pour colorer en \textbf{Vert} la ligne entière des élèves admis (Moyenne $\ge 10$). La sélection est \texttt{A2:B31}.
    \item On veut afficher une \textbf{flèche verte vers le haut} si Moyenne $\ge 14$, \textbf{flèche horizontale jaune} si $10 \le$ Moyenne $< 14$, \textbf{flèche rouge bas} si Moyenne $< 10$. Décrire la configuration exacte du \texttt{Jeu d'icônes} (Type, Seuils, Opérateurs).
    \item Quel type de graphique choisir pour visualiser la \textbf{distribution des notes} (Combien d'élèves entre 0-5, 5-10, 10-12, 12-14, 14-16, 16-20) ? Justifier.
    \item Dans la mise en page pour impression, à quoi sert l'option \texttt{Lignes à répéter en haut} ? Donner un cas concret où son absence pose problème.
\end{enumerate}
\end{exercice}

\begin{corrige}[Exercice Entraînement BEPC]
\begin{enumerate}
    \item \textbf{Formule MFC} : \texttt{=\$B2>=10}. (Sélection \texttt{A2:B31} $\to$ Nouvelle règle $\to$ Formule). \texttt{\$B} verrouille la colonne Moyenne, \texttt{2} est relatif (1ère ligne sélectionnée).
    \item \textbf{Jeu d'icônes} : Type \texttt{3 flèches (colorées)}.
    \begin{itemize}
        \item Icône 1 (Vert Haut) : Type \texttt{Nombre}, Valeur \texttt{14}, Opérateur \texttt{$\ge$}.
        \item Icône 2 (Jaune Latérale) : Type \texttt{Nombre}, Valeur \texttt{10}, Opérateur \texttt{$\ge$}.
        \item Icône 3 (Rouge Bas) : Automatique (reste $< 10$).
    \end{itemize}
    \item \textbf{Histogramme (Colonnes)}. La variable X (Classes de notes) est qualitative ordonnée, Y (Effectifs) est quantitative. On compare des effectifs par classe. La courbe (ligne) serait fausse (continuité illusoire), le camembert possible mais moins lisible pour comparer des hauteurs proches.
    \item \textbf{Rôle} : Répéter la ligne d'en-tête (Nom, Moyenne) en haut de \textbf{chaque page imprimée}. \textbf{Problème sans elle} : Sur la page 2, on voit juste des noms et des chiffres sans savoir quelle colonne est laquelle. Indispensable pour tout tableau de plus d'une page.
\end{enumerate}
\end{corrige}

\begin{aretenir}[Synthèse : MFC \& Graphiques]
\begin{itemize}
    \item \textbf{MFC} = Format \textbf{dynamique} piloté par une condition (valeur ou formule).
    \item \textbf{Règle personnalisée} : \texttt{=\$ColonneLigneActive > Seuil}. \texttt{\$} sur la \textbf{colonne de test}, \textbf{pas de \$} sur la \textbf{ligne active}.
    \item \textbf{Barres / Échelles / Icônes} : Analyse visuelle \textbf{in-cell} (sans graphique séparé). Seuils métier $\to$ Type \texttt{Nombre}.
    \item \textbf{Graphique} : \textbf{Type} dicté par nature des variables (Catégoriel+Quantitatif $\to$ Histogramme ; Temps+Quantitatif $\to$ Courbe ; Parts d'un tout $\to$ Camembert max 6 parts).
    \item \textbf{Qualité pro} : Titre contextuel, Axes nommés + unités, Étiquettes de données, Style sobre, Légende utile.
    \item \textbf{Impression} : \texttt{Zone d'impression}, \texttt{Titres d'impression (lignes à répéter)}, \texttt{Aperçu} obligatoire, \texttt{Ajustement 1 page large}.
\end{itemize}
\end{aretenir}

\coursec{Mise en page, Aperçu et Impression du rapport}

Après avoir mis en forme nos données, appliqué la mise en forme conditionnelle pour faire ressortir les informations clés et créé des graphiques parlants dans la section précédente, il reste une étape cruciale : \textbf{produire un document imprimable ou partageable de qualité}. Un tableau parfait à l'écran peut devenir illisible sur papier s'il est tronqué, s'il manque les en-têtes de colonnes sur la deuxième page, ou s'il s'étale sur dix feuilles inutiles. Cette section te donne la maîtrise complète de la chaîne « écran $\rightarrow$ papier/PDF ».

\courssub{Le mode « Mise en page » : voir avant d'imprimer}

Par défaut, le tableur affiche la vue \textbf{Normale}, optimisée pour la saisie et le calcul (grille quasi infinie, pas de marges visibles). Pour préparer l'impression, bascule en mode \textbf{Mise en page} via l'onglet \texttt{Affichage} $\rightarrow$ \texttt{Mise en page} (ou l'icône correspondante en bas à droite de la fenêtre).

\begin{popfigure}
\begin{tikzpicture}[scale=0.85, every node/.style={font=\small}]
% Fenêtre principale
\draw[popdark, thick, rounded corners=3pt] (0,0) rectangle (14,8);
\draw[popdark, fill=popblueL] (0,7.5) rectangle (14,8);
\node at (7,7.75) [text=popdark, font=\small\bfseries] {Onglets : Fichier  Accueil  Insertion  Mise en page  Formules  Données  Affichage};

% Ruban Mise en page actif
\draw[popblue, fill=popblue!20, rounded corners=2pt] (0.2,6.8) rectangle (13.8,7.3);
\node at (7,7.05) [text=popdark, font=\small\bfseries] {RUBAN MISE EN PAGE : Thèmes | Mise en page | Mise à l'échelle | Options de feuille | Organiser};

% Zone de travail avec sauts de page
\draw[popmuted, dashed] (1,1) rectangle (12,6.5);
\draw[poporange, thick, dashed] (4,1) -- (4,6.5);
\node[right, font=\tiny, text=poporange] at (4.1,3.7) {Saut de page vertical};
\draw[poporange, thick, dashed] (1,4) -- (12,4);
\node[above, font=\tiny, text=poporange] at (8,4.05) {Saut de page horizontal};

% Marges visibles
\draw[popgreen, thick, <->] (0.5,6.5) -- (1,6.5) node[midway, above, font=\tiny, text=popgreen] {Marge haut};
\draw[popgreen, thick, <->] (0.5,1) -- (1,1) node[midway, below, font=\tiny, text=popgreen] {Marge bas};

% Barre d'état
\draw[popdark, fill=popcream] (0,0) rectangle (14,0.5);
\node at (4,0.25) [font=\small] {Normal | \textbf{Mise en page} | Aperçu des sauts de page};
\node at (12,0.25) [font=\small] {100 \%};

% Légende
\node at (7,-0.5) [font=\small\bfseries, text=popdark] {Vue Mise en page : marges, en-têtes/pieds, sauts de page matérialisés en temps réel};
\end{tikzpicture}
\legende{Interface du tableur en mode Mise en page : visualisation immédiate des marges, des sauts de page et des en-têtes/pieds de page.}
\end{popfigure}

\begin{aretenir}[Pourquoi passer en mode Mise en page ?]
Ce mode affiche \textbf{exactement} ce qui sortira sur la feuille : marges physiques, en-têtes/pieds de page, sauts de page automatiques (lignes pointillées) et manuels. Tu peux \textbf{déplacer les sauts de page} à la souris pour forcer une coupure à un endroit précis (par exemple : ne pas couper une ligne de total).
\end{aretenir}

\courssub{Configuration de la page : Format, Orientation, Marges}

Ouvre la boîte de dialogue \textbf{Mise en page} : onglet \texttt{Mise en page} $\rightarrow$ groupe \texttt{Mise en page} $\rightarrow$ clic sur le \textbf{lanceur de boîte de dialogue} (petite flèche en bas à droite du groupe).

\begin{popfigure}
\begin{tikzpicture}[scale=0.8, every node/.style={font=\small}]
% Boîte de dialogue Mise en page - Onglet Page
\draw[popdark, thick, rounded corners=5pt] (0,0) rectangle (12,9);
\draw[popdark, fill=popblueL] (0,8.5) rectangle (12,9);
\node at (6,8.75) [font=\small\bfseries] {Mise en page};

% Onglets
\draw[popblue, fill=white, rounded corners=2pt] (0.3,8.0) rectangle (2.3,8.5);
\node at (1.3,8.25) {\textbf{Page}};
\draw[popmuted, fill=popcream, rounded corners=2pt] (2.5,8.0) rectangle (4.8,8.5);
\node at (3.65,8.25) {Marges};
\draw[popmuted, fill=popcream, rounded corners=2pt] (5.0,8.0) rectangle (8.0,8.5);
\node at (6.5,8.25) {En-tête/Pied};
\draw[popmuted, fill=popcream, rounded corners=2pt] (8.2,8.0) rectangle (10.5,8.5);
\node at (9.35,8.25) {Feuille};

% Contenu onglet Page
\node[align=left, anchor=north west, text width=11cm] at (0.5,7.7) {
\textbf{Orientation :} \quad \texttt{Portrait} \qquad \texttt{\color{popgreen} Paysage} (sélectionné) \\[4pt]
\textbf{Format papier :} \quad \texttt{A4 (21 x 29,7 cm)} \quad Lettre \quad Légal \\[4pt]
\textbf{Mise à l'échelle :} \\
\quad \texttt{\color{popblue} Ajuster à :} 1 page(s) de large $\times$ 1 page(s) de haut \\
\quad \texttt{Réduire/agrandir à :} \_\_ \% de la taille normale \\[4pt]
\textbf{Options :} \quad $\square$ Imprimer le quadrillage \quad $\square$ En-têtes de lignes/colonnes \\
\quad $\square$ Noir et blanc \quad $\square$ Qualité brouillon
};

% Boutons
\draw[popdark, fill=popblueL, rounded corners=2pt] (7.5,0.3) rectangle (8.8,0.8);
\node at (8.15,0.55) {OK};
\draw[popdark, fill=popblueL, rounded corners=2pt] (9.0,0.3) rectangle (10.5,0.8);
\node at (9.75,0.55) {Annuler};
\end{tikzpicture}
\legende{Boîte de dialogue Mise en page — Onglet « Page » : orientation, format papier et mise à l'échelle.}
\end{popfigure}

\begin{definition}[Orientation et Format]
\begin{itemize}
  \item \textbf{Portrait} (A4 : 21 cm de large $\times$ 29,7 cm de haut) : adapté aux listes longues comportant peu de colonnes (listes d'élèves, états de stock verticaux).
  \item \textbf{Paysage} (A4 : 29,7 cm de large $\times$ 21 cm de haut) : \textbf{indispensable pour les tableaux larges} (nombreuses colonnes : notes d'élèves avec toutes les matières, emplois du temps).
  \item \textbf{Marges} : \texttt{Normales} (2,5 cm), \texttt{Étroites} (environ 1,3 cm — permet de gagner de la place), \texttt{Larges}, \texttt{Personnalisées}. Pour un grand tableau, \texttt{Étroites} est souvent salutaire.
\end{itemize}
\end{definition}

\begin{methode}[Pour imprimer un grand tableau sur 1 page A4]
\begin{enumerate}
  \item \textbf{Orientation} : \texttt{Paysage} (onglet Page).
  \item \textbf{Marges} : \texttt{Étroites} (onglet Marges) ou \texttt{Personnalisées} (par exemple 1 cm partout).
  \item \textbf{Lignes à répéter} : onglet \texttt{Feuille} $\rightarrow$ \texttt{Lignes à répéter en haut} $\rightarrow$ sélectionner la ligne d'en-tête (par exemple \texttt{\$1:\$1}, ou \texttt{\$1:\$3} si l'en-tête occupe 3 lignes).
  \item \textbf{Mise à l'échelle} : onglet \texttt{Page} $\rightarrow$ \texttt{Ajuster à : 1 page(s) de large par 1 page(s) de haut}.
  \item \textbf{Vérifier} : \texttt{Ctrl + P} (Aperçu) $\rightarrow$ confirmer que tout tient, que l'en-tête revient sur chaque page et que la police reste lisible (au moins 8 pt).
\end{enumerate}
\end{methode}

\courssub{La mise à l'échelle : « Ajuster à 1 page » ou « Pourcentage »}

Deux stratégies s'offrent à toi dans l'onglet \texttt{Page} :

\begin{propriete}[Mise à l'échelle — Règles de choix]
\begin{itemize}
  \item \textbf{Ajuster à : 1 page(s) de large $\times$ 1 page(s) de haut} (recommandé) : le tableur \textbf{réduit automatiquement} le facteur d'échelle pour que toute la zone d'impression tienne sur une seule feuille. Idéal pour « tout sur une page ».
  \item \textbf{Réduire/agrandir à : X \% de la taille normale} : tu fixes le zoom (par exemple 75 \%). Utile si tu veux garder une taille de police précise, mais risque de tronquer le tableau ou de laisser du vide.
  \item \textbf{Attention} : si le tableau est \textbf{très grand} (par exemple 100 colonnes $\times$ 200 lignes), « Ajuster 1$\times$1 » rendra le texte microscopique. Dans ce cas, accepte plusieurs pages et utilise \texttt{Lignes à répéter en haut}.
\end{itemize}
\end{propriete}

\courssub{En-têtes et pieds de page : identifier le document}

Onglet \texttt{En-tête/Pied de page} de la boîte Mise en page $\rightarrow$ bouton \texttt{Personnaliser l'en-tête...} ou \texttt{Personnaliser le pied de page...}. Trois zones sont disponibles : \textbf{Gauche}, \textbf{Centre}, \textbf{Droite}. Les boutons d'insertion permettent d'ajouter des informations automatiques :

\begin{center}
\begin{tabularx}{\linewidth}{|l|X|l|}
\hline
\rowcolor{popblueL}
\textbf{Bouton} & \textbf{Code inséré} & \textbf{Rendu exemple} \\
\hline
Numéro de page & \texttt{\&[Page]} & Page 3 \\
\hline
Nombre de pages & \texttt{\&[Pages]} & / 12 \\
\hline
Date & \texttt{\&[Date]} & 15/03/2026 \\
\hline
Heure & \texttt{\&[Heure]} & 14:30 \\
\hline
Nom du fichier & \texttt{\&[Fichier]} & Notes\_3eme\_A.xlsx \\
\hline
Nom de la feuille & \texttt{\&[Feuille]} & Trimestre\_1 \\
\hline
Image/Logo & \texttt{\&[Image]} & [Logo du collège] \\
\hline
\end{tabularx}
\end{center}

\begin{exemplebox}[En-tête / Pied de page standard pour un bulletin officiel]
\begin{itemize}
  \item \textbf{En-tête Centre} : \texttt{COLLÈGE BILINGUE DE NGOUSSO — BULLETIN TRIMESTRIEL}
  \item \textbf{Pied Gauche} : \texttt{\&[Fichier] — \&[Feuille]}
  \item \textbf{Pied Centre} : \texttt{Imprimé le \&[Date] à \&[Heure]}
  \item \textbf{Pied Droite} : \texttt{Page \&[Page] / \&[Pages]}
\end{itemize}
\end{exemplebox}

\courssub{Titres d'impression : répéter les en-têtes sur chaque page}

C'est \textbf{la} fonction qui sépare l'amateur de l'utilisateur soigneux. Sans elle, la page 2 d'un tableau de 50 lignes n'affiche plus les noms des colonnes (Nom, Maths, Physique, Moyenne...).

\begin{popfigure}
\begin{tikzpicture}[scale=0.85, every node/.style={font=\small}]
% Boîte Mise en page - Onglet Feuille
\draw[popdark, thick, rounded corners=5pt] (0,0) rectangle (12,9);
\draw[popdark, fill=popblueL] (0,8.5) rectangle (12,9);
\node at (6,8.75) [font=\small\bfseries] {Mise en page};

% Onglets
\draw[popmuted, fill=popcream, rounded corners=2pt] (0.3,8.0) rectangle (2.3,8.5);
\node at (1.3,8.25) {Page};
\draw[popmuted, fill=popcream, rounded corners=2pt] (2.5,8.0) rectangle (4.8,8.5);
\node at (3.65,8.25) {Marges};
\draw[popmuted, fill=popcream, rounded corners=2pt] (5.0,8.0) rectangle (8.0,8.5);
\node at (6.5,8.25) {En-tête/Pied};
\draw[popblue, fill=white, rounded corners=2pt] (8.2,8.0) rectangle (10.5,8.5);
\node at (9.35,8.25) {\textbf{Feuille}};

% Contenu onglet Feuille
\node[align=left, anchor=north west, text width=11cm] at (0.5,7.7) {
\textbf{Zone d'impression :} \texttt{\$A\$1:\$O\$50} \\[4pt]
\textbf{Titres d'impression :} \\
\quad \texttt{Lignes à répéter en haut :} \texttt{\$1:\$3} \quad (cliquer les lignes 1 à 3) \\
\quad \texttt{Colonnes à répéter à gauche :} \texttt{\$A:\$A} \quad (pour garder la colonne Nom) \\[4pt]
\textbf{Imprimer :} \\
\quad $\square$ Quadrillage \quad $\square$ En-têtes de lignes et colonnes (A, B, C... 1, 2, 3...) \\
\quad $\square$ Commentaires \quad $\square$ Noir et blanc \quad $\square$ Qualité brouillon \\[4pt]
\textbf{Ordre des pages :} \quad \texttt{\color{popblue} Descendant, puis horizontal} (par défaut) \\
\quad \texttt{Horizontal, puis descendant} (pour grands tableaux paysage multi-pages)
};

% Boutons
\draw[popdark, fill=popblueL, rounded corners=2pt] (7.5,0.3) rectangle (8.8,0.8);
\node at (8.15,0.55) {OK};
\draw[popdark, fill=popblueL, rounded corners=2pt] (9.0,0.3) rectangle (10.5,0.8);
\node at (9.75,0.55) {Annuler};
\end{tikzpicture}
\legende{Onglet « Feuille » : définition de la zone d'impression, lignes/colonnes à répéter, options d'impression et ordre des pages.}
\end{popfigure}

\begin{definition}[Lignes à répéter en haut / Colonnes à répéter à gauche]
\begin{itemize}
  \item \textbf{Lignes à répéter en haut} : plage de lignes (par exemple \texttt{\$1:\$3}) imprimée en haut de \textbf{chaque page}. Syntaxe : \texttt{\$début:\$fin} (références absolues obligatoires).
  \item \textbf{Colonnes à répéter à gauche} : plage de colonnes (par exemple \texttt{\$A:\$B}) imprimée à gauche de chaque page. Utile pour garder l'identifiant (Nom, Matricule) visible quand le tableau déborde horizontalement.
  \item Ces titres ne s'affichent \textbf{pas} à l'écran en mode Normal : ils n'apparaissent qu'à l'impression et dans l'Aperçu.
\end{itemize}
\end{definition}

\courssub{Zone d'impression : ne pas gaspiller le papier}

Par défaut, le tableur imprime \textbf{toutes les cellules utilisées}, y compris les cellules vides lointaines qu'il considère comme « touchées ». La \textbf{Zone d'impression} restreint la sortie à une plage précise.

\begin{methode}[Définir / Effacer la zone d'impression]
\begin{enumerate}
  \item \textbf{Sélectionner} la plage utile (par exemple \texttt{A1:O51} pour 50 élèves + l'en-tête).
  \item Onglet \texttt{Mise en page} $\rightarrow$ groupe \texttt{Mise en page} $\rightarrow$ \texttt{Zone d'impression} $\rightarrow$ \texttt{Définir}.
  \item \textbf{Vérifier} : onglet \texttt{Affichage} $\rightarrow$ \texttt{Aperçu des sauts de page} $\rightarrow$ seule la zone sélectionnée apparaît en blanc, le reste est grisé.
  \item Pour \textbf{effacer} : \texttt{Zone d'impression} $\rightarrow$ \texttt{Effacer}.
  \item \textbf{Plusieurs zones} : sélectionner une plage, \texttt{Définir}, puis sélectionner une autre plage (non contiguë) et choisir \texttt{Ajouter à la zone d'impression}. Chaque zone s'imprime sur une page séparée.
\end{enumerate}
\end{methode}

\begin{attention}[Pièges classiques — Zone d'impression et Titres d'impression]
\begin{itemize}
  \item \textbf{Oublier « Lignes à répéter en haut »} : les pages 2, 3... n'ont plus d'en-têtes $\rightarrow$ document inutilisable pour la lecture.
  \item \textbf{Ne pas définir la zone d'impression} : impression de nombreuses pages blanches après tes 50 lignes de données (le tableur mémorise la dernière cellule modifiée, même vide).
  \item \textbf{Zone d'impression trop petite} : tu as sélectionné \texttt{A1:O50} mais tes données vont jusqu'en ligne 55 $\rightarrow$ 5 élèves manquent sur le bulletin.
  \item \textbf{Références relatives dans « Lignes à répéter »} : écrire \texttt{1:3} au lieu de \texttt{\$1:\$3} peut provoquer une erreur ou un comportement imprévisible. \textbf{Toujours utiliser le signe \$}.
\end{itemize}
\end{attention}

\courssub{Options d'impression : quadrillage, en-têtes, ordre des pages}

Toujours dans l'onglet \texttt{Feuille} de la boîte Mise en page :

\begin{propriete}[Options d'impression — Effets]
\begin{itemize}
  \item \textbf{Imprimer le quadrillage} : imprime les fines lignes grises entre les cellules. \textbf{Utile} pour les tableaux bruts sans bordures appliquées. \textbf{Inutile} si tu as déjà mis des bordures (onglet Accueil $\rightarrow$ Bordures).
  \item \textbf{Imprimer les en-têtes de lignes et colonnes} : imprime \texttt{A, B, C...} et \texttt{1, 2, 3...} sur les bords de la feuille. \textbf{Réservé aux documents techniques} (vérification, formation), \textbf{pas} pour un bulletin officiel.
  \item \textbf{Ordre des pages} :
    \begin{itemize}
      \item \texttt{Descendant, puis horizontal} (défaut) : page 1 = haut-gauche, page 2 = dessous, puis on passe à droite. Adapté aux tableaux longs (portrait).
      \item \texttt{Horizontal, puis descendant} : page 1 = haut-gauche, page 2 = à droite, puis on descend. \textbf{Préférable en Paysage} pour les grands tableaux : on lit la suite à droite avant de passer en bas.
    \end{itemize}
\end{itemize}
\end{propriete}

\courssub{Aperçu avant impression et Export PDF}

Avant toute impression physique, \textbf{ouvre l'Aperçu} : \texttt{Fichier} $\rightarrow$ \texttt{Imprimer} ou \textbf{Ctrl + P}.

\begin{savaistu}[Raccourci universel]
\textbf{Ctrl + P} (Windows/Linux) ou \textbf{Cmd + P} (Mac) ouvre instantanément l'aperçu avant impression avec tous les réglages rapides : choix de l'imprimante, plage de pages (par exemple \texttt{1-3,5}), recto/verso, nombre de copies, couleur ou noir et blanc, et même le \textbf{format PDF} proposé comme « imprimante » virtuelle.
\end{savaistu}

\begin{popfigure}
\begin{tikzpicture}[scale=0.75, every node/.style={font=\small}]
% Fenêtre Aperçu avant impression (Ctrl+P)
\draw[popdark, thick, rounded corners=3pt] (0,0) rectangle (15,9);
\draw[popdark, fill=popblueL] (0,8.5) rectangle (15,9);
\node at (7.5,8.75) [font=\small\bfseries] {Imprimer — Aperçu avant impression   [Imprimante : Microsoft Print to PDF]};

% Panneau gauche (réglages)
\draw[popmuted] (0,0.5) rectangle (4.5,8.5);
\node at (2.25,8.1) [font=\small\bfseries] {RÉGLAGES};
\node[align=left, anchor=north west, text width=4cm] at (0.3,7.8) {
\textbf{Imprimante :} Microsoft Print to PDF \\
\textbf{Pages :} \texttt{Toutes} \\
\textbf{Copies :} 1 \\
\textbf{Format :} A4 \\
\textbf{Orientation :} \texttt{\color{popgreen} Paysage} \\
\textbf{Marges :} \texttt{Étroites} \\
\textbf{Échelle :} \texttt{Ajuster la feuille sur une page}
};

% Aperçu page (tableau paysage tenu sur 1 page)
\draw[popdark, thick] (5,1) rectangle (14,8);
% En-tête
\draw[popblue, fill=popblue!10] (5,7.2) rectangle (14,7.8);
\node at (9.5,7.5) [font=\small\bfseries] {COLLÈGE BILINGUE DE NGOUSSO — BULLETIN T1 — 3ème A};
% Tableau
\foreach \y in {6.8, 6.4, 6.0, 5.6, 5.2, 4.8, 4.4, 4.0, 3.6, 3.2, 2.8, 2.4, 2.0, 1.6, 1.2} {
  \draw[popmuted, very thin] (5.2,\y) -- (13.8,\y);
}
\foreach \x in {5.2, 6.2, 7.2, 8.2, 9.2, 10.2, 11.2, 12.2, 13.2, 13.8} {
  \draw[popmuted, very thin] (\x,1.2) -- (\x,7.1);
}
% En-têtes colonnes
\node at (5.7,6.95) [font=\tiny\bfseries] {N°};
\node at (6.7,6.95) [font=\tiny\bfseries] {Nom};
\node at (7.7,6.95) [font=\tiny\bfseries] {Prénom};
\node at (8.7,6.95) [font=\tiny\bfseries] {Maths};
\node at (9.7,6.95) [font=\tiny\bfseries] {Phys.};
\node at (10.7,6.95) [font=\tiny\bfseries] {Info.};
\node at (11.7,6.95) [font=\tiny\bfseries] {SVT};
\node at (12.7,6.95) [font=\tiny\bfseries] {Angl.};
\node at (13.5,6.95) [font=\tiny\bfseries] {Moy.};
% Pied de page
\draw[popblue, fill=popblue!10] (5,1) rectangle (14,1.2);
\node at (9.5,1.1) [font=\tiny] {Page 1 / 1  —  Imprimé le 15/03/2026 à 14:30  —  Notes\_3eme\_A.xlsx};
\end{tikzpicture}
\legende{Aperçu avant impression (Ctrl+P) : un tableau de 50 lignes $\times$ 15 colonnes tenu sur 1 page A4 Paysage grâce à « Ajuster 1$\times$1 » et « Lignes à répéter en haut ». En-tête et pied de page visibles.}
\end{popfigure}

\courssub{Export PDF : diffusion numérique officielle}

Pour transmettre un bulletin, un état de stock ou un rapport à l'administration, au chef d'établissement ou aux parents, le \textbf{PDF} est le standard universel : le document est inaltérable, correctement paginé et lisible sur tout appareil (ordinateur, téléphone).

\begin{methode}[Exporter en PDF de qualité]
\begin{enumerate}
  \item \texttt{Fichier} $\rightarrow$ \texttt{Exporter} $\rightarrow$ \texttt{Créer un document PDF/XPS} $\rightarrow$ \texttt{Créer PDF/XPS}.
  \item Dans la boîte de dialogue : choisir le dossier, nommer le fichier (par exemple \texttt{Bulletin\_3eme\_A\_T1\_2026.pdf}).
  \item Cliquer \textbf{Options...} (étape cruciale) :
    \begin{itemize}
      \item \textbf{Plage} : \texttt{Tout le classeur} (toutes les feuilles), \texttt{Feuille active} ou \texttt{Page(s) 1 à 3}.
      \item \textbf{Inclure} : $\square$ \texttt{Signets} (crée une table des matières cliquable à partir des noms de feuilles) ; $\square$ \texttt{Propriétés du document} (auteur, titre).
      \item \textbf{Qualité} : \texttt{Standard (publication en ligne et impression)} — éviter « Taille minimale » qui dégrade les graphiques.
    \end{itemize}
  \item \texttt{OK} $\rightarrow$ \texttt{Publier}.
  \item \textbf{Ouvrir le PDF généré} pour vérifier : en-têtes répétés, graphiques nets, pas de coupure sauvage, numérotation correcte.
\end{enumerate}
\end{methode}

\begin{exemplebox}[Tableau 50 lignes $\times$ 15 colonnes — Comparatif de rendu]
\begin{center}
\begin{tabularx}{\linewidth}{|X|l|X|l|}
\hline
\rowcolor{popblueL}
\textbf{Configuration} & \textbf{Nb pages} & \textbf{Lisibilité} & \textbf{Utilisable ?} \\
\hline
Défaut (Portrait, marges normales, pas d'échelle, pas de répétition) & 4 (tronqué à droite et en bas) & En-têtes perdus dès la page 2, colonnes coupées & \textcolor{poporange}{Non} \\
\hline
Paysage + marges étroites + Ajuster 1$\times$1 + répéter ligne 1 & \textbf{1} & Complète, en-tête présent, police environ 9 pt & \textcolor{popgreen}{\textbf{Oui}} \\
\hline
Paysage + marges normales + 75 \% + répéter ligne 1 & 2 (largeur OK, hauteur coupée) & Correcte mais 2 pages & \textcolor{poporange}{Moyen} \\
\hline
Portrait + Ajuster 1$\times$1 (forcé) & 1 & Police environ 6 pt, illisible, colonnes écrasées & \textcolor{poporange}{Non} \\
\hline
\end{tabularx}
\end{center}
\textbf{Conclusion} : la combinaison \texttt{Paysage + marges étroites + Ajuster 1$\times$1 + répéter l'en-tête} est la « formule magique » pour la grande majorité des tableaux scolaires larges.
\end{exemplebox}

\courssub{Exercice résolu : Préparer le bulletin officiel de la 3ème A}

\begin{exoresolu}[Bulletin trimestriel — Mise en page complète]
\textbf{Énoncé :} Tu as un classeur \texttt{Bulletins\_3eme\_A.xlsx} avec la feuille \texttt{Trimestre\_1} contenant : lignes 1 à 3 = en-tête de l'établissement + titre + noms des 12 matières ; lignes 4 à 53 = 50 élèves (colonne A = N°, B = Nom, C = Prénom, D à O = notes et moyennes). Tu dois produire un PDF d'\textbf{1 page A4} prêt pour la signature du principal.
\tcblower
\textbf{Solution détaillée :}
\begin{enumerate}
  \item \textbf{Mode Mise en page} : \texttt{Affichage} $\rightarrow$ \texttt{Mise en page}. Vérifier que les données tiennent dans \texttt{A1:O53}.
  \item \textbf{Zone d'impression} : sélectionner \texttt{A1:O53} $\rightarrow$ \texttt{Mise en page} $\rightarrow$ \texttt{Zone d'impression} $\rightarrow$ \texttt{Définir}.
  \item \textbf{Boîte Mise en page} (lanceur) $\rightarrow$ onglet \texttt{Page} :
    \begin{itemize}
      \item Orientation : \texttt{Paysage}
      \item Format : \texttt{A4}
      \item Mise à l'échelle : \texttt{Ajuster à : 1 page(s) de large $\times$ 1 page(s) de haut}
    \end{itemize}
  \item Onglet \texttt{Marges} : \texttt{Étroites} (ou Personnalisées : Haut/Bas/Gauche/Droite = 1 cm).
  \item Onglet \texttt{Feuille} :
    \begin{itemize}
      \item \texttt{Lignes à répéter en haut} : cliquer dans le champ puis sélectionner les lignes \texttt{1:3} $\rightarrow$ Entrée $\rightarrow$ affiche \texttt{\$1:\$3}.
      \item \texttt{Colonnes à répéter à gauche} : \texttt{\$A:\$B} (garder N° et Nom visibles en cas de débordement horizontal).
      \item $\square$ \texttt{Quadrillage} (coché si aucune bordure n'a été appliquée) ; \texttt{En-têtes lignes/colonnes} \textbf{décoché}.
      \item Ordre : \texttt{Horizontal, puis descendant} (préférable en paysage).
    \end{itemize}
  \item Onglet \texttt{En-tête/Pied de page} $\rightarrow$ \texttt{Personnaliser...} :
    \begin{itemize}
      \item En-tête Centre : \texttt{COLLÈGE DE NGOUSSO — BULLETIN 1er TRIMESTRE 2025-2026 — 3ème A}
      \item Pied Gauche : \texttt{\&[Fichier] — \&[Feuille]}
      \item Pied Centre : \texttt{Édité le \&[Date] à \&[Heure]}
      \item Pied Droite : \texttt{Page \&[Page] / \&[Pages]}
    \end{itemize}
  \item \textbf{Ctrl + P} $\rightarrow$ vérifier l'aperçu : 1 page, en-tête visible, toutes les colonnes lisibles, pied de page correct.
  \item \textbf{Export PDF} : \texttt{Fichier} $\rightarrow$ \texttt{Exporter} $\rightarrow$ \texttt{Créer PDF/XPS} $\rightarrow$ \texttt{Options} $\rightarrow$ \texttt{Feuille active} + \texttt{Signets} + \texttt{Standard} $\rightarrow$ \texttt{Publier} $\rightarrow$ \texttt{Bulletin\_3eme\_A\_T1\_2026.pdf}.
  \item \textbf{Test final} : ouvrir le PDF $\rightarrow$ imprimer une page test sur l'imprimante du secrétariat $\rightarrow$ valider avec le surveillant général.
\end{enumerate}
\textbf{Résultat} : un PDF d'1 page A4 Paysage, 50 élèves $\times$ 15 colonnes, en-tête institutionnel répété, numérotation, date/heure, nom du fichier, prêt pour diffusion officielle.
\end{exoresolu}

\courssub{TP guidé : De l'écran au PDF officiel}

\begin{experience}[TP : Publier le rapport de stock de la cantine]
\textbf{Objectif} : transformer un tableau brut de stock (80 articles $\times$ 10 colonnes) en un PDF professionnel d'1 page A4 Paysage.

\textbf{Matériel} : fichier \texttt{Stock\_Cantine\_Brute.xlsx} (fourni), un tableur (LibreOffice Calc ou Excel), une imprimante PDF.

\textbf{Manipulation pas à pas} :
\begin{enumerate}
  \item Ouvrir le fichier. Constater : mode Normal, données en \texttt{A1:J81}, pas de bordures, pas d'en-tête/pied, orientation Portrait.
  \item \texttt{Affichage} $\rightarrow$ \texttt{Mise en page}. Observer les sauts de page (plusieurs pages).
  \item Sélectionner \texttt{A1:J81} $\rightarrow$ \texttt{Mise en page} $\rightarrow$ \texttt{Zone d'impression} $\rightarrow$ \texttt{Définir}.
  \item Lanceur de la boîte Mise en page $\rightarrow$ onglet \texttt{Page} : \texttt{Paysage}, \texttt{A4}, \texttt{Ajuster 1$\times$1}.
  \item Onglet \texttt{Marges} : \texttt{Étroites}.
  \item Onglet \texttt{Feuille} : \texttt{Lignes à répéter en haut} $\rightarrow$ \texttt{\$1:\$1} (ligne unique d'en-tête : Article, Unité, Qté Entrée, Qté Sortie, Stock, Prix, Valeur, Seuil, Fournisseur, Date Maj). \texttt{Colonnes à répéter} $\rightarrow$ \texttt{\$A:\$A} (nom de l'article). $\square$ \texttt{Quadrillage}. Ordre : \texttt{Horizontal, puis descendant}.
  \item Onglet \texttt{En-tête/Pied} : en-tête Centre \texttt{CANTINE SCOLAIRE — ÉTAT DES STOCKS AU 15/03/2026}. Pied Droite \texttt{Page \&[Page] / \&[Pages]}.
  \item \texttt{Ctrl + P} : vérifier l'aperçu. Si la police est trop petite (moins de 8 pt) : passer en \texttt{Marges Personnalisées 0,8 cm} ou accepter 2 pages (désactiver \texttt{Ajuster 1$\times$1}, mettre \texttt{85 \%}, garder \texttt{Lignes à répéter}).
  \item \texttt{Fichier} $\rightarrow$ \texttt{Exporter} $\rightarrow$ \texttt{PDF} $\rightarrow$ Options : \texttt{Feuille active}, \texttt{Signets}, \texttt{Standard} $\rightarrow$ Publier \texttt{Stock\_Cantine\_2026-03-15.pdf}.
  \item Ouvrir le PDF, vérifier : total du stock en bas visible, en-tête répété s'il y a 2 pages, logo de la cantine en en-tête gauche (optionnel : \texttt{Personnaliser l'en-tête} $\rightarrow$ bouton Image).
\end{enumerate}

\textbf{Observations attendues} :
\begin{itemize}
  \item Le mode Mise en page montre les sauts de page se réduire au fur et à mesure des réglages.
  \item « Ajuster 1$\times$1 » force l'échelle à environ 65-75 \% selon la largeur du tableau.
  \item Sans « Lignes à répéter », la page 2 (s'il y a 2 pages) n'a pas les noms des colonnes.
  \item Le PDF est léger (environ 50 Ko) et s'ouvre instantanément sur le téléphone du gestionnaire.
\end{itemize}

\textbf{Interprétation} : la chaîne \textbf{Zone d'impression $\rightarrow$ Orientation/Marges $\rightarrow$ Échelle $\rightarrow$ Titres d'impression $\rightarrow$ En-tête/Pied $\rightarrow$ PDF} est reproductible pour \textbf{tout} rapport tabulaire scolaire ou administratif.
\end{experience}

\begin{exercice}[Entraînement — Mise en page et PDF]
\begin{enumerate}
  \item Crée un classeur \texttt{Emploi\_Temps\_3eme.xlsx} avec une feuille par classe (3ème A, B, C). Chaque feuille : grille 6 jours $\times$ 7 créneaux (lignes = créneaux, colonnes = jours), matières dans les cellules. Mise en forme conditionnelle : Maths = rouge, Français = bleu, Sport = vert.
  \item Pour la feuille \texttt{3ème A} : définis la zone d'impression, passe en Paysage, marges étroites, ajuste 1$\times$1, répète la ligne 1 (jours) en haut, ajoute l'en-tête « EMPLOI DU TEMPS 2025-2026 — 3ème A » et le pied « Page x/y — Établissement X ».
  \item Exporte en PDF \texttt{EDT\_3eme\_A.pdf}. Vérifie que les couleurs de la mise en forme conditionnelle sortent en couleur (pas en noir et blanc).
  \item \textbf{Défi} : le tableau ne tient pas sur 1 page lisible (police 7 pt). Propose deux solutions alternatives (par exemple : 2 pages avec répétition, ou police 9 pt + orientation portrait + masquer les colonnes vides). Teste et choisis la meilleure.
\end{enumerate}
\end{exercice}

\begin{corrige}[Exercice n°1-4]
\begin{enumerate}
  \item Structure créée : 3 feuilles, grille 7 lignes $\times$ 6 colonnes de jours (jours en ligne 1, créneaux en colonne A), mise en forme conditionnelle appliquée sur les matières.
  \item 3ème A : zone \texttt{A1:G8} (6 jours + 1 colonne de créneaux, 7 créneaux + 1 ligne de jours). Paysage, marges étroites, Ajuster 1$\times$1, \texttt{\$1:\$1} répété. En-tête et pied de page renseignés.
  \item Export PDF : Options $\rightarrow$ \texttt{Feuille active}, \texttt{Standard}. Vérification : couleurs préservées car l'imprimante PDF virtuelle imprime en couleur par défaut.
  \item \textbf{Solutions au défi} :
    \begin{itemize}
      \item \textbf{Solution A (2 pages)} : désactiver « Ajuster 1$\times$1 », mettre 90 \%, garder Paysage + répéter \texttt{\$1:\$1}. $\rightarrow$ 2 pages lisibles (police environ 10 pt), en-tête des jours sur chaque page.
      \item \textbf{Solution B (Portrait)} : orientation Portrait, marges étroites, Ajuster 1$\times$1. $\rightarrow$ 1 page mais police environ 8 pt, colonnes des jours étroites.
      \item \textbf{Choix recommandé} : Solution A (2 pages) pour un emploi du temps affiché en salle de classe — la lisibilité prime sur la compacité. Solution B acceptable pour une consultation à l'écran uniquement.
    \end{itemize}
\end{enumerate}
\end{corrige}

\begin{aretenir}[Check-list finale avant diffusion d'un rapport tableur]
\begin{itemize}
  \item[$\square$] Zone d'impression définie (pas de pages vides).
  \item[$\square$] Orientation adaptée (Paysage pour un tableau large, Portrait pour un tableau long).
  \item[$\square$] Marges optimisées (Étroites ou Personnalisées).
  \item[$\square$] Mise à l'échelle : « Ajuster 1$\times$1 » OU pourcentage lisible (police d'au moins 8 pt).
  \item[$\square$] \textbf{Lignes à répéter en haut} définies (\texttt{\$1:\$n}).
  \item[$\square$] En-tête/Pied de page : identification (établissement, titre, date, page x/y, nom du fichier).
  \item[$\square$] Options : quadrillage si pas de bordures, PAS d'en-têtes A/1, ordre des pages cohérent.
  \item[$\square$] Aperçu (Ctrl+P) validé visuellement (toutes les pages).
  \item[$\square$] Export PDF : Options $\rightarrow$ plage correcte, signets, qualité Standard.
  \item[$\square$] Test d'ouverture du PDF sur un autre poste ou un téléphone.
\end{itemize}
\end{aretenir}

\newpage

\coursec{Activité d'intégration}

\coursec{Utilisation d'un tableur : Activité d'intégration — Tableau de bord de la cantine}

\courssub{Objectifs de l'activité}
Cette activité d'intégration vise à mobiliser l'ensemble des compétences acquises au cours de la leçon « Utilisation d'un tableur » dans une situation professionnelle réaliste. À l'issue de ce travail, tu seras capable de :
\begin{itemize}
    \item Structurer et mettre en forme un classeur selon des standards de gestion administrative (en-têtes, bordures, formats monétaires, gel des volets).
    \item Concevoir des formules de calcul robustes en distinguant les références relatives et absolues.
    \item Exploiter les fonctions mathématiques (\texttt{SOMME}, \texttt{MOYENNE}, \texttt{MAX}) et statistiques (\texttt{NB.SI}) pour synthétiser des données.
    \item Modéliser une logique de décision métier à l'aide des fonctions logiques \texttt{SI} et \texttt{SI.MULTIPLE} (imbrication de conditions).
    \item Appliquer une mise en forme conditionnelle (MFC) dynamique : barres de données et coloration de lignes entières selon une règle logique.
    \item Créer et paramétrer un graphique combiné (colonnes groupées avec axe secondaire) pour une analyse visuelle comparative.
    \item Configurer une mise en page professionnelle (orientation, zones d'impression, en-têtes/pieds de page, répétition des titres) et exporter le document final en PDF normé.
\end{itemize}

\courssub{Situation}
\textbf{Contexte :} Tu es stagiaire au service de l'intendance du \textbf{Lycée Bilingue de Nkolbisson} (Yaoundé). Le chef de service, M. \textsc{Atangana}, te confie la mission de numériser le suivi hebdomadaire de la cantine scolaire afin de remplacer le cahier de bord manuscrit, source d'erreurs et de pertes de temps. L'objectif est de produire un \textbf{Tableau de bord hebdomadaire} clair, automatisé et prêt à être imprimé pour la direction chaque vendredi.

\textbf{Données brutes de la semaine du 13 au 17 janvier 2025 :} Le fichier brut (ou la saisie guidée) contient le relevé quotidien des repas. Le prix unitaire est fixé par la direction pour chaque plat.

\begin{center}
\begin{tabularx}{\linewidth}{|c|X|c|c|c|}
\hline
\rowcolor{popblueL}
\textbf{Date} & \textbf{Plat principal} & \textbf{Qté préparée} & \textbf{Qté servie} & \textbf{Prix unitaire (FCFA)} \\
\hline
Lundi 13/01 & Ndolé & 120 & 110 & 500 \\
Mardi 14/01 & Riz sauce arachide & 100 & 100 & 450 \\
Mercredi 15/01 & Poulet DG & 80 & 60 & 1\,000 \\
Jeudi 16/01 & Spaghetti bolognaise & 90 & 85 & 400 \\
Vendredi 17/01 & Poisson braisé + Plantains & 110 & 95 & 600 \\
\hline
\end{tabularx}
\end{center}

\textbf{Contraintes de gestion :}
\begin{itemize}
    \item Une « Rupture stock » est signalée si la quantité servie égale la quantité préparée (tout a été écoulé, risque de mécontentement des derniers élèves).
    \item L'affluence est catégorisée sur la quantité servie : $\ge 100$ (Très forte), $\ge 70$ (Forte), $\ge 40$ (Moyenne), sinon (Faible).
    \item Le rapport final doit tenir sur \textbf{une seule page A4} en orientation Paysage.
\end{itemize}

\courssub{Consignes}
Réalise les tâches ci-dessous dans l'ordre, en nommant ton fichier \texttt{Tableau\_Bord\_Cantine\_PrenomNom.ods} (ou \texttt{.xlsx}) :

\begin{enumerate}
    \item \textbf{Mise en forme du tableau de saisie (plage A1:E6) :} Fusionne et centre la ligne 1 sur A1:E1 avec le titre « \textsc{Suivi Hebdomadaire Cantine - Semaine du 13/01/2025} » (Police 14, Gras, Fond bleu foncé, Police blanche). Met en forme la ligne 2 (en-têtes de colonnes) : Gras, Fond bleu clair, Bordures fines extérieures/intérieures, Alignement centré, Hauteur de ligne 30. Applique le format « Nombre personnalisé » \texttt{\# \#\#0 " FCFA"} à la colonne E. Gèle les volets sur la ligne 2.
    \item \textbf{Chiffre d'affaires journalier (Colonne F - intitulé « CA Jour (FCFA) ») :} En F3, écris la formule calculant le produit \texttt{Qté servie $\times$ Prix unitaire}. Étends vers le bas. Vérifie que la référence sur le Prix unitaire est correcte (relative ici car prix variables).
    \item \textbf{Taux de service (Colonne G - intitulé « Taux Service (\%) ») :} En G3, calcule \texttt{Qté servie / Qté préparée}. Applique le format Pourcentage avec 1 décimale. Étends.
    \item \textbf{Alerte Rupture (Colonne H - intitulé « Alerte Stock ») :} En H3, utilise la fonction \texttt{SI} : si \texttt{Qté servie = Qté préparée} alors \texttt{"Rupture stock"} sinon \texttt{"OK"}. Étends.
    \item \textbf{Catégorisation Affluence (Colonne I - intitulé « Niveau Affluence ») :} En I3, utilise \texttt{SI.MULTIPLE} (ou \texttt{SI} imbriqués) sur la \texttt{Qté servie} selon les seuils définis dans la situation. Étends.
    \item \textbf{Mise en Forme Conditionnelle (MFC) :}
    \begin{enumerate}
        \item Sur la plage \texttt{F3:F7} (CA Jour) : Barres de données (Dégradé Vert), afficher la barre seulement.
        \item Sur le tableau entier \texttt{A3:I7} : Nouvelle règle \texttt{Formule} : \texttt{=\$H3="Rupture stock"} $\rightarrow$ Format : Police Blanche, Fond Rouge, Gras. (Attention à la référence mixte \texttt{\$H3}).
    \end{enumerate}
    \item \textbf{Synthèse Hebdomadaire (Zone hors tableau, ex. à partir de A10) :} Présente proprement (étiquettes en colonne A, valeurs en B) :
    \begin{itemize}
        \item Total CA Semaine : \texttt{SOMME(F3:F7)}.
        \item Moyenne CA Jour : \texttt{MOYENNE(F3:F7)}.
        \item Max CA Jour : \texttt{MAX(F3:F7)}.
        \item Nb Jours Rupture : \texttt{NB.SI(H3:H7; "Rupture stock")}.
    \end{itemize}
    Applique le format FCFA aux montants.
    \item \textbf{Graphique :} Sélectionne \texttt{A3:A7} (Jours), \texttt{F3:F7} (CA) et \texttt{D3:D7} (Qté servie) (touche \texttt{Ctrl}). Insère un \textbf{Graphique en colonnes groupées}. Déplace la série « Qté servie » sur l'axe secondaire (Format série $\rightarrow$ Axe secondaire). Titre : « \textsc{Performance Quotidienne : CA vs Quantités Servies} ». Axe principal : « CA (FCFA) », Axe secondaire : « Quantités ». Légende en bas. Place le graphique sous la synthèse (vers A15).
    \item \textbf{Mise en page \& Impression :} Onglet \texttt{Mise en page} $\rightarrow$ Orientation \texttt{Paysage}. \texttt{Zone d'impression} $\rightarrow$ Sélectionne A1:I20 (tableau + synthèse + graphique) $\rightarrow$ Définir. \texttt{Titres d'impression} $\rightarrow$ Lignes à répéter en haut : \texttt{\$2:\$2}. \texttt{En-tête/Pied de page} $\rightarrow$ En-tête personnalisé : Centre « \textsc{Cantine Lycée Bilingue Nkolbisson - Semaine du 13/01/2025} ». Pied de page : Centre « Page 1/1 - Élève [Ton Prénom Nom] ». Aperçu avant impression $\rightarrow$ Vérifier que tout tient sur 1 page (Ajuster à 1 page de large $\times$ 1 page de haut si nécessaire).
    \item \textbf{Export Final :} Fichier $\rightarrow$ Exporter sous $\rightarrow$ PDF. Nom du fichier : \texttt{Rapport\_Cantine\_Semaine\_PrenomNom.pdf}. Vérifie l'ouverture du PDF.
\end{enumerate}

\courssub{Ressources à mobiliser}
\begin{propriete}[Syntaxe des fonctions clés]
\begin{itemize}
    \item \textbf{SI(condition; val\_si\_vrai; val\_si\_faux)} : Teste une condition unique.
    \item \textbf{SI.MULTIPLE(cond1; val1; cond2; val2; ... ; val\_par\_défaut)} : Teste une liste de conditions dans l'ordre ; renvoie la valeur de la 1ère vraie. Disponible sur LibreOffice Calc et Excel récent (sinon \texttt{SI} imbriqués : \texttt{SI(cond1; val1; SI(cond2; val2; ...))}).
    \item \textbf{SOMME(plage)}, \textbf{MOYENNE(plage)}, \textbf{MAX(plage)} : Agrégats standards.
    \item \textbf{NB.SI(plage; critère)} : Compte les cellules vérifiant le critère (ex: \texttt{"Rupture stock"}).
    \item \textbf{Référence absolue (\$)} : \texttt{\$A\$1} (ligne et colonne fixes), \texttt{A\$1} (ligne fixe), \texttt{\$A1} (colonne fixe). Indispensable pour la MFC sur ligne entière (\texttt{\$H3}) et pour les constantes.
    \item \textbf{Format personnalisé FCFA} : \texttt{\# \#\#0 " FCFA"} (espace insécable avant FCFA recommandé : \texttt{Chr(160)} ou Alt+0160).
    \item \textbf{MFC Formule} : La formule doit s'écrire pour la \textbf{première cellule de la sélection} (ici A3) avec des \texttt{\$} sur les colonnes fixes seulement.
\end{itemize}
\end{propriete}

\begin{popfigure}
\begin{tikzpicture}[
    node distance=1.2cm and 1.5cm,
    startstop/.style={rectangle, rounded corners, minimum width=3cm, minimum height=1cm, text centered, draw=popblue, fill=popblueL, font=\small},
    decision/.style={diamond, aspect=2, minimum width=3.5cm, minimum height=1.2cm, text centered, draw=poporange, fill=poporangeL, font=\small, inner sep=0},
    arrow/.style={->, >=Stealth, thick, draw=popdark},
    label/.style={font=\footnotesize, text=popdark}
]
\node (start) [startstop] {Début : Saisir Qté Servie};
\node (dec1) [decision, below of=start] {Qté Servie $\ge$ 100 ?};
\node (val1) [startstop, right of=dec1, xshift=2.5cm, fill=popgreenL, draw=popgreen] {\textbf{Très forte}};
\node (dec2) [decision, below of=dec1] {Qté Servie $\ge$ 70 ?};
\node (val2) [startstop, right of=dec2, xshift=2.5cm, fill=popgreenL, draw=popgreen] {\textbf{Forte}};
\node (dec3) [decision, below of=dec2] {Qté Servie $\ge$ 40 ?};
\node (val3) [startstop, right of=dec3, xshift=2.5cm, fill=popgreenL, draw=popgreen] {\textbf{Moyenne}};
\node (val4) [startstop, below of=dec3, fill=poppinkL, draw=poppink] {\textbf{Faible}};
\node (end) [startstop, below of=val4] {Fin : Afficher résultat};

\draw [arrow] (start) -- (dec1);
\draw [arrow] (dec1) -- node[label, pos=0.5, left] {Oui} (val1);
\draw [arrow] (dec1) -- node[label, pos=0.5, right] {Non} (dec2);
\draw [arrow] (dec2) -- node[label, pos=0.5, left] {Oui} (val2);
\draw [arrow] (dec2) -- node[label, pos=0.5, right] {Non} (dec3);
\draw [arrow] (dec3) -- node[label, pos=0.5, left] {Oui} (val3);
\draw [arrow] (dec3) -- node[label, pos=0.5, right] {Non} (val4);
\draw [arrow] (val1) |- (end);
\draw [arrow] (val2) |- (end);
\draw [arrow] (val3) |- (end);
\draw [arrow] (val4) -- (end);
\end{tikzpicture}
\legende{Organigramme de la logique décisionnelle \texttt{SI.MULTIPLE} pour la catégorisation de l'affluence (seuils décroissants).}
\end{popfigure}

\courssub{Méthode : Construire un tableau de bord automatisé}
\begin{methode}[Démarche type pour un reporting tableur]
\begin{enumerate}
    \item \textbf{Structurer} : Saisir les données brutes, nommer les colonnes explicitement, figer la ligne d'en-tête (Geler volets).
    \item \textbf{Calculer} : Écrire les formules de base (opérateurs +, -, *, /) en gérant les références (relatives pour variables, absolues pour constantes).
    \item \textbf{Modéliser la logique} : Traduire les règles de gestion en fonctions \texttt{SI} / \texttt{SI.MULTIPLE} (tester les seuils du plus grand au plus petit).
    \item \textbf{Visualiser} : Appliquer MFC (barres pour magnitude, formules pour alertes lignes) + Graphique adapté (axe secondaire si échelles différentes).
    \item \textbf{Synthétiser} : Agréger avec \texttt{SOMME}, \texttt{MOYENNE}, \texttt{MAX}, \texttt{NB.SI} dans une zone dédiée.
    \item \textbf{Finaliser} : Mise en page (Paysage, Zone impression, Titres répétés, En-têtes/Pieds), Aperçu, Export PDF.
\end{enumerate}
\end{methode}

\courssub{Démarche et corrigé commenté}
\begin{exoresolu}[Résolution guidée pas à pas]
\textbf{Étape 1 : Saisie et Structure (Consignes 1)}
Ouvre un classeur vierge. En A1, saisis le titre long. Sélectionne A1:E1 $\rightarrow$ Clic droit $\rightarrow$ \texttt{Fusionner et centrer}. Applique le style (Police 14, Gras, Blanc, Fond \texttt{popdark} ou Bleu foncé standard). En ligne 2, saisis les en-têtes exacts : \texttt{Date}, \texttt{Plat principal}, \texttt{Qté préparée}, \texttt{Qté servie}, \texttt{Prix unitaire (FCFA)}. Saisis les données du tableau (lignes 3 à 7).
\emph{Format FCFA :} Sélectionne E3:E7 $\rightarrow$ \texttt{Ctrl+1} (Format cellules) $\rightarrow$ onglet \texttt{Nombres} $\rightarrow$ \texttt{Personnalisé} $\rightarrow$ tape \texttt{\# \#\#0 " FCFA"} $\rightarrow$ OK.
\emph{Geler volets :} Clique en A3 (première cellule de données) $\rightarrow$ Menu \texttt{Affichage} $\rightarrow$ \texttt{Geler} $\rightarrow$ \texttt{Geler les lignes et colonnes} (ou \texttt{Geler la ligne du haut} selon version). La ligne 2 reste visible en défilant.

\textbf{Étape 2 : Calculs de base (Consignes 2, 3)}
En F2 : \texttt{CA Jour (FCFA)}. En F3 : \texttt{=D3*E3}. \emph{Analyse :} Le prix change chaque jour (colonne E), la quantité servie aussi (colonne D). On utilise donc des références \textbf{relatives} (D3, E3). Tire la poignée de recopie jusqu'en F7. Formate F3:F7 en FCFA (même format qu'E).
En G2 : \texttt{Taux Service (\%)}. En G3 : \texttt{=D3/C3}. Formate G3:G7 en \texttt{Pourcentage} (1 décimale). Tire vers le bas.
\begin{center}
\begin{tabular}{|c|c|c|c|c|c|c|}
\hline
\rowcolor{popblueL} \textbf{Date} & \textbf{Plat} & \textbf{Préparée} & \textbf{Servie} & \textbf{Prix} & \textbf{CA Jour} & \textbf{Taux} \\
\hline
Lundi & Ndolé & 120 & 110 & 500 & 55\,000 & 91,7\% \\
Mardi & Riz sauce & 100 & 100 & 450 & 45\,000 & 100,0\% \\
Mercredi & Poulet DG & 80 & 60 & 1\,000 & 60\,000 & 75,0\% \\
Jeudi & Spaghetti & 90 & 85 & 400 & 34\,000 & 94,4\% \\
Vendredi & Poisson & 110 & 95 & 600 & 57\,000 & 86,4\% \\
\hline
\end{tabular}
\end{center}

\textbf{Étape 3 : Logique décisionnelle (Consignes 4, 5)}
En H2 : \texttt{Alerte Stock}. En H3 :
\begin{lstlisting}
=SI(D3=C3; "Rupture stock"; "OK")
\end{lstlisting}
Tire vers le bas. \emph{Résultat :} Mardi affiche « Rupture stock » (100=100), les autres « OK ».

En I2 : \texttt{Niveau Affluence}. En I3 (avec \texttt{SI.MULTIPLE} - recommandé pour lisibilité) :
\begin{lstlisting}
=SI.MULTIPLE(D3>=100; "Très forte"; D3>=70; "Forte"; D3>=40; "Moyenne"; "Faible")
\end{lstlisting}
\emph{Alternative SI imbriqués (compatibilité ancienne) :}
\begin{lstlisting}
=SI(D3>=100; "Très forte"; SI(D3>=70; "Forte"; SI(D3>=40; "Moyenne"; "Faible")))
\end{lstlisting}
\emph{Attention :} L'ordre des tests est crucial (du plus grand seuil au plus petit). Tire vers le bas.
\emph{Résultats :} Lundi (110) $\rightarrow$ Très forte ; Mardi (100) $\rightarrow$ Très forte ; Mercredi (60) $\rightarrow$ Moyenne ; Jeudi (85) $\rightarrow$ Forte ; Vendredi (95) $\rightarrow$ Forte.

\textbf{Étape 4 : Mise en Forme Conditionnelle (Consigne 6)}
\begin{enumerate}
    \item \textbf{Barres de données (CA) :} Sélectionne F3:F7. Menu \texttt{Format} $\rightarrow$ \texttt{Conditionnel} $\rightarrow$ \texttt{Barres de données} (ou \texttt{Mise en forme conditionnelle} $\rightarrow$ \texttt{Barres de données}). Choisis un dégradé Vert. Coche « Afficher la barre seulement » si tu veux masquer le chiffre (ici on garde le chiffre lisible, donc décocher).
    \item \textbf{Ligne Rouge Rupture :} Sélectionne \textbf{tout le tableau de données A3:I7} (pas les en-têtes). Menu \texttt{Format} $\rightarrow$ \texttt{Conditionnel} $\rightarrow$ \texttt{Condition} $\rightarrow$ \texttt{La formule est} (ou \texttt{Nouvelle règle} $\rightarrow$ \texttt{Utiliser une formule}). Saisis : \texttt{=\$H3="Rupture stock"}. Clic \texttt{Nouveau style} (ou \texttt{Format}) $\rightarrow$ onglet \texttt{Effets de police} : Gras, Couleur Blanche ; onglet \texttt{Arrière-plan} : Rouge. OK. OK.
    \emph{Justification du \$H3 :} Le \$ fige la colonne H (où est l'alerte). Le 3 sans \$ suit la ligne courante (3, 4, 5...). Ainsi, pour la ligne 4, la formule évaluée est \texttt{=\$H4="Rupture stock"}.
\end{enumerate}

\textbf{Étape 5 : Synthèse Hebdomadaire (Consigne 7)}
En A10 : \texttt{Total CA Semaine (FCFA)}. En B10 : \texttt{=SOMME(F3:F7)} $\rightarrow$ \textbf{251\,000 FCFA}. Formate FCFA.
En A11 : \texttt{Moyenne CA Jour (FCFA)}. En B11 : \texttt{=MOYENNE(F3:F7)} $\rightarrow$ \textbf{50\,200 FCFA}. Formate FCFA.
En A12 : \texttt{Max CA Jour (FCFA)}. En B12 : \texttt{=MAX(F3:F7)} $\rightarrow$ \textbf{60\,000 FCFA} (Mercredi). Formate FCFA.
En A13 : \texttt{Nb Jours Rupture}. En B13 : \texttt{=NB.SI(H3:H7; "Rupture stock")} $\rightarrow$ \textbf{1} (Mardi seulement).

\textbf{Étape 6 : Graphique Combiné (Consigne 8)}
Sélectionne A3:A7 (Jours), puis \texttt{Ctrl} + sélectionne F3:F7 (CA), puis \texttt{Ctrl} + sélectionne D3:D7 (Qté servie).
Insertion $\rightarrow$ Graphique $\rightarrow$ \texttt{Colonnes groupées} (Clustered Column). Valider.
Le graphique a deux séries : « CA Jour » (valeurs ~50k) et « Qté servie » (valeurs ~100). L'échelle diffère trop $\rightarrow$ Clic droit sur la série « Qté servie » (colonnes oranges/bleues selon thème) $\rightarrow$ \texttt{Formater la série de données} $\rightarrow$ \texttt{Axe secondaire}. La série « Qté servie » bascule sur un axe à droite (échelle 0-120).
Personnalise : Titre du graphique $\rightarrow$ « \textsc{Performance Quotidienne : CA vs Quantités Servies} ». Titre axe principal (gauche) $\rightarrow$ « Chiffre d'Affaires (FCFA) ». Titre axe secondaire (droite) $\rightarrow$ « Quantités Servies ». Légende $\rightarrow$ Position Bas. Déplace/Redimensionne le graphique pour qu'il s'insère sous la synthèse (coin haut gauche en A15 environ).

\textbf{Étape 7 : Mise en page finale (Consigne 9)}
Onglet \texttt{Mise en page}.
\begin{itemize}
    \item \texttt{Orientation} $\rightarrow$ \texttt{Paysage}.
    \item \texttt{Zone d'impression} $\rightarrow$ Sélectionne A1:I25 (couvre tableau, synthèse, graphique) $\rightarrow$ \texttt{Définir la zone d'impression}.
    \item \texttt{Titres d'impression} $\rightarrow$ onglet \texttt{Feuille} $\rightarrow$ \texttt{Lignes à répéter en haut} $\rightarrow$ clique sur le n° de ligne 2 (ou tape \texttt{\$2:\$2}). Cela répète les en-têtes de colonnes si le tableau déborde (ici non, mais exigé).
    \item \texttt{En-tête/Pied de page} $\rightarrow$ \texttt{En-tête personnalisé} $\rightarrow$ Section Centre : \texttt{Cantine Lycée Bilingue Nkolbisson - Semaine du 13/01/2025}. \texttt{Pied de page personnalisé} $\rightarrow$ Section Centre : \texttt{Page 1/1 - Élève [Ton Prénom Nom]}.
    \item \texttt{Aperçu avant impression} : Vérifie que tout tient sur 1 page. Si graphique coupé : \texttt{Mise en page} $\rightarrow$ \texttt{Échelle} $\rightarrow$ \texttt{Ajuster à} 1 page(s) de large $\times$ 1 page(s) de haut.
\end{itemize}

\textbf{Étape 8 : Export PDF (Consigne 10)}
Fichier $\rightarrow$ Exporter au format $\rightarrow$ \texttt{Exportation directe en PDF} (ou \texttt{Enregistrer sous} $\rightarrow$ Type PDF). Nom : \texttt{Rapport\_Cantine\_Semaine\_PrenomNom.pdf}. Ouvre le PDF pour valider la mise en page, la lisibilité des barres de données, la couleur rouge sur la ligne Mardi, la présence des en-têtes/pieds de page.
\tcblower
\textbf{Solution finale commentée (Extrait des formules clés en F3, G3, H3, I3) :}
\begin{lstlisting}
F3: =D3*E3                    // CA = Qté Servie * Prix Unitaire (Relatif)
G3: =D3/C3                    // Taux = Servie / Préparée (Relatif) -> Format %
H3: =SI(D3=C3; "Rupture stock"; "OK")  // Test égalité
I3: =SI.MULTIPLE(D3>=100; "Très forte"; D3>=70; "Forte"; D3>=40; "Moyenne"; "Faible")
\end{lstlisting}
\textbf{Synthèse (valeurs attendues) :}
\begin{center}
\begin{tabular}{|l|r|}
\hline
\rowcolor{popblueL} \textbf{Indicateur} & \textbf{Valeur} \\
\hline
Total CA Semaine & 251\,000 FCFA \\
Moyenne CA Jour & 50\,200 FCFA \\
Max CA Jour & 60\,000 FCFA \\
Nb Jours Rupture & 1 \\
\hline
\end{tabular}
\end{center}
\textbf{Vérifications visuelles attendues :}
\begin{itemize}
    \item Ligne Mardi (ligne 4) : Fond rouge, police blanche, gras (MFC ligne entière).
    \item Colonne F : Barres vertes proportionnelles (Mercredi et Vendredi les plus longues).
    \item Graphique : Colonnes bleues (CA) lisibles sur axe gauche, colonnes oranges (Qté) sur axe droit, pas d'écrasement.
    \item PDF : 1 page, Paysage, En-tête/Pied de page présents, Titre tableau répété si multi-pages (non nécessaire ici mais configuré).
\end{itemize}
\end{exoresolu}

\courssub{Bilan de l'activité}
Cette activité a permis de valider l'acquisition des compétences transversales du tableur en situation de \textbf{gestion réelle}. Les points critiques maîtrisés sont :
\begin{enumerate}
    \item \textbf{La rigueur des références :} La distinction relative/absolue (\texttt{\$H3} en MFC) est la source n°1 d'erreurs ; sa maîtrise prouve la compréhension du modèle de calcul matriciel.
    \item \textbf{La modélisation logique :} Le passage du langage naturel (« Si stock épuisé alors alerte ») à la syntaxe formelle (\texttt{SI}, \texttt{SI.MULTIPLE}) structure la pensée algorithmique. L'ordre des seuils dans \texttt{SI.MULTIPLE} (décroissant) est un piège classique évité.
    \item \textbf{La lisibilité professionnelle :} La MFC (barres + coloration ligne) transforme une grille de chiffres en tableau de bord décisionnel (« \textit{Data Visualization} » légère). Le format FCFA personnalisé ancre le document dans le contexte monétaire local.
    \item \textbf{La chaîne de production documentaire :} De la saisie brute à l'export PDF normé (mise en page, répétition titres, en-têtes/pieds), l'élève a produit un \textbf{livrable final} conforme aux exigences administratives (archivage, transmission hiérarchique).
\end{enumerate}
\begin{attention}[Pièges à éviter pour le BEPC]
\begin{itemize}
    \item Oublier le \texttt{\$} dans la formule MFC (\texttt{=H3="..."} au lieu de \texttt{=\$H3="..."}) $\rightarrow$ seule la colonne H se colorie, pas la ligne entière.
    \item Confondre \texttt{SI.MULTIPLE} (conditions successives) et \texttt{ET/OU} (conditions simultanées). Ici, les seuils sont exclusifs et ordonnés $\rightarrow$ \texttt{SI.MULTIPLE} ou \texttt{SI} imbriqués.
    \item Ne pas formater la colonne Taux en \% $\rightarrow$ affichage « 0,917 » au lieu de « 91,7\% ».
    \item Oublier de mettre la série « Quantités » sur l'axe secondaire $\rightarrow$ graphique illisible (colonnes CA écrasées par l'échelle des quantités).
    \item Ne pas définir la zone d'impression $\rightarrow$ impression de pages blanches ou tronquées.
\end{itemize}
\end{attention}
\begin{savaistu}[Le numérique au service de la cantine scolaire au Cameroun]
Au-delà de l'exercice scolaire, de nombreux lycées camerounais (ex: Lycée Leclerc, Lycée Bilingue d'Essos) déploient des solutions de gestion de cantine intégrées (badgeuses, cartes à puce, logiciels de stock). Le tableur reste l'outil « \textit{low-tech} » universel pour le prototypage rapide, l'audit de données ou le reporting hors-ligne. Savoir produire un \texttt{Rapport\_Cantine\_Semaine\_...pdf} structuré, signé (en-tête/pied de page) et infalsifiable (valeurs figées à l'export) est une compétence recherchée pour les stages en intendance ou en secrétariat de direction. Le format PDF/A (archivage) serait l'étape suivante pour garantir la lisibilité à 10 ans.
\end{savaistu}

\newpage

\coursec{Méthodes et exercices résolus}

\courssub{Méthodes et exercices résolus}

\begin{methode}[Structurer et mettre en forme un classeur de gestion]
\textbf{Objectif :} créer un tableau de suivi des ventes d'une boutique (ex. : \og{}Boutique École\fg{} au quartier Mvog-Ada) lisible, structuré et prêt au calcul.
\begin{enumerate}
    \item \textbf{Planifier les en-têtes :} en ligne 1, saisir les titres de colonnes : \texttt{A1 : Date}, \texttt{B1 : Article}, \texttt{C1 : Quantité}, \texttt{D1 : Prix Unitaire (FCFA)}, \texttt{E1 : Total (FCFA)}.
    \item \textbf{Mise en forme des en-têtes :} sélectionner la plage \texttt{A1:E1} $\rightarrow$ \textbf{Gras} (\texttt{Ctrl+G} sous Excel, \texttt{Ctrl+B} sous LibreOffice), \textbf{Centrer}, \textbf{Couleur de remplissage} (ex. bleu foncé), \textbf{Couleur de police blanche}, \textbf{Toutes les bordures}.
    \item \textbf{Format des nombres :} sélectionner la colonne D $\rightarrow$ clic droit \og{}Format de cellule\fg{} $\rightarrow$ catégorie \textbf{Nombre} $\rightarrow$ 0 décimale $\rightarrow$ cocher \og{}Séparateur de milliers\fg{}. Faire de même pour la colonne E.
    \item \textbf{Ajustement des largeurs :} double-cliquer sur la séparation entre deux en-têtes de colonnes (A/B, B/C, etc.) pour \textbf{ajuster automatiquement la largeur} au contenu.
    \item \textbf{Figer la ligne d'en-tête :} onglet \textbf{Affichage} $\rightarrow$ \textbf{Figer les volets} $\rightarrow$ \og{}Figer la première ligne\fg{} : les titres restent visibles quand on fait défiler la feuille vers le bas.
\end{enumerate}
\end{methode}

\begin{exoresolu}[Création du classeur \og{}Ventes Juin 2025\fg{}]
\textbf{Énoncé :} vous êtes chargé de tenir le journal de vente d'une papeterie scolaire. Créez le classeur permettant de saisir les 5 premières ventes du mois. Appliquez une mise en forme professionnelle (en-têtes stylisés, format monétaire FCFA, gel de la ligne 1). Donnez la procédure exacte pour figer la ligne d'en-tête sous LibreOffice Calc et sous Microsoft Excel.
\tcblower
\textbf{Solution détaillée :}
\begin{enumerate}
    \item \textbf{Saisie des en-têtes (ligne 1) :}
    \begin{lstlisting}
A1 : Date        B1 : Article          C1 : Quantite
D1 : Prix Unitaire (FCFA)   E1 : Total (FCFA)
    \end{lstlisting}
    \item \textbf{Mise en forme des en-têtes :} sélection \texttt{A1:E1}. Onglet Accueil $\rightarrow$ Gras, Centrer, couleur de remplissage bleu foncé, couleur de police blanche, bordures \og{}Toutes les bordures\fg{}.
    \item \textbf{Format monétaire (colonnes D et E) :} sélectionner la colonne D puis, en maintenant \texttt{Ctrl}, la colonne E. Clic droit \og{}Format de cellule\fg{} $\rightarrow$ onglet \textbf{Nombres} $\rightarrow$ catégorie \textbf{Nombre}, 0 décimale, séparateur de milliers (ou format personnalisé \texttt{\# \#\#0 " FCFA"}). \textbf{Valider}.
    \item \textbf{Figer la ligne 1 :}
    \begin{itemize}
        \item \textbf{Excel :} onglet \textbf{Affichage} $\rightarrow$ groupe \textbf{Fenêtre} $\rightarrow$ \textbf{Figer les volets} $\rightarrow$ \textbf{Figer la première ligne}.
        \item \textbf{LibreOffice Calc :} placer le curseur en \texttt{A2}, puis menu \textbf{Affichage} $\rightarrow$ \textbf{Figer les lignes et les colonnes} (tout ce qui est au-dessus et à gauche de la cellule active est figé).
    \end{itemize}
    \item \textbf{Saisie des données (exemple) :}
    \begin{center}
    \begin{tabularx}{\linewidth}{|l|X|c|r|r|}\hline
    \rowcolor{popblueL} \textbf{Date} & \textbf{Article} & \textbf{Qté} & \textbf{Prix Unit.} & \textbf{Total} \\\hline
    01/06/2025 & Stylo Bic & 50 & 100 & 5\,000 \\\hline
    01/06/2025 & Cahier 100p & 20 & 500 & 10\,000 \\\hline
    02/06/2025 & Gomme & 30 & 50 & 1\,500 \\\hline
    03/06/2025 & Règle 30cm & 15 & 200 & 3\,000 \\\hline
    03/06/2025 & Crayon HB & 100 & 50 & 5\,000 \\\hline
    \end{tabularx}
    \end{center}
    \item \textbf{Vérification :} faire défiler vers le bas (flèche bas ou molette) : la ligne 1 reste visible. La mise en page est propre, les nombres sont alignés à droite avec séparateur de milliers.
\end{enumerate}
\textbf{Conclusion :} le classeur est opérationnel pour la saisie journalière. La ligne d'en-tête figée évite les erreurs de colonne sur de longues listes.
\end{exoresolu}

\begin{methode}[Construire une formule de calcul avec références relatives et absolues]
\textbf{Objectif :} automatiser le calcul \texttt{Total = Quantité $\times$ Prix Unitaire} et le recopier vers le bas sans erreur, puis calculer une TVA avec un taux fixe placé dans une cellule unique.
\begin{enumerate}
    \item \textbf{Formule simple (références relatives) :} en \texttt{E2}, taper \texttt{=C2*D2}. Valider (\texttt{Entrée}). Les références \texttt{C2} et \texttt{D2} sont \textbf{relatives} : en tirant la poignée de recopie vers le bas, elles deviennent \texttt{C3*D3}, \texttt{C4*D4}, etc.
    \item \textbf{Recopie :} cliquer sur \texttt{E2}, placer le curseur sur la \textbf{poignée de recopie} (petit carré en bas à droite de la cellule active), tirer vers le bas jusqu'à \texttt{E6}.
    \item \textbf{Référence absolue (taux de TVA) :} saisir le taux en \texttt{H1} : \texttt{19,25\%}. En \texttt{F2}, pour la TVA : \texttt{=E2*\$H\$1}. Le symbole \texttt{\$} devant la lettre de colonne et devant le numéro de ligne \textbf{bloque} la référence lors de la recopie.
    \item \textbf{Test :} tirer \texttt{F2} vers le bas. Vérifier en \texttt{F3} : la formule est \texttt{=E3*\$H\$1} (\texttt{E2} est devenu \texttt{E3}, mais \texttt{\$H\$1} est resté fixe).
    \item \textbf{Référence mixte (pour information) :} \texttt{\$H1} (colonne fixe, ligne variable) ou \texttt{H\$1} (ligne fixe, colonne variable).
\end{enumerate}
\end{methode}

\begin{exoresolu}[Calcul des totaux et de la TVA sur le journal de vente]
\textbf{Énoncé :} reprenez le classeur précédent.
\begin{enumerate}
    \item En \texttt{E2}, écrivez la formule pour calculer le total de la première vente. Recopiez jusqu'en \texttt{E6}.
    \item En \texttt{H1}, saisissez le taux de TVA \texttt{19,25\%}. En \texttt{F1}, écrivez \og{}TVA\fg{}. En \texttt{F2}, écrivez la formule pour calculer la TVA de la première vente en verrouillant la référence au taux. Recopiez jusqu'en \texttt{F6}.
    \item En \texttt{G1}, écrivez \og{}Total TTC\fg{}. En \texttt{G2}, écrivez la formule \texttt{Total + TVA}. Recopiez.
    \item Que se passe-t-il si l'on modifie la valeur en \texttt{H1} en \texttt{20\%} ?
\end{enumerate}
\tcblower
\textbf{Solution détaillée :}
\begin{enumerate}
    \item \textbf{Total HT (colonne E) :}
    \begin{itemize}
        \item \texttt{E2} : \texttt{=C2*D2} $\rightarrow$ résultat : \texttt{5000}.
        \item Recopie \texttt{E2:E6} : \texttt{E3 = 10000}, \texttt{E4 = 1500}, \texttt{E5 = 3000}, \texttt{E6 = 5000}.
    \end{itemize}
    \item \textbf{TVA (colonne F) :}
    \begin{itemize}
        \item \texttt{H1} : \texttt{19,25\%} (le tableur stocke la valeur $0,1925$ et l'affiche en pourcentage).
        \item \texttt{F2} : \texttt{=E2*\$H\$1}. \textbf{Explication :} \texttt{E2} est relative (change à chaque ligne), \texttt{\$H\$1} est absolue (reste \texttt{H1} partout).
        \item Recopie \texttt{F2:F6} :
        \begin{align*}
        \texttt{F2} &= 5000 \times 0,1925 = 962,5 \\
        \texttt{F3} &= 10000 \times 0,1925 = 1925 \\
        \texttt{F4} &= 1500 \times 0,1925 = 288,75 \\
        \texttt{F5} &= 3000 \times 0,1925 = 577,5 \\
        \texttt{F6} &= 5000 \times 0,1925 = 962,5
        \end{align*}
    \end{itemize}
    \item \textbf{Total TTC (colonne G) :}
    \begin{itemize}
        \item \texttt{G2} : \texttt{=E2+F2} (ou directement \texttt{=E2*(1+\$H\$1)}). Recopier vers le bas.
    \end{itemize}
    \item \textbf{Test de modification du taux :} si \texttt{H1} devient \texttt{20\%}, \textbf{toutes} les valeurs des colonnes F et G se recalculent instantanément. C'est l'intérêt de la référence absolue : \textbf{un seul point de modification} pour tout le classeur.
\end{enumerate}
\textbf{Conclusion :} l'usage de \texttt{\$} sur le taux (\texttt{\$H\$1}) garantit l'intégrité du calcul lors de la recopie. Sans \texttt{\$}, la référence glisserait vers \texttt{H2}, \texttt{H3}... (cellules vides), donnant des résultats nuls ou erronés.
\end{exoresolu}

\begin{methode}[Calculer des indicateurs statistiques avec les fonctions usuelles]
\textbf{Objectif :} résumer un jeu de données (notes d'élèves, stock, ventes) avec \texttt{SOMME}, \texttt{MOYENNE}, \texttt{NB}, \texttt{NBVAL}, \texttt{MIN}, \texttt{MAX}, \texttt{GRANDE.VALEUR}, \texttt{PETITE.VALEUR}.
\begin{enumerate}
    \item \textbf{Zone de synthèse :} réserver une zone en bas du tableau (ex. lignes 10 à 16) ou sur une feuille \og{}Bilan\fg{}.
    \item \textbf{Syntaxe générale :} \texttt{=NOM\_FONCTION(plage\_de\_cellules)}.
    \item \textbf{Fonctions clés :}
    \begin{itemize}
        \item \texttt{=SOMME(E2:E6)} : total des ventes HT.
        \item \texttt{=MOYENNE(E2:E6)} : panier moyen.
        \item \texttt{=NB(E2:E6)} : nombre de cellules \textbf{numériques} (ignore le texte et les vides).
        \item \texttt{=NBVAL(E2:E6)} : nombre de cellules \textbf{non vides} (texte + nombres).
        \item \texttt{=MIN(E2:E6)} / \texttt{=MAX(E2:E6)} : vente minimale / maximale.
        \item \texttt{=GRANDE.VALEUR(E2:E6;2)} : 2\up{e} plus grosse vente.
        \item \texttt{=PETITE.VALEUR(E2:E6;1)} : plus petite vente (équivaut ici à \texttt{MIN}).
    \end{itemize}
    \item \textbf{Insertion assistée :} onglet \textbf{Formules} $\rightarrow$ \textbf{Insérer une fonction} (bouton \texttt{fx}) $\rightarrow$ catégorie \og{}Statistiques\fg{} $\rightarrow$ choisir la fonction $\rightarrow$ sélectionner la plage à la souris $\rightarrow$ OK.
\end{enumerate}
\end{methode}

\begin{exoresolu}[Tableau de bord statistique de la papeterie]
\textbf{Énoncé :} à partir des données de l'exercice précédent (colonne E : Total HT), construisez un tableau de synthèse en \texttt{A10:B16} avec les indicateurs suivants. Donnez la formule exacte à saisir en \texttt{B10} à \texttt{B16}.
\begin{center}
\begin{tabularx}{\linewidth}{|l|X|}\hline
\rowcolor{popblueL} \textbf{Cellule} & \textbf{Indicateur} \\\hline
B10 & Total Chiffre d'Affaires HT \\\hline
B11 & Nombre de transactions \\\hline
B12 & Panier moyen HT \\\hline
B13 & Vente la plus élevée \\\hline
B14 & Vente la plus faible \\\hline
B15 & 2\up{e} meilleure vente \\\hline
B16 & Écart Max $-$ Min \\\hline
\end{tabularx}
\end{center}
\tcblower
\textbf{Solution détaillée :}
\begin{enumerate}
    \item \textbf{Libellés (colonne A) :} saisir les textes en \texttt{A10:A16}.
    \item \textbf{Formules (colonne B) :}
    \begin{lstlisting}
B10 : =SOMME(E2:E6)            -> 24 500
B11 : =NB(E2:E6)               -> 5 (5 nombres dans la plage)
B12 : =MOYENNE(E2:E6)          -> 4 900 (ou =B10/B11)
B13 : =MAX(E2:E6)              -> 10 000
B14 : =MIN(E2:E6)              -> 1 500
B15 : =GRANDE.VALEUR(E2:E6;2)  -> 5 000
B16 : =B13-B14                 -> 8 500
    \end{lstlisting}
    \item \textbf{Vérification manuelle :}
    \begin{align*}
    \text{Somme} &= 5000+10000+1500+3000+5000 = 24\,500 \quad \checkmark \\
    \text{Moyenne} &= 24500 / 5 = 4\,900 \quad \checkmark \\
    \text{Max} &= 10\,000 \quad ; \quad \text{Min} = 1\,500 \quad \checkmark
    \end{align*}
    Pour \texttt{B15} : valeurs triées par ordre décroissant : $10000$, $5000$, $5000$, $3000$, $1500$. La 2\up{e} plus grande valeur est bien $5000$ (la fonction tient compte des doublons). $\checkmark$
    \item \textbf{Mise en forme :} appliquer le format \og{}Nombre avec séparateur de milliers\fg{} (ou monétaire FCFA) sur \texttt{B10}, \texttt{B12}, \texttt{B13}, \texttt{B14}, \texttt{B15}, \texttt{B16}. Format \og{}Nombre\fg{} (0 décimale) sur \texttt{B11}.
\end{enumerate}
\textbf{Conclusion :} ces quelques fonctions permettent de produire un rapport de synthèse complet en quelques secondes, sans recopier de formules complexes.
\end{exoresolu}

\begin{methode}[Automatiser une décision avec la fonction SI imbriquée]
\textbf{Objectif :} classer automatiquement les ventes (\og{}Forte\fg{}, \og{}Moyenne\fg{}, \og{}Faible\fg{}) ou déterminer une décision (\og{}Admis\fg{}, \og{}Ajourné\fg{}) selon des seuils.
\begin{enumerate}
    \item \textbf{Syntaxe de base :} \texttt{=SI(test\_logique; valeur\_si\_vrai; valeur\_si\_faux)}. Le séparateur d'arguments est le \textbf{point-virgule (;)} en version française.
    \item \textbf{Imbrication (SI dans SI) :} pour 3 cas ou plus, on place un nouveau \texttt{SI} dans l'argument \og{}valeur\_si\_faux\fg{} du précédent.
    \item \textbf{Structure type à 3 cas :}
    \texttt{=SI(condition\_1; resultat\_1; SI(condition\_2; resultat\_2; resultat\_3))}
    \begin{itemize}
        \item si \texttt{condition\_1} est vraie $\rightarrow$ \texttt{resultat\_1} ;
        \item sinon, si \texttt{condition\_2} est vraie $\rightarrow$ \texttt{resultat\_2} ;
        \item sinon (toutes fausses) $\rightarrow$ \texttt{resultat\_3}.
    \end{itemize}
    \item \textbf{Opérateurs de comparaison :} \texttt{>}, \texttt{<}, \texttt{>=}, \texttt{<=}, \texttt{=}, \texttt{<>} (différent de).
    \item \textbf{Combinaison avec ET / OU :} \texttt{=SI(ET(cond1;cond2); ...)} (toutes les conditions vraies) ; \texttt{=SI(OU(cond1;cond2); ...)} (au moins une condition vraie). La fonction \texttt{NON(condition)} inverse un test logique.
    \item \textbf{Astuce lisibilité :} pour les formules longues, construire d'abord la logique au brouillon (arbre de décision), puis saisir la formule dans la barre de formule.
\end{enumerate}
\end{methode}

\begin{exoresolu}[Classification des ventes et décision sur les notes]
\textbf{Énoncé :}
\textbf{Partie A (Ventes) :} en colonne H (libellé \og{}Catégorie\fg{} en \texttt{H1}), classer chaque vente (colonne E) :
\begin{itemize}
    \item \texttt{>= 8000} $\rightarrow$ \og{}Forte\fg{} ;
    \item \texttt{>= 3000} $\rightarrow$ \og{}Moyenne\fg{} ;
    \item sinon $\rightarrow$ \og{}Faible\fg{}.
\end{itemize}
\textbf{Partie B (Notes) :} sur une nouvelle feuille \og{}Notes\fg{}, colonne A : Noms, colonne B : Note/20. En \texttt{C1}, écrire \og{}Décision\fg{}. Règle :
\begin{itemize}
    \item note $\ge 10$ $\rightarrow$ \og{}Admis\fg{} ;
    \item note $\ge 8$ $\rightarrow$ \og{}Repêchage\fg{} ;
    \item sinon $\rightarrow$ \og{}Ajourné\fg{}.
\end{itemize}
Donnez les formules exactes en \texttt{H2} et \texttt{C2}.
\tcblower
\textbf{Solution détaillée :}

\textbf{Partie A : classification des ventes}
\begin{itemize}
    \item \textbf{Logique :} on teste d'abord le seuil le plus haut (\texttt{>=8000}), puis le suivant (\texttt{>=3000}). L'ordre est crucial : si l'on testait \texttt{>=3000} en premier, la valeur \texttt{10000} serait classée \og{}Moyenne\fg{}.
    \item \textbf{Formule en H2 :}
    \begin{lstlisting}
=SI(E2>=8000; "Forte"; SI(E2>=3000; "Moyenne"; "Faible"))
    \end{lstlisting}
    \item \textbf{Résultats attendus :}
    \begin{align*}
    \texttt{E2}=5000 &\rightarrow 5000 \ge 8000\ ?\ \text{FAUX} \rightarrow 5000 \ge 3000\ ?\ \text{VRAI} \rightarrow \text{\og{}Moyenne\fg{}} \\
    \texttt{E3}=10000 &\rightarrow \text{VRAI} \rightarrow \text{\og{}Forte\fg{}} \\
    \texttt{E4}=1500 &\rightarrow \text{FAUX, FAUX} \rightarrow \text{\og{}Faible\fg{}} \\
    \texttt{E5}=3000 &\rightarrow \text{FAUX, VRAI (car } 3000 \ge 3000\text{)} \rightarrow \text{\og{}Moyenne\fg{}} \\
    \texttt{E6}=5000 &\rightarrow \text{\og{}Moyenne\fg{}}
    \end{align*}
\end{itemize}
\textbf{Partie B : décision sur les notes}
\begin{itemize}
    \item \textbf{Formule en C2 :}
    \begin{lstlisting}
=SI(B2>=10; "Admis"; SI(B2>=8; "Repêchage"; "Ajourné"))
    \end{lstlisting}
    \item \textbf{Test avec \texttt{B2}=12, \texttt{B3}=9, \texttt{B4}=7 :} \texttt{C2} = \og{}Admis\fg{}, \texttt{C3} = \og{}Repêchage\fg{}, \texttt{C4} = \og{}Ajourné\fg{}.
\end{itemize}
\textbf{Conclusion :} l'imbrication \texttt{SI(SI(...))} gère les seuils ordonnés. Au BEPC, on exige de savoir écrire un \texttt{SI} simple et un \texttt{SI} imbriqué à deux niveaux, avec les bons opérateurs et le point-virgule comme séparateur.
\end{exoresolu}

\begin{methode}[Analyser visuellement les données : mise en forme conditionnelle (MFC) et graphique]
\textbf{Objectif :} repérer instantanément les valeurs remarquables (MFC) et présenter les ventes sous forme de graphique.
\begin{enumerate}
    \item \textbf{MFC, règle \og{}Supérieur à\fg{} :}
    \begin{enumerate}
        \item sélectionner la plage cible (ex. \texttt{E2:E6} pour les totaux HT) ;
        \item onglet Accueil $\rightarrow$ \textbf{Mise en forme conditionnelle} $\rightarrow$ \textbf{Règles de mise en surbrillance des cellules} $\rightarrow$ \textbf{Supérieur à...} ;
        \item saisir la valeur (ex. \texttt{5000}) $\rightarrow$ choisir un format (ex. remplissage vert, police blanche) $\rightarrow$ OK.
    \end{enumerate}
    \item \textbf{MFC, \og{}Barres de données\fg{} (visualisation dans la cellule) :} sélection \texttt{E2:E6} $\rightarrow$ Mise en forme conditionnelle $\rightarrow$ \textbf{Barres de données} $\rightarrow$ choisir une couleur. \textbf{Effet :} une barre proportionnelle à la valeur s'affiche dans chaque cellule.
    \item \textbf{Création d'un graphique (histogramme ou secteurs) :}
    \begin{enumerate}
        \item sélectionner les données \textbf{avec les en-têtes} : \texttt{B1:B6} (articles) puis, en maintenant \texttt{Ctrl}, \texttt{E1:E6} (totaux) ;
        \item onglet \textbf{Insertion} $\rightarrow$ groupe \textbf{Graphiques} $\rightarrow$ choisir \textbf{Histogramme groupé} (comparaison de valeurs) ou \textbf{Secteurs} (répartition en parts) ;
        \item \textbf{personnalisation indispensable :}
        \begin{itemize}
            \item \textbf{titre du graphique :} ex. \og{}Chiffre d'affaires par article -- Juin 2025\fg{} ;
            \item \textbf{titre des axes} (histogramme) : axe vertical \og{}Montant (FCFA)\fg{} ;
            \item \textbf{légende :} la supprimer s'il n'y a qu'une seule série de données ;
            \item \textbf{étiquettes de données :} les activer pour afficher les valeurs sur le graphique.
        \end{itemize}
    \end{enumerate}
\end{enumerate}
\end{methode}

\begin{exoresolu}[Mise en évidence des meilleures ventes et graphique de synthèse]
\textbf{Énoncé :} sur le classeur \og{}Ventes Juin 2025\fg{} :
\begin{enumerate}
    \item Appliquez une MFC \og{}Barres de données\fg{} bleue sur la colonne \texttt{Total (E2:E6)}.
    \item Appliquez une MFC \og{}Supérieur à 5000\fg{} sur la même plage : remplissage vert foncé, police blanche, gras.
    \item Créez un graphique en \textbf{secteurs (camembert)} représentant la part de chaque article dans le chiffre d'affaires total. Donnez le titre exact et décrivez les 3 étapes de personnalisation indispensables.
\end{enumerate}
\tcblower
\textbf{Solution détaillée :}
\begin{enumerate}
    \item \textbf{MFC Barres de données (E2:E6) :} sélection \texttt{E2:E6} $\rightarrow$ Accueil $\rightarrow$ Mise en forme conditionnelle $\rightarrow$ Barres de données $\rightarrow$ remplissage bleu.
    \textbf{Résultat visuel :} la cellule \texttt{E3} (10000) a une barre pleine ; \texttt{E2} et \texttt{E6} (5000) une barre à moitié remplie ; \texttt{E4} (1500) une barre très fine. La comparaison est immédiate.
    \item \textbf{MFC \og{}Supérieur à 5000\fg{} (E2:E6) :} sélection \texttt{E2:E6} $\rightarrow$ Mise en forme conditionnelle $\rightarrow$ Règles de mise en surbrillance $\rightarrow$ Supérieur à $\rightarrow$ saisir \texttt{5000} $\rightarrow$ format personnalisé : remplissage vert foncé, police blanche, gras $\rightarrow$ OK.
    \textbf{Résultat :} seule la cellule \texttt{E3} (10000) se colore en vert. \texttt{E2} et \texttt{E6} (égales à 5000) ne sont \textbf{pas} colorées, car la condition est \textbf{strictement supérieure}. \textbf{Piège :} pour inclure la borne, choisir \og{}Supérieur ou égal à\fg{} (via \og{}Autres règles\fg{}).
    \item \textbf{Graphique en secteurs (répartition) :}
    \begin{enumerate}
        \item \textbf{Sélection :} \texttt{B1:B6} (noms des articles) puis, \texttt{Ctrl} maintenu, \texttt{E1:E6} (totaux).
        \item \textbf{Insertion :} onglet Insertion $\rightarrow$ Graphiques $\rightarrow$ Secteurs $\rightarrow$ \textbf{Secteur 2D} (premier icône).
        \item \textbf{Personnalisation (3 étapes clés) :}
        \begin{enumerate}
            \item \textbf{Titre :} cliquer sur le titre du graphique $\rightarrow$ saisir : \textbf{\og{}Répartition du chiffre d'affaires par article -- Juin 2025\fg{}}.
            \item \textbf{Étiquettes de données :} cliquer sur le graphique $\rightarrow$ bouton \texttt{+} (Éléments du graphique) $\rightarrow$ cocher \textbf{Étiquettes de données} $\rightarrow$ ouvrir les options et cocher \textbf{Pourcentage} et \textbf{Nom de la catégorie}.
            \item \textbf{Légende :} décocher \textbf{Légende} dans le bouton \texttt{+}, puisque les étiquettes portent déjà les noms des catégories.
        \end{enumerate}
    \end{enumerate}
    \textbf{Analyse :} le secteur \og{}Cahier 100p\fg{} ($10000/24500 \approx 41\%$) domine visuellement. \og{}Stylo Bic\fg{} et \og{}Crayon HB\fg{} (5000 chacun, soit $\approx 20\%$ chacun) ont des parts égales.
\end{enumerate}
\textbf{Conclusion :} la MFC permet une analyse rapide directement dans le tableau ; le graphique sert à communiquer les résultats (présentation orale, rapport imprimé).
\end{exoresolu}

\begin{methode}[Préparer la mise en page pour l'impression d'un rapport professionnel]
\textbf{Objectif :} imprimer le tableau de ventes et le graphique sur une seule page A4, lisible, avec en-tête et pied de page identifiant le document.
\begin{enumerate}
    \item \textbf{Onglet Mise en page :}
    \begin{itemize}
        \item \textbf{Orientation :} \textbf{Paysage} (largeur supérieure à la hauteur, idéal pour les tableaux larges) ;
        \item \textbf{Taille du papier :} \textbf{A4} ;
        \item \textbf{Marges :} \textbf{Étroites} (ou Normales) pour maximiser l'espace utile.
    \end{itemize}
    \item \textbf{Mise à l'échelle et titres répétés :}
    \begin{itemize}
        \item \textbf{Ajustement :} \og{}Ajuster la feuille sur une page\fg{} (force la réduction pour que tout tienne) ;
        \item \textbf{Lignes à répéter en haut :} onglet Mise en page $\rightarrow$ \textbf{Imprimer les titres} $\rightarrow$ champ \og{}Lignes à répéter en haut\fg{} $\rightarrow$ sélectionner la ligne \texttt{\$1:\$1} : les en-têtes apparaîtront sur chaque page si le tableau dépasse une page.
    \end{itemize}
    \item \textbf{En-tête / pied de page (boîte \og{}Mise en page\fg{}, onglet En-tête/Pied de page $\rightarrow$ Personnaliser) :}
    \begin{itemize}
        \item \textbf{en-tête centre :} \og{}Boutique École -- Journal des Ventes\fg{} ;
        \item \textbf{pied de page centre :} \og{}Page \&[Page] sur \&[Pages]\fg{} (pagination automatique) ;
        \item \textbf{pied de page droit :} \og{}Imprimé le \&[Date]\fg{}.
    \end{itemize}
    \item \textbf{Aperçu avant impression :} \texttt{Ctrl+P} (ou Fichier $\rightarrow$ Imprimer). Vérifier : rien n'est coupé, le graphique est visible, les marges sont correctes.
    \item \textbf{Zone d'impression (optionnel) :} sélectionner la zone utile (tableau + graphique) $\rightarrow$ Mise en page $\rightarrow$ \textbf{Zone d'impression} $\rightarrow$ \textbf{Définir}.
\end{enumerate}
\end{methode}

\begin{exoresolu}[Finalisation et impression du rapport mensuel]
\textbf{Énoncé :} vous devez remettre au gérant une version papier du suivi de juin.
\begin{enumerate}
    \item Configurez la page en \textbf{Paysage}, marges \textbf{Étroites}, ajustement sur \textbf{1 page}.
    \item Programmez la \textbf{répétition de la ligne 1} (en-têtes) en haut de chaque page imprimée.
    \item Définissez un \textbf{en-tête} : \og{}RAPPORT MENSUEL -- JUIN 2025\fg{} (centré, gras, taille 14).
    \item Définissez un \textbf{pied de page} : \og{}Page 1 sur 1 -- Confidentiel\fg{} (centré, numérotation automatique).
    \item Dans l'aperçu avant impression, vous constatez que le graphique est coupé en bas de page. Quelle action corrige cela sans changer l'orientation ?
\end{enumerate}
\tcblower
\textbf{Solution détaillée :}
\begin{enumerate}
    \item \textbf{Configuration de la page :} onglet Mise en page $\rightarrow$ Orientation \textbf{Paysage} ; Marges $\rightarrow$ \textbf{Étroites} ; groupe \textbf{Mise à l'échelle} $\rightarrow$ Largeur : \textbf{1 page}, Hauteur : \textbf{1 page} (ou cocher \og{}Ajuster la feuille sur une page\fg{} dans la boîte Mise en page).
    \item \textbf{Répétition des titres :} Mise en page $\rightarrow$ \textbf{Imprimer les titres} $\rightarrow$ champ \textbf{Lignes à répéter en haut} $\rightarrow$ cliquer sur la ligne 1 de la feuille (la référence \texttt{\$1:\$1} s'inscrit automatiquement) $\rightarrow$ OK.
    \item \textbf{En-tête :} dans la boîte \og{}Mise en page\fg{}, onglet \textbf{En-tête/Pied de page} $\rightarrow$ \textbf{Personnaliser l'en-tête...} $\rightarrow$ section \textbf{Centre} : taper \texttt{RAPPORT MENSUEL - JUIN 2025}, sélectionner le texte, cliquer sur l'icône \textbf{A} (police) $\rightarrow$ Gras, taille 14 $\rightarrow$ OK.
    \item \textbf{Pied de page :} \textbf{Personnaliser le pied de page...} $\rightarrow$ section \textbf{Centre} : taper \texttt{Page }, cliquer sur l'icône d'insertion du \textbf{numéro de page}, taper \texttt{ sur }, cliquer sur l'icône d'insertion du \textbf{nombre total de pages}, taper \texttt{ - Confidentiel} $\rightarrow$ OK $\rightarrow$ OK. Le code généré est \texttt{Page \&[Page] sur \&[Pages] - Confidentiel}.
    \item \textbf{Correction du graphique coupé :} deux solutions sans changer l'orientation :
    \begin{enumerate}
        \item \textbf{redimensionner le graphique :} cliquer sur le graphique $\rightarrow$ tirer une poignée d'angle vers l'intérieur pour le réduire légèrement afin qu'il entre dans la zone imprimable ;
        \item \textbf{vérifier les sauts de page :} onglet Affichage $\rightarrow$ \textbf{Aperçu des sauts de page} $\rightarrow$ glisser la ligne pointillée bleue pour inclure le graphique dans la page.
    \end{enumerate}
    \textbf{Recommandation :} réduire légèrement le graphique est la méthode la plus rapide et la plus propre.
\end{enumerate}
\textbf{Conclusion :} un rapport imprimé professionnel tient sur une page, porte l'identité de l'organisation, est paginé, et conserve ses repères (ligne d'en-têtes répétée) s'il s'étend sur plusieurs pages.
\end{exoresolu}

\begin{attention}[Les 5 erreurs qui coûtent des points au BEPC]
\begin{enumerate}
    \item \textbf{Oublier le \texttt{=} en début de formule.} \og{}C2*D2\fg{} s'affiche comme du texte et ne calcule rien. \textbf{Réflexe :} toute formule commence par le signe \texttt{=}.
    \item \textbf{Oublier les \texttt{\$} sur une valeur fixe.} Écrire \texttt{=E2*H1} puis recopier vers le bas fait glisser la référence vers \texttt{H2}, \texttt{H3}... (cellules vides) $\rightarrow$ résultats nuls ou erronés. \textbf{Règle :} tout paramètre constant (taux de TVA, seuil, coefficient) se référence en \textbf{absolu} : \texttt{\$H\$1}.
    \item \textbf{Inverser l'ordre des tests dans un \texttt{SI} imbriqué.} Tester \texttt{>=3000} avant \texttt{>=8000} classe la valeur \texttt{10000} en \og{}Moyenne\fg{}. \textbf{Règle :} toujours tester du seuil le \textbf{plus élevé} vers le seuil le \textbf{plus bas}.
    \item \textbf{Utiliser la virgule (,) au lieu du point-virgule (;) comme séparateur d'arguments.} En version française, on écrit \texttt{=SI(A1>10;"Ok";"Non")}. La virgule provoque une erreur de formule.
    \item \textbf{Imprimer sans vérifier l'aperçu (\texttt{Ctrl+P}).} Tableau coupé, graphique sur deux pages, en-têtes absents en page 2. \textbf{Réflexe obligatoire :} \texttt{Ctrl+P} $\rightarrow$ vérifier l'ajustement sur une page, la répétition des titres et la lisibilité $\rightarrow$ seulement ensuite Imprimer.
\end{enumerate}
\end{attention}

\newpage

\coursec{Exercices}

\courssub{Je vérifie mes connaissances}

\begin{exercice}[QCM : Formules et références]
Dans un tableur, la cellule \texttt{D4} contient la formule \texttt{=B2*\$C\$3}. On recopie cette formule vers la cellule \texttt{F6} (deux colonnes à droite, deux lignes en bas). Quelle sera la formule affichée dans \texttt{F6} ?
\begin{enumerate}
    \item \texttt{=D4*\$E\$5}
    \item \texttt{=D4*\$C\$3}
    \item \texttt{=D4*E5}
    \item \texttt{=B2*\$C\$3}
\end{enumerate}
\end{exercice}

\begin{exercice}[Vrai ou Faux : Fonctions statistiques]
Pour chaque affirmation, répondre \textbf{Vrai} ou \textbf{Faux} et justifier brièvement.
\begin{enumerate}
    \item La fonction \texttt{MOYENNE(A1:A10)} ignore les cellules contenant du texte dans la plage.
    \item La fonction \texttt{NBVAL(A1:A10)} compte uniquement les cellules contenant des nombres.
    \item La fonction \texttt{NB.SI(A1:A10;\textgreater=10)} compte le nombre de cellules dont la valeur est supérieure ou égale à 10.
    \item \texttt{MAX(A1:A10)} renvoie 0 si la plage ne contient que des nombres négatifs.
\end{enumerate}
\end{exercice}

\begin{exercice}[Définitions et syntaxe]
Compléter les définitions suivantes :
\begin{enumerate}
    \item La référence de cellule qui ne change pas lors d'une recopie (ex: \texttt{\$B\$5}) s'appelle une référence \underline{\hspace{4cm}}.
    \item La fonction \texttt{SI(condition; valeur\_si\_vrai; valeur\_si\_faux)} permet de réaliser un \underline{\hspace{4cm}}.
    \item Pour additionner les valeurs de la plage \texttt{B2:B20} \textbf{uniquement si} les cellules correspondantes de la plage \texttt{A2:A20} contiennent le mot "Ventes", on utilise la fonction \underline{\hspace{4cm}} avec la syntaxe \underline{\hspace{4cm}}.
    \item L'outil qui permet de colorier automatiquement les cellules dont la valeur est inférieure à la moyenne s'appelle \underline{\hspace{4cm}}.
\end{enumerate}
\end{exercice}

\begin{exercice}[Calcul direct et analyse d'erreur]
On considère les cellules suivantes :
\texttt{A1}=10, \texttt{A2}=20, \texttt{A3}=30, \texttt{A4}=Texte, \texttt{A5}=40.
Calculer le résultat affiché par chacune des formules ci-dessous (ou indiquer l'erreur) :
\begin{enumerate}
    \item \texttt{=SOMME(A1:A5)}
    \item \texttt{=MOYENNE(A1:A5)}
    \item \texttt{=NB(A1:A5)}
    \item \texttt{=NBVAL(A1:A5)}
\end{enumerate}
\end{exercice}

\courssub{Je m'entraîne}

\begin{exercice}[Gestion des notes d'une classe]
Un professeur saisit les notes sur 20 de 5 élèves en \texttt{B2:B6} : 12, 8, 15, 9, 14.
\begin{enumerate}
    \item Écrire la formule en \texttt{B7} pour calculer la moyenne de la classe.
    \item Écrire la formule en \texttt{B8} pour afficher "Admis" si la moyenne (cellule \texttt{B7}) est $\ge$ 10, sinon "Recalé".
    \item Écrire la formule en \texttt{B9} pour compter le nombre d'élèves ayant une note $\ge$ 10.
    \item Écrire la formule en \texttt{B10} pour afficher la meilleure note.
\end{enumerate}
\end{exercice}

\begin{exercice}[Facture avec références mixtes et absolues]
On construit une table de multiplication en \texttt{B2:K11}.
\begin{itemize}
    \item Ligne 1 (en-têtes) : \texttt{B1}=1, \texttt{C1}=2, ..., \texttt{K1}=10.
    \item Colonne A (en-têtes) : \texttt{A2}=1, \texttt{A3}=2, ..., \texttt{A11}=10.
\end{itemize}
On veut saisir une formule unique en \texttt{B2} et la recopier sur tout le bloc \texttt{B2:K11} pour obtenir le produit de l'en-tête de ligne par l'en-tête de colonne.
\begin{enumerate}
    \item Écrire la formule exacte à saisir en \texttt{B2} (utiliser les signes \$ nécessaires).
    \item Expliquer le rôle du \$ devant la lettre et devant le chiffre dans cette formule.
\end{enumerate}
\end{exercice}

\begin{exercice}[SI imbriquée et mention au BEPC]
La note d'un élève est en cellule \texttt{C3}. On veut afficher la mention correspondante dans \texttt{D3} selon le barème officiel :
\begin{center}
\begin{tabular}{|c|c|}
\hline
Note & Mention \\ \hline
$[0, 10[$ & "Échec" \\ \hline
$[10, 12[$ & "Passable" \\ \hline
$[12, 14[$ & "Assez Bien" \\ \hline
$[14, 16[$ & "Bien" \\ \hline
$[16, 20]$ & "Très Bien" \\ \hline
\end{tabular}
\end{center}
\begin{enumerate}
    \item Écrire la formule \texttt{SI} imbriquée complète à saisir en \texttt{D3}.
    \item Si la version du tableur le permet, écrire la formule équivalente avec la fonction \texttt{SI.MULTIPLE} (ou \texttt{IFS}).
\end{enumerate}
\end{exercice}

\begin{exercice}[Graphique et mise en forme conditionnelle]
On dispose d'un tableau des ventes mensuelles (Janvier à Décembre) en \texttt{A1:B13}.
\begin{enumerate}
    \item Quel type de graphique est le plus adapté pour visualiser l'évolution des ventes dans le temps ? Justifier.
    \item Décrire la procédure pour créer ce graphique (sélection, onglet, type).
    \item On veut que les mois où les ventes sont inférieures à 5000 apparaissent en police rouge et fond jaune clair. Décrire les étapes exactes de la mise en forme conditionnelle (règle "Formule" ou "Valeur de la cellule").
\end{enumerate}
\end{exercice}

\courssub{Je me prépare au Probatoire}

\begin{exercice}[Sujet type Probatoire : Bulletin de notes automatisé]
\textbf{Contexte :} Vous êtes délégué de classe au Lycée Bilingue de Nkolbisson. Vous devez automatiser le bulletin de notes du 3ème trimestre pour 30 élèves. Le tableau brut (feuille "Données") contient :
\begin{itemize}
    \item Colonne A : Nom (A2:A31)
    \item Colonne B : Prénom (B2:B31)
    \item Colonnes C à F : Notes sur 20 en Maths, PC, SVT, Anglais (C2:F31)
    \item Colonne G : Coefficient Maths (fixe 4), PC (fixe 3), SVT (fixe 2), Anglais (fixe 2) — \textbf{les coefficients sont en G1:J1}.
\end{itemize}

\textbf{Travail demandé :}
\begin{enumerate}
    \item \textbf{Moyenne pondérée :} Écrire la formule exacte à saisir en \texttt{K2} (à recopier vers \texttt{K31}) pour calculer la moyenne pondérée de l'élève. Utiliser les références absolues pour les coefficients.
    \item \textbf{Rang :} Écrire la formule en \texttt{L2} pour déterminer le rang de l'élève (1 pour le meilleur). Gérer les ex-aequo (même rang).
    \item \textbf{Mention :} En \texttt{M2}, écrire la formule \texttt{SI} imbriquée (ou \texttt{SI.MULTIPLE}) pour afficher la mention selon : $<10$ : "Échec" ; $[10,12[$ : "Passable" ; $[12,14[$ : "AB" ; $[14,16[$ : "B" ; $\ge 16$ : "TB".
    \item \textbf{Statistiques de la classe (cellules fixes) :}
    \begin{itemize}
        \item \texttt{P2} : Moyenne générale de la classe (moyenne des moyennes pondérées).
        \item \texttt{P3} : Nombre d'admis (moyenne $\ge 10$).
        \item \texttt{P4} : Pourcentage de réussite.
        \item \texttt{P5} : Meilleure moyenne (MAX).
    \end{itemize}
    Écrire les formules correspondantes.
    \item \textbf{Mise en forme conditionnelle :} Sur la plage \texttt{K2:K31}, décrire la règle pour colorier en vert les moyennes $\ge 14$, en orange celles $\ge 10$, en rouge les autres.
    \item \textbf{Graphique :} On veut un graphique en secteurs (camembert) montrant la répartition des mentions (Échec, Passable, AB, B, TB). Expliquer pourquoi il faut d'abord créer un tableau de synthèse (comptage par mention) avant de tracer le graphique, et indiquer la fonction à utiliser pour ce comptage.
\end{enumerate}
\end{exercice}

\begin{exercice}[Sujet type Probatoire : Gestion de stock d'une boutique]
\textbf{Contexte :} La boutique "Au Bon Prix" à Mfoundi gère son stock. Feuille "Stock" :
\begin{itemize}
    \item \texttt{A2:A51} : Code article (ex: "ART001")
    \item \texttt{B2:B51} : Désignation
    \item \texttt{C2:C51} : Prix unitaire (FCFA)
    \item \texttt{D2:D51} : Quantité en stock
    \item \texttt{E2:E51} : Seuil d'alerte (quantité min)
\end{itemize}

\textbf{Travail demandé :}
\begin{enumerate}
    \item \textbf{Valeur du stock :} En \texttt{F2}, formule pour la valeur du stock de l'article (Prix $\times$ Qté). Recopier.
    \item \textbf{Alerte stock :} En \texttt{G2}, formule affichant "COMMANDER" si Qté $\le$ Seuil, sinon "OK". Recopier.
    \item \textbf{Valorisation totale :} En \texttt{H1} (hors tableau), formule pour la valeur totale du stock (somme des valeurs).
    \item \textbf{Recherche article :} En \texttt{J2}, on saisit un Code article. En \texttt{K2}, écrire la formule \texttt{RECHERCHEV} (ou \texttt{RECHERCHE.X}) pour afficher la Désignation correspondante. En \texttt{L2}, afficher le Prix unitaire. Préciser le 4ème argument (VRAI/FAUX) et justifier.
    \item \textbf{Analyse par catégorie :} On ajoute une colonne \texttt{H} "Catégorie" (Boisson, Alimentaire, Hygiène...).
    \begin{itemize}
        \item En \texttt{P2:P4}, lister les 3 catégories.
        \item En \texttt{Q2}, formule \texttt{SOMME.SI} pour calculer la valeur totale du stock (colonne F) pour la catégorie en \texttt{P2}. Recopier.
        \item En \texttt{R2}, formule \texttt{NB.SI} pour compter le nombre d'articles différents par catégorie.
    \end{itemize}
    \item \textbf{Mise en page :} Le tableau fait 50 lignes $\times$ 8 colonnes. Décrire la procédure complète (4 étapes) pour l'imprimer sur une seule page A4 en mode Paysage, avec les en-têtes de colonnes (ligne 1) répétés sur chaque page si débordement, et des marges étroites.
\end{enumerate}
\end{exercice}

\begin{exercice}[Sujet type Probatoire : Analyse d'erreurs et correction]
\textbf{Contexte :} Un élève de 3ème a réalisé le tableau ci-dessous pour calculer la moyenne de 3 notes (B2, C2, D2) et décider de l'admission. Il a recopié les formules vers le bas. Il obtient des résultats aberrants.

\textbf{Extrait du tableur (Mode Affichage Formules \texttt{Ctrl+`} ) :}
\begin{center}
\begin{tabular}{|c|c|c|c|c|c|}
\hline
\textbf{A} & \textbf{B} & \textbf{C} & \textbf{D} & \textbf{E} & \textbf{F} \\ \hline
1 & Note1 & Note2 & Note3 & Moyenne & Décision \\ \hline
2 & 12 & 14 & 10 & \texttt{=(B2+C2+D2)/3} & \texttt{=SI(E2>10;"Admis";"Recalé")} \\ \hline
3 & 8 & 9 & 11 & \texttt{=(B3+C3+D3)/3} & \texttt{=SI(E3>10;"Admis";"Recalé")} \\ \hline
4 & 15 & 16 & 14 & \texttt{=(B4+C4+D4)/3} & \texttt{=SI(E4>10;"Admis";"Recalé")} \\ \hline
\end{tabular}
\end{center}

L'élève signale deux problèmes :
\begin{enumerate}
    \item Pour la ligne 3 (moyenne = 9,33), la décision affiche "Recalé". L'élève trouve cela normal. Le professeur dit que le seuil d'admission est 10 \textbf{sur 20}, mais que la décision devrait être "Admis" si moyenne $\ge$ 10. Corriger la formule en \texttt{F2}.
    \item L'élève a saisi une 4ème note en colonne \texttt{E} (déplaçant l'ancienne colonne E vers F). Il a modifié la formule de moyenne en \texttt{F2} en \texttt{=MOYENNE(B2:E2)}. Mais pour l'élève 2, la cellule \texttt{E2} contient le texte "Absent". Quel est le résultat affiché par \texttt{MOYENNE(B2:E2)} ? Est-ce correct ? Proposer une formule robuste qui ignore le texte et calcule la moyenne sur les 3 notes numériques uniquement (supposons que "Absent" ne doit pas compter comme 0).
    \item L'élève veut ajouter un graphique en barres groupées comparant les 3 notes de chaque élève (lignes 2 à 4). Il sélectionne \texttt{A1:D4} et insère un histogramme. Le graphique affiche une 4ème série "Moyenne" qu'il ne veut pas. Expliquer comment modifier la sélection source du graphique pour exclure la colonne Moyenne.
    \item L'élève veut imprimer ce petit tableau. Il clique sur Imprimer et découvre que le tableau est tout petit en haut à gauche de la page A4. Donner les 3 réglages principaux à modifier dans "Mise en page" / "Aperçu avant impression" pour centrer le tableau (horizontalement et verticalement) et l'agrandir pour qu'il remplisse la page.
\end{enumerate}
\end{exercice}

\newpage

\coursec{Corrigés des exercices}

\courssub{Je vérifie mes connaissances}

\begin{corrige}[Exercice 1 — QCM : Formules et références]
\textbf{Bonne réponse : b) \texttt{=D4*\$C\$3}.}

\textbf{Justification :} La formule \texttt{=B2*\$C\$3} contient deux références de nature différente :
\begin{itemize}
    \item \texttt{B2} est une référence \cle{relative} : lors d'une recopie, elle se décale du même nombre de colonnes et de lignes que le déplacement de la cellule. De \texttt{D4} vers \texttt{F6}, on se déplace de \textbf{2 colonnes vers la droite} et de \textbf{2 lignes vers le bas}. Donc \texttt{B2} devient \texttt{D4} (B $\rightarrow$ D pour les colonnes, 2 $\rightarrow$ 4 pour les lignes).
    \item \texttt{\$C\$3} est une référence \cle{absolue} : le signe \texttt{\$} « fige » la colonne C et la ligne 3. Elle ne change \textbf{jamais} lors d'une recopie, quelle que soit la destination.
\end{itemize}
La formule affichée dans \texttt{F6} est donc \texttt{=D4*\$C\$3}.

\textbf{Point de vigilance :} les réponses a) et c) sont des pièges classiques : on ne doit jamais faire varier une référence absolue (a), et on ne doit jamais « oublier » les \texttt{\$} (c).
\end{corrige}

\begin{corrige}[Exercice 2 — Vrai ou Faux : Fonctions statistiques]
\begin{enumerate}
    \item \textbf{Vrai.} La fonction \texttt{MOYENNE} ne prend en compte que les cellules contenant des \textbf{nombres} ; les cellules vides et les cellules contenant du texte sont ignorées (elles ne comptent ni au numérateur ni au dénominateur).
    \item \textbf{Faux.} \texttt{NBVAL} compte toutes les cellules \textbf{non vides} : nombres, textes, dates, valeurs logiques. C'est la fonction \texttt{NB} qui compte uniquement les nombres.
    \item \textbf{Vrai.} La syntaxe est \texttt{NB.SI(plage ; critère)}. Le critère \texttt{">=10"} (à écrire entre guillemets : \texttt{NB.SI(A1:A10;">=10")}) compte les cellules dont la valeur est supérieure ou égale à 10.
    \item \textbf{Faux.} \texttt{MAX} renvoie la \textbf{plus grande valeur} de la plage, même si toutes les valeurs sont négatives. Par exemple, \texttt{MAX(-5;-2;-8)} renvoie $-2$, et non 0.
\end{enumerate}
\textbf{Point de vigilance :} ne pas confondre \texttt{NB} (nombres seulement) et \texttt{NBVAL} (tout contenu non vide) : c'est une question fréquente au BEPC.
\end{corrige}

\begin{corrige}[Exercice 3 — Définitions et syntaxe]
\begin{enumerate}
    \item La référence de cellule qui ne change pas lors d'une recopie (ex : \texttt{\$B\$5}) s'appelle une référence \textbf{absolue}.
    \item La fonction \texttt{SI(condition; valeur\_si\_vrai; valeur\_si\_faux)} permet de réaliser un \textbf{test logique (ou choix conditionnel)}.
    \item On utilise la fonction \textbf{\texttt{SOMME.SI}} avec la syntaxe \texttt{=SOMME.SI(A2:A20;"Ventes";B2:B20)} (plage de critère ; critère ; plage à additionner).
    \item L'outil qui permet de colorier automatiquement les cellules selon leur valeur s'appelle la \textbf{mise en forme conditionnelle}.
\end{enumerate}
\end{corrige}

\begin{corrige}[Exercice 4 — Calcul direct et analyse d'erreur]
Données : \texttt{A1}=10, \texttt{A2}=20, \texttt{A3}=30, \texttt{A4}="Texte", \texttt{A5}=40.
\begin{enumerate}
    \item \texttt{=SOMME(A1:A5)} : la fonction \texttt{SOMME} ignore le texte. Résultat : $10+20+30+40 = \mathbf{100}$.
    \item \texttt{=MOYENNE(A1:A5)} : la moyenne est calculée sur les 4 valeurs \textbf{numériques} uniquement (le texte est ignoré, il ne compte pas comme 0) : $\dfrac{100}{4} = \mathbf{25}$.
    \item \texttt{=NB(A1:A5)} : \texttt{NB} compte uniquement les nombres : il y en a 4. Résultat : \textbf{4}.
    \item \texttt{=NBVAL(A1:A5)} : \texttt{NBVAL} compte toutes les cellules non vides (4 nombres + 1 texte). Résultat : \textbf{5}.
\end{enumerate}
\textbf{Point de vigilance :} l'erreur fréquente est de croire que \texttt{MOYENNE} divise par 5 ici : non, le texte est purement ignoré. En revanche, une cellule contenant 0 compterait bien dans la moyenne.
\end{corrige}

\courssub{Je m'entraîne}

\begin{corrige}[Exercice 5 — Gestion des notes d'une classe]
Notes en \texttt{B2:B6} : 12, 8, 15, 9, 14.
\begin{enumerate}
    \item Moyenne de la classe en \texttt{B7} : \texttt{=MOYENNE(B2:B6)}. Vérification : $\dfrac{12+8+15+9+14}{5} = \dfrac{58}{5} = 11,6$.
    \item Décision en \texttt{B8} : \texttt{=SI(B7>=10;"Admis";"Recalé")}. Ici $11,6 \ge 10$, donc la cellule affiche « Admis ».
    \item Nombre d'élèves ayant une note $\ge 10$ en \texttt{B9} : \texttt{=NB.SI(B2:B6;">=10")}. Vérification : 12, 15 et 14 sont $\ge 10$, donc le résultat est \textbf{3}.
    \item Meilleure note en \texttt{B10} : \texttt{=MAX(B2:B6)}, qui affiche \textbf{15}.
\end{enumerate}
\textbf{Point de vigilance :} le critère de \texttt{NB.SI} doit être écrit \textbf{entre guillemets} : \texttt{">=10"}. Écrire \texttt{NB.SI(B2:B6;>=10)} provoque une erreur de formule.
\end{corrige}

\begin{corrige}[Exercice 6 — Facture avec références mixtes et absolues]
\begin{enumerate}
    \item Formule exacte à saisir en \texttt{B2} :
    \begin{center}\texttt{=\$A2*B\$1}\end{center}
    \item \textbf{Explication du rôle des \texttt{\$} :}
    \begin{itemize}
        \item \texttt{\$A2} : le \texttt{\$} devant la lettre \textbf{fige la colonne A}. Quand on recopie la formule vers la droite (de \texttt{B2} vers \texttt{K2}), la référence doit toujours pointer sur l'en-tête de ligne situé en colonne A. En revanche, le numéro de ligne reste libre : recopiée vers le bas, elle devient \texttt{\$A3}, \texttt{\$A4}, etc., ce qui permet de changer de ligne d'en-tête.
        \item \texttt{B\$1} : le \texttt{\$} devant le chiffre \textbf{fige la ligne 1}. Quand on recopie vers le bas (de \texttt{B2} vers \texttt{B11}), la référence doit toujours pointer sur l'en-tête de colonne situé en ligne 1. La lettre de colonne reste libre : recopiée vers la droite, elle devient \texttt{C\$1}, \texttt{D\$1}, etc.
    \end{itemize}
    Ainsi, chaque cellule du bloc \texttt{B2:K11} calcule bien le produit (en-tête de sa ligne) $\times$ (en-tête de sa colonne). Par exemple, \texttt{F5} contiendra \texttt{=\$A5*F\$1}, soit $4 \times 5 = 20$.
\end{enumerate}
\textbf{Point de vigilance :} c'est le cas d'école des \cle{références mixtes} : on ne fige que ce qui doit rester fixe. Écrire \texttt{=\$A\$2*B\$1} bloquerait la ligne 2 et fausserait toute la table.
\end{corrige}

\begin{corrige}[Exercice 7 — SI imbriquée et mention au BEPC]
\begin{enumerate}
    \item Formule \texttt{SI} imbriquée complète en \texttt{D3} :
\begin{lstlisting}
=SI(C3<10;"Échec";SI(C3<12;"Passable";SI(C3<14;"Assez Bien";SI(C3<16;"Bien";"Très Bien"))))
\end{lstlisting}
    \textbf{Vérification du raisonnement :} les tests sont enchaînés du plus petit seuil au plus grand. Si \texttt{C3} $< 10$, on affiche « Échec » ; sinon, on sait déjà que \texttt{C3} $\ge 10$, donc le test \texttt{C3<12} correspond exactement à l'intervalle $[10;12[$, et ainsi de suite. Le dernier « sinon » correspond à \texttt{C3} $\ge 16$, soit $[16;20]$ : « Très Bien ».
    \item Avec \texttt{SI.MULTIPLE} (versions récentes d'Excel) :
\begin{lstlisting}
=SI.MULTIPLE(C3<10;"Échec";C3<12;"Passable";C3<14;"Assez Bien";C3<16;"Bien";C3>=16;"Très Bien")
\end{lstlisting}
\end{enumerate}
\textbf{Point de vigilance :} il faut autant de parenthèses fermantes que de fonctions \texttt{SI} ouvertes (ici 4). L'ordre des tests est essentiel : commencer par \texttt{C3<16} rendrait les mentions « Échec », « Passable » et « Assez Bien » inaccessibles.
\end{corrige}

\begin{corrige}[Exercice 8 — Graphique et mise en forme conditionnelle]
\begin{enumerate}
    \item \textbf{Type de graphique :} la \textbf{courbe} (graphique en lignes). Justification : on veut visualiser l'\emph{évolution} d'une grandeur (les ventes) \emph{en fonction du temps} (les 12 mois) ; la courbe est le type de graphique adapté à la représentation d'une évolution chronologique. (Un histogramme serait acceptable pour comparer les mois, mais il met moins bien en évidence la tendance.)
    \item \textbf{Procédure de création :}
    \begin{enumerate}
        \item Sélectionner la plage \texttt{A1:B13} (en-têtes incluses pour les légendes).
        \item Ouvrir l'onglet \textbf{Insertion}.
        \item Dans le groupe Graphiques, choisir \textbf{Insérer un graphique en courbes} (ou « Graphique recommandé » $\rightarrow$ Courbes).
        \item Le graphique s'insère dans la feuille ; on peut ensuite lui ajouter un titre (« Évolution des ventes mensuelles ») et déplacer la légende.
    \end{enumerate}
    \item \textbf{Mise en forme conditionnelle (ventes $< 5000$) :}
    \begin{enumerate}
        \item Sélectionner la plage des ventes \texttt{B2:B13}.
        \item Onglet \textbf{Accueil} $\rightarrow$ \textbf{Mise en forme conditionnelle} $\rightarrow$ \textbf{Règles de mise en surbrillance des cellules} $\rightarrow$ \textbf{Inférieur à...}
        \item Saisir la valeur \textbf{5000}.
        \item Choisir un format personnalisé : \textbf{police rouge} et \textbf{remplissage jaune clair}, puis valider par OK.
    \end{enumerate}
    Toute cellule de la plage dont la valeur devient inférieure à 5000 sera désormais automatiquement affichée en rouge sur fond jaune.
\end{enumerate}
\end{corrige}

\courssub{Je me prépare au Probatoire}

\begin{corrige}[Exercice 9 — Sujet type Probatoire : Bulletin de notes automatisé]
\textbf{Barème indicatif (sur 20) :} Q1 : 4 pts ; Q2 : 3 pts ; Q3 : 4 pts ; Q4 : 4 pts ; Q5 : 2 pts ; Q6 : 3 pts.

\begin{enumerate}
    \item \textbf{Moyenne pondérée en \texttt{K2} :} les coefficients sont en ligne 1 (\texttt{G1:J1}) et doivent rester fixes lors de la recopie vers le bas : on utilise des références absolues.
\begin{lstlisting}
=(C2*$G$1+D2*$H$1+E2*$I$1+F2*$J$1)/($G$1+$H$1+$I$1+$J$1)
\end{lstlisting}
    Vérification : la somme des coefficients vaut $4+3+2+2 = 11$ ; pour un élève ayant 12 partout : $\dfrac{12\times 11}{11} = 12$. La formule se recopie de \texttt{K2} vers \texttt{K31} sans modification des coefficients grâce aux \texttt{\$}.
    \item \textbf{Rang en \texttt{L2} :}
\begin{lstlisting}
=RANG(K2;$K$2:$K$31;0)
\end{lstlisting}
    Le 3\textsuperscript{e} argument \texttt{0} (ordre décroissant) donne le rang 1 à la meilleure moyenne. La plage \texttt{\$K\$2:\$K\$31} est absolue pour rester fixe à la recopie. \textbf{Gestion des ex-aequo :} la fonction \texttt{RANG} attribue automatiquement le \emph{même rang} aux élèves ayant la même moyenne (puis saute le rang suivant), ce qui répond à la consigne.
    \item \textbf{Mention en \texttt{M2} :}
\begin{lstlisting}
=SI(K2<10;"Échec";SI(K2<12;"Passable";SI(K2<14;"AB";SI(K2<16;"B";"TB"))))
\end{lstlisting}
    \item \textbf{Statistiques de la classe :}
    \begin{itemize}
        \item \texttt{P2} (moyenne générale) : \texttt{=MOYENNE(K2:K31)}
        \item \texttt{P3} (nombre d'admis) : \texttt{=NB.SI(K2:K31;">=10")}
        \item \texttt{P4} (pourcentage de réussite) : \texttt{=P3/30} au format pourcentage, ou directement \texttt{=NB.SI(K2:K31;">=10")/NB(K2:K31)}. Exemple : si 24 admis, $\dfrac{24}{30} = 0,8 = 80\,\%$.
        \item \texttt{P5} (meilleure moyenne) : \texttt{=MAX(K2:K31)}
    \end{itemize}
    \item \textbf{Mise en forme conditionnelle sur \texttt{K2:K31} :} sélectionner \texttt{K2:K31}, puis Accueil $\rightarrow$ Mise en forme conditionnelle $\rightarrow$ Nouvelle règle, et créer \textbf{dans cet ordre} :
    \begin{enumerate}
        \item « Valeur de la cellule supérieure ou égale à 14 » $\rightarrow$ remplissage vert ;
        \item « Valeur de la cellule supérieure ou égale à 10 » $\rightarrow$ remplissage orange ;
        \item « Valeur de la cellule inférieure à 10 » $\rightarrow$ remplissage rouge.
    \end{enumerate}
    \textbf{Point de vigilance :} l'ordre des règles est crucial : la règle « $\ge 14$ » doit être testée \emph{avant} « $\ge 10$ », sinon toutes les moyennes $\ge 10$ (y compris celles $\ge 14$) seraient coloriées en orange.
    \item \textbf{Graphique en secteurs :} un camembert représente des \emph{parts d'un total} ; il lui faut donc une série de \textbf{valeurs numériques} (le nombre d'élèves par mention), et non la colonne de textes \texttt{M2:M31} que le tableur ne sait pas exploiter directement. On crée d'abord un \textbf{tableau de synthèse} : en colonne, les 5 mentions (Échec, Passable, AB, B, TB) ; en regard, l'effectif de chaque mention calculé par :
\begin{lstlisting}
=NB.SI($M$2:$M$31;"Échec")
\end{lstlisting}
    (puis « Passable », « AB », « B », « TB »). On sélectionne ensuite ce petit tableau de synthèse et on insère le graphique en secteurs.
\end{enumerate}
\textbf{Points de vigilance (rédaction attendue) :} justifier les références absolues pour les coefficients et pour la plage du rang ; écrire les critères entre guillemets ; préciser le rôle du 3\textsuperscript{e} argument de \texttt{RANG}.
\end{corrige}

\begin{corrige}[Exercice 10 — Sujet type Probatoire : Gestion de stock d'une boutique]
\textbf{Barème indicatif (sur 20) :} Q1 : 2 pts ; Q2 : 3 pts ; Q3 : 2 pts ; Q4 : 5 pts ; Q5 : 4 pts ; Q6 : 4 pts.

\begin{enumerate}
    \item \textbf{Valeur du stock en \texttt{F2} :} \texttt{=C2*D2}, recopiée vers \texttt{F51} (références relatives : elles s'adaptent ligne par ligne).
    \item \textbf{Alerte en \texttt{G2} :} \texttt{=SI(D2<=E2;"COMMANDER";"OK")}, recopiée vers \texttt{G51}.
    \item \textbf{Valorisation totale en \texttt{H1} :} \texttt{=SOMME(F2:F51)}.
    \item \textbf{Recherche article :}
    \begin{itemize}
        \item Désignation en \texttt{K2} : \texttt{=RECHERCHEV(J2;A2:F51;2;FAUX)}
        \item Prix unitaire en \texttt{L2} : \texttt{=RECHERCHEV(J2;A2:F51;3;FAUX)}
    \end{itemize}
    \textbf{4\textsuperscript{e} argument : FAUX.} Justification : on cherche une correspondance \textbf{exacte} du code article (un code « ART007 » ne doit pas renvoyer la désignation d'un autre article). L'argument \texttt{VRAI} (ou omis) exigerait une colonne triée et autoriserait des correspondances approximatives, inadaptées ici. Les colonnes sont numérotées depuis la première colonne de la plage : A = 1, B (désignation) = 2, C (prix) = 3.
    \item \textbf{Analyse par catégorie} (colonne H ajoutée) :
    \begin{itemize}
        \item \texttt{P2} = « Boisson », \texttt{P3} = « Alimentaire », \texttt{P4} = « Hygiène ».
        \item \texttt{Q2} : \texttt{=SOMME.SI(H2:H51;P2;F2:F51)} — on additionne les valeurs de \texttt{F} pour les lignes dont la catégorie en \texttt{H} égale le contenu de \texttt{P2}. Recopier vers \texttt{Q4}.
        \item \texttt{R2} : \texttt{=NB.SI(H2:H51;P2)} — compte le nombre d'articles de la catégorie. Recopier vers \texttt{R4}.
    \end{itemize}
    \textbf{Point de vigilance :} dans \texttt{SOMME.SI}, le critère est ici une \emph{référence de cellule} (\texttt{P2}), donc \textbf{sans guillemets} ; les guillemets ne servent que pour un texte ou un critère écrit directement (\texttt{">=10"}).
    \item \textbf{Mise en page pour l'impression (4 étapes) :}
    \begin{enumerate}
        \item Onglet \textbf{Mise en page} $\rightarrow$ \textbf{Orientation} : choisir \textbf{Paysage} ; puis \textbf{Marges} $\rightarrow$ \textbf{Marges étroites}.
        \item \textbf{Mise à l'échelle} : dans le groupe « Mise à l'échelle », régler « Ajuster à » sur \textbf{1 page en largeur} (et 1 en hauteur, ou laisser la hauteur automatique si le tableau tient en une page).
        \item \textbf{Imprimer les titres} : onglet Mise en page $\rightarrow$ \textbf{Imprimer les titres} $\rightarrow$ dans « Lignes à répéter en haut », saisir \texttt{\$1:\$1} (la ligne d'en-têtes sera répétée sur chaque page en cas de débordement).
        \item Vérifier avec \textbf{Fichier $\rightarrow$ Imprimer} (Aperçu avant impression) que le tableau tient sur une page A4 et que les en-têtes apparaissent bien, puis lancer l'impression.
    \end{enumerate}
\end{enumerate}
\end{corrige}

\begin{corrige}[Exercice 11 — Sujet type Probatoire : Analyse d'erreurs et correction]
\textbf{Barème indicatif (sur 20) :} Q1 : 4 pts ; Q2 : 6 pts ; Q3 : 5 pts ; Q4 : 5 pts.

\begin{enumerate}
    \item \textbf{Correction du seuil d'admission :} la formule \texttt{=SI(E2>10;"Admis";"Recalé")} utilise l'inégalité \textbf{stricte} : une moyenne exactement égale à 10 affiche « Recalé », ce qui contredit la règle « Admis si moyenne $\ge 10$ ». Formule corrigée en \texttt{F2} :
    \begin{center}\texttt{=SI(E2>=10;"Admis";"Recalé")}\end{center}
    (Pour la ligne 3, la moyenne 9,33 reste bien « Recalé » : le problème ne concernait que le cas limite 10/20.)
    \item \textbf{Effet du texte « Absent » :} la fonction \texttt{MOYENNE(B2:E2)} \textbf{ignore} la cellule contenant le texte « Absent » : elle calcule la moyenne des 3 notes numériques uniquement. Pour l'élève 2 : $\dfrac{8+9+11}{3} = \dfrac{28}{3} \approx 9,33$. \textbf{C'est correct} au sens où « Absent » ne compte pas comme 0 ; c'est même le comportement souhaité. La formule \texttt{=MOYENNE(B2:E2)} est donc déjà « robuste » face au texte. Si l'on voulait une formule explicitement sécurisée (par exemple pour compter les notes réellement saisies), on pourrait écrire :
\begin{lstlisting}
=SOMME(B2:E2)/NB(B2:E2)
\end{lstlisting}
    car \texttt{SOMME} ignore le texte et \texttt{NB} ne compte que les nombres : le quotient donne bien la moyenne sur les seules notes numériques. \textbf{Point de vigilance :} le piège serait d'utiliser \texttt{NBVAL} au dénominateur (il compterait « Absent » et fausserait la moyenne : $\frac{28}{4} = 7$).
    \item \textbf{Exclure la série « Moyenne » du graphique :} le graphique a été construit sur la plage \texttt{A1:D4}, qui inclut la colonne Moyenne ; le tableur l'a donc ajoutée comme 4\textsuperscript{e} série. Pour corriger : cliquer sur le graphique, puis clic droit $\rightarrow$ \textbf{Sélectionner des données} (ou onglet Création de graphique $\rightarrow$ Sélectionner des données). Dans la boîte de dialogue, \textbf{supprimer la série « Moyenne »} de la liste des séries, ou redéfinir la plage de données en \texttt{A1:C4} (colonnes Note1 à Note3 uniquement). Le graphique ne comparera alors que les 3 notes de chaque élève.
    \item \textbf{Mise en page pour centrer et agrandir le tableau :}
    \begin{enumerate}
        \item \textbf{Centrage :} onglet Mise en page $\rightarrow$ lanceur de boîte de dialogue (petite flèche en bas du groupe) $\rightarrow$ onglet \textbf{Marges} $\rightarrow$ cocher « \textbf{Centrer horizontalement} » et « \textbf{Centrer verticalement} ».
        \item \textbf{Agrandissement :} dans l'onglet \textbf{Page} de la même boîte (ou dans l'Aperçu avant impression), régler la \textbf{mise à l'échelle} : choisir « Ajuster à 1 page de large × 1 page de haut », ou augmenter manuellement le pourcentage d'échelle (par exemple 150 à 200\,\%) jusqu'à ce que le tableau remplisse la page.
        \item \textbf{Vérification :} passer par \textbf{Fichier $\rightarrow$ Imprimer} (Aperçu avant impression) pour contrôler le rendu (tableau centré, lisible, sur une seule page) avant de lancer l'impression.
    \end{enumerate}
\end{enumerate}
\textbf{Points de vigilance (rédaction attendue) :} distinguer clairement \texttt{>} et \texttt{>=} dans un test logique ; savoir que \texttt{MOYENNE}, \texttt{SOMME} et \texttt{NB} ignorent le texte alors que \texttt{NBVAL} le compte ; savoir modifier la source de données d'un graphique sans le recréer.
\end{corrige}

\newpage

\coursec{Fiche bilan}

\begin{popfigure}
\begin{tikzpicture}[mindmap, grow cyclic, every node/.style={concept, text=white, font=\small\bfseries},
root concept/.append style={concept color=popblue, font=\large\bfseries},
level 1/.append style={level distance=3.6cm, sibling angle=60},
level 2/.append style={level distance=2.2cm, sibling angle=45, font=\scriptsize}]
\node[root concept]{Utilisation d'un tableur}
  child[concept color=poporange]{ node {Environnement et mise en forme}
    child { node {Cellules, feuilles, classeur} }
    child { node {Formats, bordures} }
  }
  child[concept color=popgreen]{ node {Formules et références}
    child { node {Commence par =} }
    child { node {A1, \$A\$1} }
  }
  child[concept color=poppurple]{ node {Fonctions maths et stats}
    child { node {SOMME, MOYENNE} }
    child { node {MAX, MIN, NB} }
  }
  child[concept color=poppink]{ node {Fonctions logiques}
    child { node {SI, ET, OU, NON} }
    child { node {SI imbriqués} }
  }
  child[concept color=popgold]{ node {Mise en forme conditionnelle}
    child { node {Règles, couleurs} }
    child { node {Graphiques} }
  }
  child[concept color=popmuted]{ node {Mise en page et impression}
    child { node {Marges, orientation} }
    child { node {Aperçu avant impression} }
  };
\end{tikzpicture}
\legende{Carte mentale de la leçon : les six compétences clés de l'utilisation d'un tableur.}
\end{popfigure}

\begin{aretenir}[L'essentiel de la leçon -- Définitions clés]
\begin{itemize}
  \item \cle{Classeur} : fichier d'un tableur (Microsoft Excel, LibreOffice Calc) contenant une ou plusieurs \cle{feuilles de calcul}.
  \item \cle{Feuille de calcul} : grille formée de lignes (repérées par des numéros : 1, 2, 3\ldots) et de colonnes (repérées par des lettres : A, B, C\ldots).
  \item \cle{Cellule} : intersection d'une ligne et d'une colonne, repérée par sa référence, ex. \texttt{B3} (colonne B, ligne 3). La cellule en cours d'utilisation est appelée \cle{cellule active}.
  \item \cle{Plage de cellules} : ensemble rectangulaire de cellules adjacentes, noté avec le signe deux-points, ex. \texttt{A1:C10} (de la cellule A1 à la cellule C10).
  \item \cle{Formule} : expression saisie dans une cellule, qui commence toujours par le signe \texttt{=} et qui calcule un résultat à partir de valeurs ou d'autres cellules.
  \item \cle{Fonction} : formule prédéfinie du tableur, repérée par un nom suivi de parenthèses contenant des arguments, ex. \texttt{SOMME(B2:B10)}, \texttt{SI(B2>=10;"Admis";"Ajourné")}.
\end{itemize}
\end{aretenir}

\begin{aretenir}[Fonctions à connaître par c\oe ur]
\begin{center}
\begin{tabularx}{\linewidth}{|l|X|l|}
\hline
\rowcolor{popblueL} \textbf{Fonction} & \textbf{Rôle} & \textbf{Exemple} \\
\hline
\texttt{SOMME(plage)} & Additionne les valeurs & \texttt{=SOMME(B2:B10)} \\
\hline
\texttt{MOYENNE(plage)} & Calcule la moyenne & \texttt{=MOYENNE(B2:B10)} \\
\hline
\texttt{MAX(plage)} & Plus grande valeur & \texttt{=MAX(B2:B10)} \\
\hline
\texttt{MIN(plage)} & Plus petite valeur & \texttt{=MIN(B2:B10)} \\
\hline
\texttt{NB(plage)} & Compte les cellules contenant des nombres & \texttt{=NB(B2:B10)} \\
\hline
\texttt{SI(test;vrai;faux)} & Teste une condition et affiche un résultat selon qu'elle est vraie ou fausse & \texttt{=SI(B2>=10;"Admis";"Ajourné")} \\
\hline
\texttt{ET(cond1;cond2)} & VRAI si toutes les conditions sont vraies & \texttt{=ET(B2>=10;C2>=10)} \\
\hline
\texttt{OU(cond1;cond2)} & VRAI si au moins une condition est vraie & \texttt{=OU(B2>=10;C2>=10)} \\
\hline
\texttt{NON(cond)} & Inverse le résultat logique & \texttt{=NON(B2>=10)} \\
\hline
\end{tabularx}
\end{center}
\textbf{Attention :} dans les fonctions, les arguments sont séparés par le point-virgule \texttt{;} et un texte s'écrit toujours entre guillemets, ex. \texttt{"Admis"}.
\end{aretenir}

\begin{aretenir}[Références de cellules]
\begin{itemize}
  \item \cle{Référence relative} : \texttt{A1} ; elle change automatiquement lors de la recopie de la formule vers d'autres cellules.
  \item \cle{Référence absolue} : \texttt{\$A\$1} ; elle reste fixe lors de la recopie : le symbole \texttt{\$} bloque la colonne et/ou la ligne. On l'utilise pour désigner une valeur unique (taux, coefficient) commune à tous les calculs.
\end{itemize}
\textbf{Exemple :} si la formule \texttt{=B2*\$E\$1} est recopiée vers le bas, \texttt{B2} devient \texttt{B3}, \texttt{B4}\ldots{} mais \texttt{\$E\$1} ne change pas.
\end{aretenir}

\begin{methode}[Méthodes express à retenir]
\begin{enumerate}
  \item \textbf{Saisir une formule} : cliquer sur la cellule $\rightarrow$ taper le signe \texttt{=} $\rightarrow$ écrire la formule $\rightarrow$ valider par la touche Entrée.
  \item \textbf{Recopier une formule} : sélectionner la cellule puis tirer la \cle{poignée de recopie} (petit carré en bas à droite de la cellule active) vers le bas ou vers la droite.
  \item \textbf{Mettre en forme un tableau} : sélectionner les cellules $\rightarrow$ appliquer police, gras, couleurs, bordures et format d'affichage (nombre, monétaire en FCFA, pourcentage).
  \item \textbf{Appliquer une mise en forme conditionnelle} : sélectionner la plage $\rightarrow$ menu \emph{Mise en forme conditionnelle} $\rightarrow$ choisir la règle (ex. colorer en rouge les notes inférieures à 10).
  \item \textbf{Créer un graphique} : sélectionner les données $\rightarrow$ menu \emph{Insertion} $\rightarrow$ choisir le type adapté : histogramme (comparaison), courbes (évolution), secteurs (répartition).
  \item \textbf{Imprimer} : régler les marges et l'orientation (portrait ou paysage) $\rightarrow$ vérifier le résultat avec l'\cle{aperçu avant impression} $\rightarrow$ lancer l'impression.
\end{enumerate}
\end{methode}

\begin{attention}[Erreurs fréquentes à éviter]
\begin{itemize}
  \item Oublier le signe \texttt{=} au début d'une formule : le tableur affiche alors du texte au lieu de calculer.
  \item Écrire \texttt{=SOMME(B2+B10)} au lieu de \texttt{=SOMME(B2:B10)} : la première formule n'additionne que deux cellules.
  \item Oublier les guillemets autour d'un texte dans la fonction \texttt{SI}.
  \item Utiliser une référence relative là où une référence absolue (\texttt{\$}) est nécessaire : la recopie donne alors des résultats faux.
  \item L'affichage \texttt{\#\#\#\#\#} n'est pas une erreur de calcul : la colonne est simplement trop étroite, il faut l'élargir.
\end{itemize}
\end{attention}

\begin{aretenir}[Checklist avant le BEPC]
\begin{itemize}
  \item $\square$ Je sais repérer une cellule par sa référence (ex. \texttt{C5}) et une plage (ex. \texttt{A1:C10}).
  \item $\square$ Je sais que toute formule commence par le signe \texttt{=}.
  \item $\square$ Je sais écrire et recopier une formule simple avec la poignée de recopie.
  \item $\square$ Je connais les opérateurs de calcul : \texttt{+}, \texttt{-}, \texttt{*}, \texttt{/} et les opérateurs de comparaison : \texttt{=}, \texttt{<}, \texttt{>}, \texttt{>=}, \texttt{<=}, \texttt{<>} (différent de).
  \item $\square$ Je distingue référence relative (\texttt{A1}) et référence absolue (\texttt{\$A\$1}).
  \item $\square$ Je maîtrise \texttt{SOMME}, \texttt{MOYENNE}, \texttt{MAX}, \texttt{MIN} et \texttt{NB}.
  \item $\square$ Je sais construire une fonction \texttt{SI} et des \texttt{SI} imbriqués (ex. mention d'un bulletin).
  \item $\square$ Je sais utiliser \texttt{ET}, \texttt{OU} et \texttt{NON} dans un test logique.
  \item $\square$ Je sais appliquer une mise en forme (police, bordures, format nombre) et une mise en forme conditionnelle.
  \item $\square$ Je sais créer un graphique adapté aux données.
  \item $\square$ Je sais régler la mise en page et vérifier l'aperçu avant d'imprimer.
  \item $\square$ Je sais justifier ma démarche sur une situation concrète (bulletin de notes, facture d'une boutique, budget d'une classe).
\end{itemize}
\end{aretenir}

\findecours

\end{document}
